\documentclass[11pt,a4paper,oneside,openany]{book}
\usepackage{fontspec}
\setmainfont{Times New Roman}
\setsansfont{Arial}
\setmonofont[Scale=.86]{Consolas}
\usepackage{amsmath,amssymb,array,longtable,booktabs,graphicx,xcolor}
\usepackage[margin=20mm,headheight=15pt,headsep=7mm,footskip=11mm]{geometry}
\usepackage{fancyhdr}
\usepackage[unicode,colorlinks=true,linkcolor=teal!65!black,urlcolor=blue!60!black,
  bookmarksnumbered=true,bookmarksopen=false,pdfencoding=auto]{hyperref}
\hypersetup{pdftitle={Sổ tay thuật toán PPS},pdfauthor={Tài liệu đối chiếu mã nguồn PPS},
  pdfsubject={Đầu vào, công thức, phép tính, vòng lặp và giới hạn của công cụ}}
\renewcommand{\chaptername}{Chương}
\renewcommand{\contentsname}{Mục lục}
\renewcommand{\appendixname}{Phụ lục}
\renewcommand{\tablename}{Bảng}
\renewcommand{\figurename}{Hình}
\renewcommand{\bibname}{Nguồn trong repo}
\setcounter{tocdepth}{1}
\setcounter{secnumdepth}{2}
\setcounter{LTchunksize}{200}
\setlength{\parindent}{0pt}
\setlength{\parskip}{5pt plus 1pt}
\setlength{\tabcolsep}{4pt}
\renewcommand{\arraystretch}{1.18}
\setlength{\emergencystretch}{2em}
\allowdisplaybreaks[2]
\raggedbottom
% Keep a worked example's heading with its input and first table.
\newcommand{\examplebreak}{\par\ifdim\dimexpr\pagegoal-\pagetotal\relax<230pt\relax\newpage\fi}
\pagestyle{fancy}
\fancyhf{}
\fancyhead[L]{\small\sffamily SỔ TAY THUẬT TOÁN PPS}
\fancyhead[R]{\small\sffamily\nouppercase{\leftmark}}
\fancyfoot[L]{\footnotesize Đối chiếu mã nguồn trong repo}
\fancyfoot[R]{\thepage}
\renewcommand{\chaptermark}[1]{\markboth{\thechapter. #1}{}}
\fancypagestyle{plain}{\fancyhf{}\fancyfoot[L]{\footnotesize Sổ tay PPS}\fancyfoot[R]{\thepage}\renewcommand{\headrulewidth}{0pt}}
\begin{document}
\frontmatter
\begin{titlepage}
\color{teal!65!black}\sffamily
\vspace*{18mm}
{\large PHƯƠNG PHÁP SỐ · MI3042\par}
\vspace{12mm}
{\Huge\bfseries Sổ tay\par thuật toán PPS\par}
\vspace{8mm}
{\Large Input của tool $\longrightarrow$ quá trình tính $\longrightarrow$ output\par}
\vspace{12mm}
\rule{\textwidth}{1.2pt}
\vspace{4mm}
{\large Nội suy · Horner · Nội suy ngược · Spline\par
 Bình phương tối thiểu · Đạo hàm · Tích phân\par
 Phương trình vi phân · Bài toán biên · Trị riêng\par}
\vfill
\color{black}\rmfamily
Mỗi nhóm bắt đầu bằng \textbf{tên ô nhập trên GUI}, cách biến đổi dữ liệu, \textbf{khuôn bảng tổng quát}, thứ tự tính và \textbf{output cần đọc}. Các công thức tổng/tích có dạng khai triển để điền số; ví dụ cụ thể đặt sau khuôn làm chung.

Ví dụ nhỏ được tự dựng để thấy hết phép tính. Ví dụ đề 20241 và dữ liệu 20251 được ghi riêng nguồn, tập mốc và mục tiêu. Các con số được đối chiếu bằng công cụ trong repo; phần suy diễn không được nhận là nhật ký nội bộ nếu app không lưu nhật ký đó.
\vspace{8mm}
{\small Bản đối chiếu ngày 02/11/2026. Nguồn chỉnh sửa trực tiếp: \texttt{So\_tay\_PPS.tex} và thư mục \texttt{chapters/}.\par}
\end{titlepage}
\tableofcontents
\chapter*{Cách đọc và chuẩn bị một bài thi}
\addcontentsline{toc}{chapter}{Cách đọc và chuẩn bị một bài thi}
\markboth{Cách đọc}{}

\section*{Đọc theo luồng của tool, không bắt đầu từ ví dụ số}
Mỗi chương bắt đầu bằng một phần \textbf{Input -- quá trình -- output}. Tìm đúng tên nhóm và lựa chọn trên GUI, đọc tên các ô, rồi điền \emph{khuôn bảng tổng quát} theo dữ liệu đề. Hàng $i$ chỉ đại diện cho các dòng giữa; khi làm bài phải điền mọi dòng thật, không cộng thêm hàng ký hiệu này vào dữ liệu.

Khi bảng có hàng tổng, công thức ngay dưới cho biết mỗi tổng cộng những cột nào và được chuyển vào hệ/phép tính nào. Với tích, dấu ba chấm phải giữ đúng thứ tự và quy tắc bỏ mốc. Khi có lặp, đọc rõ đại lượng nhận từ vòng trước, đại lượng mới, tiêu chuẩn dừng và thứ gì được chuyển sang vòng sau.

Chương Bình phương tối thiểu tách riêng từng mô hình và cách khôi phục tham số. Gauss, QR và SVD được coi là những bộ giải đã biết; tài liệu tập trung bài toán mà tool đưa vào bộ giải và kết quả mà tool lấy ra.

\section*{Mỗi bài giải cần công khai sáu điều}
\begin{enumerate}
\item \textbf{Dữ liệu}: tên tệp, sheet, cột $x,y$, số dòng, đơn vị; dữ liệu đo hay giá trị hàm. Không lấy cột tham số làm quan sát, không bỏ dòng âm chỉ để lấy được log.
\item \textbf{Câu hỏi}: tìm $y$ từ $x$, tìm $x$ từ $y$, đạo hàm, tích phân hay khớp mô hình? Tìm trên nhánh nào, đoạn nào? Đề có quy định bậc, số mốc và điều kiện biên không?
\item \textbf{Mô hình được dùng}: đa thức qua những mốc nào, spline với biên nào, cơ sở hồi quy nào, hay hệ sai phân nào. Hai tập mốc khác nhau có thể cho hai đáp số khác nhau mà đều tính đúng mô hình đã khai báo.
\item \textbf{Thuật toán}: ghi công thức trước, định nghĩa biến rồi mới thế số. Phân biệt $x$ vật lý với $t=(x-x_0)/h$, hệ số theo bậc tăng với hệ số theo bậc giảm.
\item \textbf{Tính toán}: bảng mốc, sai phân, hệ số, khử/thế ngược, các tầng RK hoặc các vòng lặp. Hàng nào nhận đầu vào từ hàng nào phải chỉ rõ. Trong các ví dụ nhỏ của sổ tay, không lược vòng lặp ở giữa.
\item \textbf{Kết luận có điều kiện}: đáp số, đơn vị, mô hình, phần dư và lý do dừng. Nêu riêng sai số dữ liệu, sai số mô hình và sai số giải; một phần dư nhỏ không tự chứng minh cả ba đều nhỏ.
\end{enumerate}

\clearpage
\section*{Chọn chương theo dạng đề}
\begin{center}\small\begin{tabular}{@{}p{.30\textwidth}p{.43\textwidth}p{.19\textwidth}@{}}
\toprule
Đề hỏi & Quyết định đầu tiên & Đọc\\\midrule
Giá trị/đạo hàm từ bảng & Mốc đều hay bất kỳ? Bậc, tập mốc, gốc? & Chương \ref{ch:interp}\\
Giá trị đa thức đã biết & Nhập đủ hệ số, kể cả các hệ số 0 & Horner, Chương \ref{ch:interp}\\
Tìm $x$ sao cho $y=\bar y$ & Quét mọi khoảng đổi dấu; xác định nhánh; chọn hàm ngược hay giải đa thức thuận & Chương \ref{ch:inverse}\\
Hàm từng đoạn qua bảng & Bậc spline và đúng số điều kiện biên độc lập & Chương \ref{ch:spline}\\
Khớp dữ liệu nhiễu & Mô hình, trọng số, thang sai số gốc hay biến đổi & Chương \ref{ch:ls}\\
Đạo hàm trực tiếp từ bảng & Bậc $p$, cấp $s$, $p+1$ mốc được dùng & Chương \ref{ch:calculus}\\
Tích phân từ hàm/bảng & Quy tắc, lưới hợp lệ; cận đạo hàm hay Runge & Chương \ref{ch:calculus}\\
Phương trình có điều kiện đầu & Biến độc lập $t$, trạng thái $x,y,z$, bước $h$ & Chương \ref{ch:ode}\\
Điều kiện tại hai đầu / trị riêng & Đúng dấu phương trình, thứ tự ba số biên & Chương \ref{ch:bvp}\\
Tên hàm Python cũ & Phân biệt API bổ trợ với luồng GUI đang dùng & Phụ lục \ref{app:api}\\
Áp dụng vào đề thật & Kiểm tra sheet, biến thể đề và giả thiết dữ liệu & Phụ lục \ref{app:exam}\\
\bottomrule
\end{tabular}\end{center}

\section*{Ký hiệu và cách đọc kết quả}
$m$ là số quan sát được dùng, $n$ thường là bậc đa thức hoặc số khoảng chia (sẽ định nghĩa lại ngay tại từng thuật toán). $h$ là bước lưới; $z$ là điểm đánh giá; $\bar y$ là mức cần tìm ngược. Trong hồi quy, $z$ còn có thể là vector đáp ứng đã biến đổi, được định nghĩa trong chương đó. $c$ là vector hệ số, $\Phi$ là ma trận cơ sở. Không đồng nhất $c$ với vector hoành độ $\mathbf{x}$ của dữ liệu.

Đầu vào số dùng dấu chấm thập phân. Các hàm cơ sở ngăn bằng dấu chấm phẩy. Ví dụ \texttt{1; x; x\^{}2}; hệ ODE \texttt{y; -x}. Dấu trừ, log, căn phải nằm trong miền xác định tại mọi điểm mà thuật toán cần đánh giá, kể cả các điểm trung gian.

\textbf{Tổng hợp} phục vụ đọc nhanh hệ số và bảng vòng lặp; \textbf{Chi tiết} mở các phép thế số khi tính tay. Số chữ số hiển thị là cách viết kết quả, không phải yêu cầu làm tròn mọi phép tính trung gian. Các bảng trong sổ tay làm tròn để dễ đọc; bộ tính giữ độ chính xác nội bộ của từng mô-đun.

\section*{Khi nào được nói ``hội tụ''?}
Nội suy hữu hạn, Horner, khử Gauss, Thomas và dựng spline là thuật toán hữu hạn: hoàn tất khi tính hết hàng và kiểm tra hệ. Với nội suy ngược hoặc LM, số vòng phụ thuộc dữ liệu và tiêu chuẩn dừng. Với tích phân tự chia lưới, vòng lặp tăng số khoảng để ước lượng sai số. Với ODE, $h$ quyết định sai số rời rạc; dung sai bước ẩn chỉ kiểm soát việc giải phương trình trong một bước. Không tráo các ý nghĩa này.

Tài liệu mô tả cách tính và giới hạn đã kiểm tra, không khẳng định mọi dữ liệu tùy ý đều có nghiệm duy nhất hay mọi phép lặp đều hội tụ. Nếu thuật toán không phù hợp giả thiết đề, cần công khai lý do trước khi đổi phương pháp.

\mainmatter
\chapter{Nội suy đa thức, chọn mốc và Horner}
\label{ch:interp}
\section{Bắt đầu từ các ô nhập: ba tool liên quan đến nội suy}
\subsection*{Input: các ô đúng tên trên giao diện}
Bảng này được lấy từ biểu mẫu của tool. Giá trị mặc định là ví dụ nhập, cần thay bằng dữ liệu đề. Hai dãy $x,y$ là bảng dữ liệu riêng bên dưới nếu chức năng dùng bảng.
\subsubsection*{Ô chung của các lựa chọn trong nhóm}
\begin{center}\small\begin{tabular}{@{}p{.32\linewidth}p{.17\linewidth}p{.43\linewidth}@{}}\toprule
Ô nhập & Khóa JSON & Cách nhập / mặc định\\\midrule
Điểm cần tính x & \texttt{target} & \texttt{1.5}\\[3pt]
Cấp đạo hàm (0 = giá trị hàm) & \texttt{order} & \texttt{0}
\par Khi cấp > 0,\allowbreak{} lời giải trình bày cả giá trị hàm và đạo hàm yêu cầu trên cùng tập mốc.\\[3pt]
\bottomrule\end{tabular}\end{center}
\subsubsection*{Ô riêng: Lagrange /\allowbreak{} Gauss I /\allowbreak{} Gauss II /\allowbreak{} Stirling /\allowbreak{} Bessel}
\begin{center}\small\begin{tabular}{@{}p{.32\linewidth}p{.17\linewidth}p{.43\linewidth}@{}}\toprule
Ô nhập & Khóa JSON & Cách nhập / mặc định\\\midrule
Số mốc dùng (để trống: toàn bảng) & \texttt{count} & Để trống
\par Stirling: số lẻ;\allowbreak{} Bessel: số chẵn;\allowbreak{} Gauss dùng được cả hai. Ví dụ nhập 9 để chọn 9 mốc gần điểm tính từ bảng lớn.\\[3pt]
\bottomrule\end{tabular}\end{center}
\subsubsection*{Ô riêng: Newton}
\begin{center}\small\begin{tabular}{@{}p{.32\linewidth}p{.17\linewidth}p{.43\linewidth}@{}}\toprule
Ô nhập & Khóa JSON & Cách nhập / mặc định\\\midrule
Loại mốc & \texttt{newton\_nodes} & \texttt{Mốc cách đều}
\par Mốc cách đều dùng sai phân tiến;\allowbreak{} mốc bất kỳ dùng tỷ sai phân.\\[3pt]
Hướng nội suy & \texttt{direction} & \texttt{Tiến}\\[3pt]
Số mốc dùng (để trống: toàn bảng) & \texttt{count} & Để trống
\par Nhập điểm tính khi chọn ít mốc hơn. Không tự thay bậc đề yêu cầu.\\[3pt]
\bottomrule\end{tabular}\end{center}

\paragraph{Input dữ liệu.} Bảng $m=n+1$ cặp phải giữ đúng thứ tự ghép $x_i,y_i$. Tool sắp theo $x$ cùng cả cặp. Hoành độ phân biệt; Lagrange/Newton bất kỳ không cần đều, các công thức sai phân đều cần bước $h$ chung. Ô \emph{Số mốc dùng} trống dùng cả bảng; nhập $K$ thì chọn đúng $K$ mốc và công khai nhóm được chọn. $K$ là số mốc, không phải bậc.

\subsection{Bảng tổng quát phải lập trước khi chọn hệ số}
Lập từ trái sang phải, hết một cấp rồi sang cấp sau:
\begin{center}\small\begin{tabular}{clllll}\toprule
Dòng&$x_i$&$y_i$&Cấp 1&Cấp 2&$\cdots$, cấp $n$\\\midrule
0&$x_0$&$y_0$&$A_{0,1}$&$A_{0,2}$&$\cdots,A_{0,n}$\\
1&$x_1$&$y_1$&$A_{1,1}$&$A_{1,2}$&\\
$i$&$x_i$&$y_i$&$A_{i,1}$ nếu $i<n$&$A_{i,2}$ nếu $i<n-1$&\\
$n$&$x_n$&$y_n$&&&\\\bottomrule
\end{tabular}\end{center}
\textbf{Với mốc đều}: $A_{i,0}=y_i$, $A_{i,k}=A_{i+1,k-1}-A_{i,k-1}$. Chẳng hạn
\[
A_{0,1}=y_1-y_0,\quad A_{1,1}=y_2-y_1,\quad
A_{0,2}=(y_2-y_1)-(y_1-y_0)=y_2-2y_1+y_0.
\]
\textbf{Với mốc bất kỳ}: dùng bảng cùng khuôn nhưng ký hiệu $D$:
\[
D_{i,k}=\frac{D_{i+1,k-1}-D_{i,k-1}}{x_{i+k}-x_i},\quad
D_{0,1}=\frac{y_1-y_0}{x_1-x_0},\quad
D_{0,2}=\frac{\frac{y_2-y_1}{x_2-x_1}-\frac{y_1-y_0}{x_1-x_0}}{x_2-x_0}.
\]
Ô ngoài tam giác không tồn tại, để trống; ô đã tính bằng 0 phải ghi 0. Output của bước lập bảng là nguồn hệ số cho những bước tiếp theo, chưa phải giá trị $P(z)$.

\subsection{Lagrange: bảng tổng quát và đường đi của một hàng}
\textbf{Input riêng}: chọn Lagrange; không có hướng tiến/lùi. Với mỗi mốc $i$, làm đủ một hàng:
\begin{center}\small\begin{tabular}{cllll}\toprule
$i$&Mẫu $d_i$&Hệ số $C_i$&Tích $W_i(z)$&Đóng góp $v_i$\\\midrule
0&$\prod_{j\ne0}(x_0-x_j)$&$y_0/d_0$&$\prod_{j\ne0}(z-x_j)$&$C_0W_0(z)$\\
$i$&$\prod_{j\ne i}(x_i-x_j)$&$y_i/d_i$&$\prod_{j\ne i}(z-x_j)$&$C_iW_i(z)$\\
$n$&$\prod_{j\ne n}(x_n-x_j)$&$y_n/d_n$&$\prod_{j\ne n}(z-x_j)$&$C_nW_n(z)$\\\bottomrule
\end{tabular}\end{center}
Tích ở hàng $i$ viết rõ bằng
\[
W_i(x)=(x-x_0)\cdots(x-x_{i-1})(x-x_{i+1})\cdots(x-x_n),
\]
bỏ đúng mốc $i$. Tính mẫu $d_i=W_i(x_i)$ $\to$ chia $C_i=y_i/d_i$ $\to$ thế $z$ vào tích $\to$ cộng
\[
P(z)=\sum_{i=0}^nC_iW_i(z)=C_0W_0(z)+C_1W_1(z)+\cdots+C_nW_n(z).
\]
\textbf{Output}: đa thức qua tập mốc, bảng các đóng góp, $P(z)$ nếu nhập điểm và $P^{(s)}(z)$ nếu cấp $s>0$. Khi $z$ bằng một mốc, giá trị phải khớp tung độ; không tính dạng $0/0$.

\subsection{Newton: tách loại mốc và hướng trước khi lấy hệ số}
\textbf{Input riêng}: ô \emph{Loại mốc} và \emph{Hướng nội suy}. Bảng tổng quát dưới đây định nghĩa từng hạng $c_kB_k$; hàng tổng luôn là $P=\sum c_kB_k=c_0B_0+c_1B_1+\cdots+c_nB_n$.
\begin{center}\small\begin{tabular}{@{}p{.20\linewidth}p{.21\linewidth}p{.49\linewidth}@{}}\toprule
Lựa chọn&$c_k$ lấy ở đâu&$B_k$ dùng để nhân\\\midrule
Bất kỳ, tiến&$D_{0,k}$&$B_0=1$; $B_k=(x-x_0)(x-x_1)\cdots(x-x_{k-1})$\\
Bất kỳ, lùi&$D_{n-k,k}$&$B_0=1$; $B_k=(x-x_n)(x-x_{n-1})\cdots(x-x_{n-k+1})$\\
Đều, tiến&$A_{0,k}/k!$&$t=(x-x_0)/h$; $B_k=t(t-1)\cdots(t-k+1)$\\
Đều, lùi&$A_{n-k,k}/k!$&$t=(x-x_n)/h$; $B_k=t(t+1)\cdots(t+k-1)$\\\bottomrule
\end{tabular}\end{center}
\textbf{Quá trình}: bảng $A/D$ $\to$ lấy đúng đường hệ số $\to$ dựng tích theo đúng hướng $\to$ tính đóng góp $c_kB_k$ $\to$ cộng. Với mốc bất kỳ không chia thêm $k!$. \textbf{Output}: cùng loại với Lagrange, kèm gốc đổi biến và hướng đang dùng. Nếu dùng cùng mốc đủ số hạng, các lựa chọn biểu diễn cùng đa thức.

\subsection{Gauss I, Gauss II, Stirling, Bessel: mỗi hạng cần ghi đủ ba thứ}
\textbf{Input}: chọn đúng tên công thức; kiểm lưới đều. Stirling dùng số mốc lẻ; Bessel số mốc chẵn. Gauss dùng được cả hai. Ghi gốc $a$ và $t=(z-a)/h$ trước khi điền bảng.
\begin{center}\small\begin{tabular}{cllll}\toprule
Cấp $k$&Ô sai phân đã lấy $\sigma_k$&$c_k=\sigma_k/k!$&Tích $B_k(t)$&Đóng góp $c_kB_k(t)$\\\midrule
0&Ô/giá trị giữa theo loại&$c_0$&1&$c_0$\\
1&Nguồn cấp 1&$c_1$&$B_1(t)$&$c_1B_1(t)$\\
$k$&Nguồn cấp $k$&$c_k$&$B_k(t)$&$c_kB_k(t)$\\
$n$&Nguồn cấp cuối&$c_n$&$B_n(t)$&$c_nB_n(t)$\\\bottomrule
\end{tabular}\end{center}
Không ghi chỉ một dãy hệ số mà thiếu nguồn ô hoặc thiếu tích. Quy tắc lấy nguồn và các tích đầy đủ nằm ở các mục từng phương pháp bên dưới. Các hạng đầu được khai triển để nhận dạng:
\begin{align*}
\text{Gauss I: }&y_0+t\Delta y_0+\tfrac{t(t-1)}2\Delta^2y_{-1}
 +\tfrac{t(t-1)(t+1)}6\Delta^3y_{-1}+\cdots,\\
\text{Gauss II: }&y_0+t\Delta y_{-1}+\tfrac{t(t+1)}2\Delta^2y_{-1}
 +\tfrac{t(t+1)(t-1)}6\Delta^3y_{-2}+\cdots,\\
\text{Stirling: }&y_0+t\tfrac{\Delta y_{-1}+\Delta y_0}2+\tfrac{t^2}2\Delta^2y_{-1}\\
 &\quad+\tfrac{t(t^2-1)}6\tfrac{\Delta^3y_{-2}+\Delta^3y_{-1}}2
 +\tfrac{t^2(t^2-1)}{24}\Delta^4y_{-2}+\cdots,\\
\text{Bessel: }&\tfrac{y_0+y_1}2+(t-\tfrac12)\Delta y_0
 +\tfrac{t(t-1)}2\tfrac{\Delta^2y_{-1}+\Delta^2y_0}2\\
 &\quad+\tfrac{(t-\tfrac12)t(t-1)}6\Delta^3y_{-1}+\cdots.
\end{align*}
\textbf{Output}: bảng mốc đã chọn, bảng sai phân, nguồn/hệ số/tích từng hạng, $P(z)$ và đạo hàm yêu cầu. Khi lấy đạo hàm theo $t$, phải đổi $P^{(s)}(z)=Q^{(s)}(t)/h^s$. Một hạng có giá trị 0 tại điểm hỏi vẫn có thể có đạo hàm khác 0.

\subsection{Tool Chọn mốc: output là dữ liệu để đưa sang tool tiếp theo}
\subsection*{Input: các ô đúng tên trên giao diện}
Bảng này được lấy từ biểu mẫu của tool. Giá trị mặc định là ví dụ nhập, cần thay bằng dữ liệu đề. Hai dãy $x,y$ là bảng dữ liệu riêng bên dưới nếu chức năng dùng bảng.
\subsubsection*{Ô riêng: Mốc cách đều}
\begin{center}\small\begin{tabular}{@{}p{.32\linewidth}p{.17\linewidth}p{.43\linewidth}@{}}\toprule
Ô nhập & Khóa JSON & Cách nhập / mặc định\\\midrule
Điểm x & \texttt{target} & \texttt{1.5}\\[3pt]
Số mốc cần lấy & \texttt{count} & \texttt{3}\\[3pt]
\bottomrule\end{tabular}\end{center}
\subsubsection*{Ô riêng: Mốc nội suy ngược}
\begin{center}\small\begin{tabular}{@{}p{.32\linewidth}p{.17\linewidth}p{.43\linewidth}@{}}\toprule
Ô nhập & Khóa JSON & Cách nhập / mặc định\\\midrule
Giá trị y & \texttt{target} & \texttt{1.5}\\[3pt]
\bottomrule\end{tabular}\end{center}
\subsubsection*{Ô riêng: Chebyshev}
\begin{center}\small\begin{tabular}{@{}p{.32\linewidth}p{.17\linewidth}p{.43\linewidth}@{}}\toprule
Ô nhập & Khóa JSON & Cách nhập / mặc định\\\midrule
Bậc đa thức n (sinh n+1 mốc) & \texttt{degree} & \texttt{4}\\[3pt]
Đầu đoạn a & \texttt{a} & \texttt{-1}\\[3pt]
Cuối đoạn b & \texttt{b} & \texttt{1}\\[3pt]
\bottomrule\end{tabular}\end{center}

\textbf{Mốc cách đều}: input toàn bảng, điểm $z$, số mốc $K$ $\to$ xét các cửa sổ liên tiếp $j,\ldots,j+K-1$ $\to$ tính $c_j=(x_j+x_{j+K-1})/2$, $d_j=|c_j-z|$ $\to$ lấy cửa sổ gần tâm nhất, hòa lấy trái. Trên mốc bất kỳ bộ chọn có thể dùng nhóm gần điểm rồi sắp lại. Output là $K$ cặp, không phải $P(z)$.

\textbf{Mốc nội suy ngược}: input bảng và mức $\bar y$ $\to$ tìm cặp kề nhau có tung độ khác nhau và kẹp mức $\to$ mở rộng trái/phải khi các hiệu tung độ giữ cùng dấu $\to$ xuất từng bảng \emph{đoạn đơn điệu chứa mức}, bỏ các đoạn trùng. Output là các tập mốc ứng viên để chọn nhánh, chưa phải nghiệm nội suy ngược. Một đoạn phẳng bằng mức cần khảo sát riêng; tool chọn mốc này không thay chức năng quét mức đầy đủ trong bộ giải ngược.

\textbf{Chebyshev}: input bậc $n$, đoạn $[a,b]$ $\to$ sinh $m=n+1$ mốc
\[
x_k=\tfrac{a+b}2+\tfrac{b-a}2\cos\tfrac{(2k+1)\pi}{2m},\quad k=0,\ldots,m-1.
\]
Output là dãy $x$ sắp tăng. Cần tự có $y_k=f(x_k)$ rồi mới nội suy; không thể lấy tùy tiện tung độ của bảng cũ gắn vào các mốc mới.

\subsection{Tool Horner: input hệ số, output giá trị và các đạo hàm}
\subsection*{Input: các ô đúng tên trên giao diện}
Bảng này được lấy từ biểu mẫu của tool. Giá trị mặc định là ví dụ nhập, cần thay bằng dữ liệu đề. Hai dãy $x,y$ là bảng dữ liệu riêng bên dưới nếu chức năng dùng bảng.
\subsubsection*{Ô chung của các lựa chọn trong nhóm}
\begin{center}\small\begin{tabular}{@{}p{.32\linewidth}p{.17\linewidth}p{.43\linewidth}@{}}\toprule
Ô nhập & Khóa JSON & Cách nhập / mặc định\\\midrule
Điểm x & \texttt{target} & \texttt{1.5}\\[3pt]
\bottomrule\end{tabular}\end{center}
\subsubsection*{Ô riêng: Đạo hàm từ bảng}
\begin{center}\small\begin{tabular}{@{}p{.32\linewidth}p{.17\linewidth}p{.43\linewidth}@{}}\toprule
Ô nhập & Khóa JSON & Cách nhập / mặc định\\\midrule
Bậc p (dùng p+1 mốc) & \texttt{degree} & \texttt{2}\\[3pt]
Cấp đạo hàm & \texttt{order} & \texttt{1}\\[3pt]
\bottomrule\end{tabular}\end{center}
\subsubsection*{Ô riêng: Horner}
\begin{center}\small\begin{tabular}{@{}p{.32\linewidth}p{.17\linewidth}p{.43\linewidth}@{}}\toprule
Ô nhập & Khóa JSON & Cách nhập / mặc định\\\midrule
Hệ số theo bậc giảm dần & \texttt{coef} & \texttt{1 0 0}
\par Ví dụ 1 0 0 biểu diễn x².\\[3pt]
Đạo hàm đến cấp & \texttt{order} & \texttt{2}\\[3pt]
\bottomrule\end{tabular}\end{center}

\textbf{Horner nhận hệ số giảm dần}: $P=a_0x^n+a_1x^{n-1}+\cdots+a_n$, điểm $c$, cấp đạo hàm đến $s$. Giữ hệ số 0 giữa/cuối. Mẫu bảng ngang tổng quát:
\begin{center}\small\setlength{\tabcolsep}{10pt}\begin{tabular}{@{}lrrrr|r@{}}\toprule
Điểm $c$&Bậc $n$&Bậc $n-1$&$\cdots$&Bậc 1&Dư\\\midrule
Hệ số&$a_0$&$a_1$&$\cdots$&$a_{n-1}$&$a_n$\\
Nhân, chuyển phải&---&$cb_0$&$\cdots$&$cb_{n-2}$&$cb_{n-1}$\\\midrule
Cộng xuống&$b_0$&$b_1$&$\cdots$&$b_{n-1}$&$b_n$\\\bottomrule
\end{tabular}\end{center}
Tính từ trái sang phải: $b_0=a_0$, $b_1=a_1+cb_0$, $b_2=a_2+cb_1$, $\ldots$, $b_n=a_n+cb_{n-1}$. Hàng Nhân lấy ô \emph{Cộng xuống của cột trước} nhân $c$ rồi chuyển sang cột kế; sau đó cộng với hệ số phía trên.
Khuôn trên minh họa nhiều cột. Đa thức bậc 1 chỉ có hai cột; đa thức hằng chỉ hạ hệ số, không tạo phép nhân thừa. Cấp đạo hàm lớn hơn bậc thực cho 0.
Ô cuối $b_n=P(c)$; các ô trước là hệ số thương. Đưa \emph{chỉ thương} vào lượt sau. Nếu các số dư lần lượt $s_0,s_1,\ldots$, output đạo hàm là
\[
P(c)=s_0,\quad P'(c)=s_1,\quad P''(c)=2s_2,\quad P'''(c)=6s_3,\quad P^{(r)}(c)=r!s_r.
\]
Các mục tiếp theo giải thích công thức và đưa ví dụ số trên đúng các khuôn bảng này. Tool \emph{Đạo hàm từ bảng} cùng nhóm dùng cặp $x,y$, được hướng dẫn riêng ở Chương \ref{ch:calculus}.


\section{Đọc bài toán trước khi bấm Tính}

Cho bảng $m$ cặp $(x_i,y_i)$, ta dựng một đa thức $P$ đi qua các cặp được chọn rồi dùng $P$ để tính giá trị hoặc đạo hàm. Đặt $n=m-1$: bậc của $P$ \emph{không quá} $n$, không nhất thiết bằng $n$. Các hệ số bậc cao có thể bằng 0.

\begin{center}\small
\begin{tabular}{@{}p{.25\linewidth}p{.32\linewidth}p{.31\linewidth}@{}}
\toprule
Đề cho/hỏi & Lựa chọn trong tool & Kiểm tra trước khi tính\\
\midrule
Mốc bất kỳ, cần $P(z)$ & Lagrange hoặc Newton, loại mốc bất kỳ & $x_i$ phân biệt; giữ đúng cặp $x,y$\\
Mốc đều, gần đầu/cuối bảng & Newton tiến/lùi & Có bước $h$ chung; gốc đúng đầu đã chọn\\
Mốc đều quanh một mốc giữa & Gauss I, Gauss II, Stirling & Stirling cần số mốc lẻ\\
Mốc đều quanh giữa hai mốc & Bessel & Số mốc chẵn, ít nhất 2\\
Đề đã cho đa thức & Đạo hàm \& Horner $\to$ Horner & Nhập đủ các hệ số theo bậc giảm\\
Đề hỏi $f^{(s)}(z)$ qua nội suy & Giữ đúng $P$ vừa dựng & Đạo hàm theo $t$ phải chia $h^s$\\
\bottomrule
\end{tabular}\end{center}

\paragraph{Điều kiện toán học.} Với $m$ hoành độ phân biệt có đúng một đa thức bậc không quá $m-1$ nội suy các cặp. Nếu hai đa thức cùng thỏa điều kiện thì hiệu của chúng có $m$ nghiệm phân biệt nhưng bậc không quá $m-1$, nên hiệu bằng 0. Vì vậy Lagrange, các dạng Newton, Gauss, Stirling và Bessel cho cùng một $P$ nếu dùng \textbf{cùng tập mốc và đủ số hạng}. Thay tên công thức không tự làm sai số nội suy nhỏ hơn.

\paragraph{Không có vòng hội tụ.} Lập bảng sai phân, nhân các tích và cộng các hạng là những phép tính hữu hạn. Số bước do số mốc quyết định; không cần $\varepsilon$ để dừng. Vấn đề sai số là độ chính xác của dữ liệu, độ nhạy của nội suy và số học máy, không phải ``lặp chưa hội tụ''.

\subsection{Các ô nhập chung}
\begin{enumerate}
\item Nhập hai dãy $x,y$ cùng độ dài. Dấu thập phân là dấu chấm. Tool sắp \emph{cả cặp} theo $x$ tăng; các hoành độ trùng phải được xử lý ở dữ liệu nguồn.
\item \textbf{Điểm cần tính} là $z$. Để trống thì chỉ dựng đa thức; nếu yêu cầu chọn ít mốc hơn bảng nhập thì phải có $z$.
\item \textbf{Cấp đạo hàm}: $s=0$ là giá trị, $s=1$ là đạo hàm một, v.v. Khi $s>0$, báo cáo mới có cả $P(z)$ và $P^{(s)}(z)$ trên cùng tập mốc.
\item \textbf{Số mốc dùng}: để trống nghĩa là toàn bộ bảng. Nhập $K$ nghĩa là chọn đúng $K$ cặp, bậc tối đa $K-1$. Đề chỉ định rõ các mốc thì nhập đúng các cặp đó.
\item Với Newton, còn phải chọn \textbf{mốc cách đều/mốc bất kỳ} và \textbf{tiến/lùi}. Hai lựa chọn này quyết định cả bảng hệ số và tích cơ sở.
\end{enumerate}

Một bảng 100 mốc không tự dùng được với Stirling toàn bảng vì 100 là số chẵn. Có thể chọn công khai 9 mốc theo đề; tool không tự bỏ một mốc để tạo 99 mốc. Một mốc cho đa thức hằng; Bessel cần ít nhất hai mốc.

\section{Chọn mốc: biết số nào thực sự được đưa vào đa thức}
\label{sec:node-selection}

\subsection{Cửa sổ liên tiếp trên lưới đều}
\textbf{Đầu vào}: toàn bảng, điểm $z$ và số mốc $K$. Với mỗi cửa sổ $K$ mốc liên tiếp, tính tâm hai đầu rồi so khoảng cách đến $z$:
\[
 c_j=\frac{x_j+x_{j+K-1}}2,\qquad d_j=|c_j-z|.
\]
Chọn cửa sổ có $d_j$ nhỏ nhất; hòa thì lấy cửa sổ bên trái. Kết quả vẫn là lưới đều, có đủ $K$ mốc và không làm lệch cặp.

\textbf{Ví dụ tự dựng}: $x=0,1,\ldots,6$, $y=x^2$, $z=3.2$, $K=3$.
\begin{center}\small
\begin{tabular}{@{}crrr@{}}\toprule
Cửa sổ & Tâm & Khoảng cách đến $3.2$ & Kết luận\\\midrule
$[0,1,2]$ & 1 & 2.2 & bỏ\\
$[1,2,3]$ & 2 & 1.2 & bỏ\\
$[2,3,4]$ & 3 & 0.2 & chọn\\
$[3,4,5]$ & 4 & 0.8 & bỏ\\
$[4,5,6]$ & 5 & 1.8 & bỏ\\\bottomrule
\end{tabular}\end{center}
Đưa các cặp $(2,4),(3,9),(4,16)$ sang thuật toán nội suy. Bảng gốc có 7 cặp nhưng đa thức của lượt này chỉ dựng từ 3 cặp.

\subsection{Mốc bất kỳ và điểm nằm sát biên}
Nếu lưới không đều, tool xếp các mốc theo $|x_i-z|$, hòa theo thứ tự $x$ tăng; lấy $K$ mốc đầu rồi sắp tập chọn tăng dần. Ví dụ $x=[0,1,3,8]$, $z=2.2$, $K=3$ có các khoảng cách $[2.2,1.2,0.8,5.8]$, nên chọn $x=[0,1,3]$. Tập này dùng Lagrange hoặc Newton tỉ sai phân, không dùng công thức sai phân đều.

Ở đầu bảng, cửa sổ có thể lệch về một phía. Không yêu cầu thêm mốc nằm ngoài tập đã chọn. Nếu $z$ nằm ngoài đoạn của tập mốc dùng thì đang \textbf{ngoại suy}; đa thức vẫn đánh giá được nhưng không có bảo đảm chính xác của hàm thật chỉ từ bảng.

\subsection{Quét mốc cho mức cần tìm ngược}
Lựa chọn \textbf{Chọn mốc $\to$ Mốc nội suy ngược} nhận mức $\bar y$. Tool tìm các cặp kề có mức nằm giữa hai tung độ, mở rộng theo dấu của sai phân cấp một và liệt kê các đoạn đơn điệu trên \emph{mẫu}. Ví dụ $x=[0,1,2,3,4]$, $y=[0,1,2,1,0]$, mức $1.5$ tạo hai nhánh mẫu $[0,2]$ và $[2,4]$. Phải chọn nhánh của bài toán trước khi giải ngược.

Đơn điệu trên các mẫu chưa chứng minh hàm thật đơn điệu giữa mọi mốc. Quét đổi dấu có thể bỏ sót nghiệm tiếp xúc. Phần này chọn dữ liệu; chưa tính nghiệm.

\subsection{Sinh mốc Chebyshev khi được quyền lấy mẫu mới}
\label{sec:chebyshev}
\textbf{Đầu vào}: bậc $n\ge0$, đoạn $a<b$. Tool sinh $m=n+1$ mốc nghiệm:
\[
 \theta_i=\frac{(2i+1)\pi}{2m},\qquad
 x_i=\frac{a+b}{2}+\frac{b-a}{2}\cos\theta_i,
 \qquad i=0,\ldots,m-1.
\]
\begin{enumerate}
\item Tính $m=n+1$, trung điểm và nửa độ dài đoạn.
\item Với từng $i$, tính góc $\theta_i$, cosin, rồi $x_i$.
\item Sắp các $x_i$ tăng. Kiểm tra chúng còn phân biệt trong biểu diễn máy.
\item Tính hoặc đo \emph{tung độ mới} tại các mốc mới; lúc đó mới nội suy.
\end{enumerate}
Ví dụ $n=2$, $[a,b]=[-1,1]$:
\begin{center}\begin{tabular}{@{}crr@{}}\toprule
$i$ & $\theta_i$ & $x_i$ trước sắp\\\midrule
0 & $\pi/6$ & $\sqrt3/2$\\
1 & $\pi/2$ & $0$\\
2 & $5\pi/6$ & $-\sqrt3/2$\\\bottomrule
\end{tabular}\end{center}
Đầu ra tăng là $[-\sqrt3/2,0,\sqrt3/2]$. Đây là mốc \emph{nghiệm}, không chứa hai đầu đoạn; khác mốc cực trị Chebyshev. Không được ghép lại các $y$ cũ vào các $x$ mới.

Nếu $|f^{(m)}|\le M_m$ trên $[a,b]$, dùng đúng $m$ nghiệm trên thì
\[
 \|f-P\|_\infty\le
 \frac{M_m(b-a)^m}{2^{2m-1}m!}.
\]
Tính tối ưu ở đây là chuẩn của tích mốc monic; không khẳng định sai số thực nhỏ nhất cho mọi hàm. Công cụ Python \texttt{chebyshev\_value} tính $T_n$ qua $T_0=1,T_1=u,T_{k+1}=2uT_k-T_{k-1}$; đó là phép tính hữu hạn khác với việc sinh các mốc.

\section{Hai bảng phải phân biệt: sai phân và tỉ sai phân}
\label{sec:difference-tables}
Giữ chỉ số mốc bắt đầu của mỗi ô. Ô ngoài tam giác không được tính, để trống chứ không ghi 0.

\subsection{Sai phân tiến cho mốc cách đều}
\[
 A_{i,0}=y_i,\qquad
 A_{i,k}=A_{i+1,k-1}-A_{i,k-1}=\Delta^ky_i.
\]
\textbf{Thứ tự}: tính hết cột cấp 1 từ $y$; sau đó hết cấp 2 từ cấp 1; tiếp tục đến cấp $n$. Cấp $k$ có $m-k$ phép trừ.

\textbf{Bảng mẫu dùng xuyên suốt các công thức trung tâm}:
\[
 x=[-2,-1,0,1,2],\qquad y=[17,1,1,5,25].
\]
Chọn mốc giữa làm $x_0=0$, đánh chỉ số $j=-2,-1,0,1,2$.
\begin{center}\small
\begin{tabular}{@{}rrrrrrr@{}}\toprule
$j$ & $x_j$ & $y_j$ & $\Delta y_j$ & $\Delta^2y_j$ & $\Delta^3y_j$ & $\Delta^4y_j$\\\midrule
$-2$ & $-2$ & 17 & $-16$ & 16 & $-12$ & 24\\
$-1$ & $-1$ & 1 & 0 & 4 & 12 & \\
0 & 0 & 1 & 4 & 16 & & \\
1 & 1 & 5 & 20 & & & \\
2 & 2 & 25 & & & & \\\bottomrule
\end{tabular}\end{center}
Các phép tính thực hiện theo thứ tự:
\begin{align*}
\Delta^1 &: 1-17=-16,\quad1-1=0,\quad5-1=4,\quad25-5=20;\\
\Delta^2 &: 0-(-16)=16,\quad4-0=4,\quad20-4=16;\\
\Delta^3 &: 4-16=-12,\quad16-4=12;\\
\Delta^4 &: 12-(-12)=24.
\end{align*}
Một ô $\Delta^2y_{-1}=4$ có nghĩa là bắt đầu ở $j=-1$ và dùng ba tung độ $y_{-1},y_0,y_1$.

\subsection{Tỉ sai phân cho mốc bất kỳ}
\[
 D_{i,0}=y_i,\qquad
 D_{i,k}=\frac{D_{i+1,k-1}-D_{i,k-1}}{x_{i+k}-x_i}
        =f[x_i,\ldots,x_{i+k}].
\]
Tính như sai phân nhưng chia cho khoảng giữa hai mốc ngoài cùng của ô. Không chia thêm $k!$. Với lưới đều mới có đẳng thức $D_{i,k}=\Delta^ky_i/(k!h^k)$.

Ví dụ $x=[0,1,3]$, $y=[1,3,13]$:
\begin{center}\begin{tabular}{@{}rrrr@{}}\toprule
$x_i$ & $y_i$ & $D_{i,1}$ & $D_{i,2}$\\\midrule
0 & 1 & 2 & 1\\
1 & 3 & 5 & \\
3 & 13 & & \\\bottomrule
\end{tabular}\end{center}
Ở đây $(3-1)/(1-0)=2$, $(13-3)/(3-1)=5$, rồi $(5-2)/(3-0)=1$. Toàn bộ hệ số cần dùng đã có trong bảng.

\section{Lagrange: dựng từng đa thức cơ sở}
\label{sec:lagrange}
\textbf{Đầu vào}: các $x_i$ phân biệt, $y_i$, điểm $z$ nếu cần đánh giá; không cần lưới đều. Đặt
\[
 W_i(x)=\prod_{j\ne i}(x-x_j),\quad
 d_i=\prod_{j\ne i}(x_i-x_j),\quad
 L_i(x)=\frac{W_i(x)}{d_i},\quad
 P(x)=\sum_{i=0}^ny_iL_i(x).
\]
Vì $L_i(x_j)=\delta_{ij}$ nên từng hạng chỉ giữ đúng tung độ của mốc tương ứng.

\subsection{Các bước và số chuyển sang bước sau}
\begin{enumerate}
\item Với từng $i$, bỏ mốc $x_i$ khỏi tích để lập $W_i$. Tử số này là một đa thức bậc $n$.
\item Thế $x=x_i$ vào $W_i$ để có mẫu $d_i\ne0$; giữ dấu các hiệu.
\item Tính $C_i=y_i/d_i$. Hạng của đa thức là $C_iW_i(x)$.
\item Cộng tất cả các hạng để có $P$. Có thể giữ dạng tích, không cần khai triển theo $x$.
\item Tại $z$, tính từng $W_i(z)$ rồi từng đóng góp $C_iW_i(z)$ và cộng. Nếu $z=x_i$, dùng tích bỏ mốc trực tiếp, không tính biểu thức $W(z)/(z-x_i)$ dạng $0/0$.
\end{enumerate}

\subsection{Ví dụ đầy đủ trên ba mốc không đều}
Dùng bảng $[0,1,3]$, $[1,3,13]$ để tính $P(2)$.
\begin{center}\small
\begin{tabular}{@{}crrlr@{}}\toprule
$i$ & $d_i$ & $C_i$ & $W_i(x)$ & $C_iW_i(2)$\\\midrule
0 & 3 & $1/3$ & $(x-1)(x-3)$ & $-1/3$\\
1 & $-2$ & $-3/2$ & $x(x-3)$ & $3$\\
2 & 6 & $13/6$ & $x(x-1)$ & $13/3$\\\bottomrule
\end{tabular}\end{center}
Mẫu $d_1=(1-0)(1-3)=-2$ phải âm. Cộng các hạng:
\[
 P(x)=\frac{(x-1)(x-3)}3-\frac32x(x-3)+\frac{13}6x(x-1)
     =x^2+x+1.
\]
Tại $z=2$: $P(2)=-1/3+3+13/3=7$. Kiểm tra $P(0)=1$, $P(1)=3$, $P(3)=13$ và $\sum_iL_i(2)=1$.

\paragraph{Bảng tích/thương.} Một cách tương đương là lập $W(x)=\prod_j(x-x_j)$ rồi chia Horner cho từng $(x-x_i)$ để lấy $W_i$. Số dư phải bằng 0. Hệ số của bảng tích phải được đọc đúng nhãn bậc tăng/giảm; không coi một dãy hệ số chưa định nghĩa thứ tự là công thức hoàn chỉnh.

\section{Newton với mốc bất kỳ: tiến và lùi}
\label{sec:newton-divided}
\subsection{Tiến: lấy đường hệ số đầu bảng}
Sau khi lập bảng $D$, lấy $c_k=D_{0,k}$. Đặt $B_0=1$, $B_{k+1}(x)=B_k(x)(x-x_k)$:
\[
 P(x)=\sum_{k=0}^nc_kB_k(x)
 =y_0+f[x_0,x_1](x-x_0)+\cdots.
\]
Mỗi bước $k$ lấy một ô hệ số, nhân tích đang có, cộng vào tổng; tích cho bước sau được nhân thêm $(x-x_k)$.

Với ví dụ ba mốc, $c=[1,2,1]$ nên
\[
 P(x)=1+2x+x(x-1).
\]
Tại 2, các đóng góp là $1,4,2$, tổng 7. Ở hạng cuối, hệ số là $1$, không phải $1/2!$.

\subsection{Lùi: lấy hệ số theo đường chéo cuối bảng}
Từ cùng bảng $D$, lấy $\tilde c_k=D_{n-k,k}$ và
\[
 \tilde B_0=1,\qquad
 \tilde B_k(x)=\prod_{j=0}^{k-1}(x-x_{n-j}),\qquad
 P(x)=\sum_{k=0}^n\tilde c_k\tilde B_k(x).
\]
Thứ tự mốc trong các tích là từ phải sang trái. Với bảng mẫu, $\tilde c=[13,5,1]$:
\[
 P(x)=13+5(x-3)+(x-3)(x-1).
\]
Tại 2, các đóng góp $13,-5,-1$ cho tổng 7. Hệ số lùi phải đi cùng tích lùi; không trộn với đường hệ số tiến.

\subsection{Horner Newton: đánh giá dạng lồng}
Viết theo đúng thứ tự mốc $\xi_0,\xi_1,\ldots$:
\[
 P(x)=c_0+(x-\xi_0)\{c_1+(x-\xi_1)[\cdots+(x-\xi_{n-1})c_n]\}.
\]
\begin{enumerate}
\item Bắt đầu $v=c_n$.
\item Với $k=n-1,n-2,\ldots,0$, tính $d=z-\xi_k$, tích $d\,v_{\rm cũ}$, rồi $v_{\rm mới}=c_k+d\,v_{\rm cũ}$.
\item Số cuối cùng là $P(z)$.
\end{enumerate}
Ví dụ tiến tại $z=2$: $v=1$; bước $k=1$ cho $v=2+(2-1)1=3$; bước $k=0$ cho $v=1+(2-0)3=7$. Đây là Horner theo cơ sở Newton, không phải sơ đồ lấy hệ số lũy thừa và nhân cùng một $c$ ở mọi lượt.

\subsection{Thêm một mốc mà không phá thứ tự cũ}
Nếu cần bổ sung $(x_{n+1},y_{n+1})$, tính
\[
 c_{n+1}=\frac{y_{n+1}-P_n(x_{n+1})}{\prod_{j=0}^n(x_{n+1}-\xi_j)}.
\]
Ví dụ từ $(0,1),(1,2)$ có $P_1=1+x$. Thêm $(2,5)$: $P_1(2)=3$, tích mẫu $2(2-1)=2$, nên hệ số mới $(5-3)/2=1$ và $P_2=1+x+x(x-1)=1+x^2$. Mốc mới phải khác mọi mốc cũ. API \texttt{newton\_append} trả các mốc và hệ số đã nối; đây không phải một vòng tìm nghiệm.

\section{Newton trên lưới đều: giai thừa và đổi gốc}
\label{sec:newton-uniform}
Gọi $h=x_{i+1}-x_i>0$. Tool kiểm tra các bước theo thang $h$, không theo trị tuyệt đối của $x$. Trong lời giải, luôn ghi gốc, $h$ và $t$ trước khi lấy hệ số.

\subsection{Tiến từ đầu trái}
\[
 t=\frac{x-x_0}{h},\qquad
 P(x_0+ht)=\sum_{k=0}^n\frac{\Delta^ky_0}{k!}
              \prod_{j=0}^{k-1}(t-j).
\]
Lấy hàng đầu của mỗi cấp; chia $k!$; tích lần lượt nhân $t,t-1,t-2,\ldots$. Dùng bảng năm mốc ở trên, lần này gốc \emph{đầu trái} là $-2$, $h=1$, nên tại $z=1/4$ có $t=9/4$.
\begin{center}\small
\begin{tabular}{@{}rrrrr@{}}\toprule
$k$ & $\Delta^ky_{\rm đầu}$ & $\beta_k=\Delta^k/k!$ & $B_k(9/4)$ & $\beta_kB_k$\\\midrule
0 & 17 & 17 & 1 & 17\\
1 & $-16$ & $-16$ & $9/4$ & $-36$\\
2 & 16 & 8 & $45/16$ & $45/2$\\
3 & $-12$ & $-2$ & $45/64$ & $-45/32$\\
4 & 24 & 1 & $-135/256$ & $-135/256$\\\bottomrule
\end{tabular}\end{center}
Tổng bằng $401/256=1.56640625$. Giá trị $t=2.25$ không làm công thức mất hiệu lực; khoảng $[0,1]$ chỉ là miền tính thuận lợi gần đầu, không phải điều kiện tồn tại của đa thức.

\subsection{Lùi từ đầu phải}
\[
 t=\frac{x-x_n}{h},\qquad
 P(x_n+ht)=\sum_{k=0}^n\frac{\Delta^ky_{n-k}}{k!}
              \prod_{j=0}^{k-1}(t+j).
\]
Đẳng thức $\nabla^ky_n=\Delta^ky_{n-k}$ cho phép đọc bảng $\Delta$ mà không dựng thêm bảng. Gốc đầu phải là 2; tại $z=1/4$, $t=-7/4$.
\begin{center}\small
\begin{tabular}{@{}rrrrr@{}}\toprule
$k$ & $\Delta^ky_{n-k}$ & $\beta_k$ & $B_k(-7/4)$ & $\beta_kB_k$\\\midrule
0 & 25 & 25 & 1 & 25\\
1 & 20 & 20 & $-7/4$ & $-35$\\
2 & 16 & 8 & $21/16$ & $21/2$\\
3 & 12 & 2 & $21/64$ & $21/32$\\
4 & 24 & 1 & $105/256$ & $105/256$\\\bottomrule
\end{tabular}\end{center}
Tổng cũng bằng $401/256$. Dùng $t=-7/4$ nhưng lấy các sai phân của gốc trái sẽ cho một công thức sai.

\section{Gauss I và Gauss II: thứ tự kết nạp quyết định tích}
\label{sec:gauss-central}
Đổi sang chỉ số quanh gốc: $x_j=a+jh$, $t=(x-a)/h$. Đặt $r_0,r_1,\ldots$ là thứ tự các độ lệch mốc được kết nạp, $B_0=1$, $B_{k+1}=B_k(t-r_k)$.
\begin{center}\small
\begin{tabular}{@{}lll@{}}\toprule
 & Gauss I & Gauss II\\\midrule
Thứ tự $r_j$ & $0,1,-1,2,-2,\ldots$ & $0,-1,1,-2,2,\ldots$\\
Nguồn sai phân bậc $k$ & $\Delta^ky_{-\lfloor k/2\rfloor}$ & $\Delta^ky_{-\lceil k/2\rceil}$\\
Hệ số nhân $B_k$ & sai phân chia $k!$ & sai phân chia $k!$\\\bottomrule
\end{tabular}\end{center}
Giữ các mốc có trong tập. Với $m$ mốc tăng, vị trí gốc trong mảng là $c=\lfloor(m-1)/2\rfloor$ cho Gauss I, $c=\lfloor m/2\rfloor$ cho Gauss II. Khi $m$ lẻ, hai gốc trùng nhau; khi $m$ chẵn, chúng là hai mốc giữa khác nhau. Công thức đầy đủ:
\[
 Q_{\rm I}(t)=\sum_{k=0}^n\frac{\Delta^ky_{-\lfloor k/2\rfloor}}{k!}B_k(t),\qquad
 Q_{\rm II}(t)=\sum_{k=0}^n\frac{\Delta^ky_{-\lceil k/2\rceil}}{k!}B_k(t).
\]

\subsection{Năm mốc: theo dõi từng nguồn và từng tích}
Với bảng mẫu năm mốc, cả hai có $a=0,h=1,t_z=1/4$.
\begin{center}\small
\begin{tabular}{@{}rllrr@{}}\toprule
$k$ & Nguồn Gauss I & $B_k(t)$ & $\beta_k$ & đóng góp tại $1/4$\\\midrule
0 & $y_0=1$ & $1$ & 1 & 1\\
1 & $\Delta y_0=4$ & $t$ & 4 & 1\\
2 & $\Delta^2y_{-1}=4$ & $t(t-1)$ & 2 & $-3/8$\\
3 & $\Delta^3y_{-1}=12$ & $t(t-1)(t+1)$ & 2 & $-15/32$\\
4 & $\Delta^4y_{-2}=24$ & $t(t-1)(t+1)(t-2)$ & 1 & $105/256$\\\bottomrule
\end{tabular}\end{center}
Tổng $1+1-3/8-15/32+105/256=401/256$.
\begin{center}\small
\begin{tabular}{@{}rllrr@{}}\toprule
$k$ & Nguồn Gauss II & $B_k(t)$ & $\beta_k$ & đóng góp tại $1/4$\\\midrule
0 & $y_0=1$ & 1 & 1 & 1\\
1 & $\Delta y_{-1}=0$ & $t$ & 0 & 0\\
2 & $\Delta^2y_{-1}=4$ & $t(t+1)$ & 2 & $5/8$\\
3 & $\Delta^3y_{-2}=-12$ & $t(t+1)(t-1)$ & $-2$ & $15/32$\\
4 & $\Delta^4y_{-2}=24$ & $t(t+1)(t-1)(t+2)$ & 1 & $-135/256$\\\bottomrule
\end{tabular}\end{center}
Tổng vẫn bằng $401/256$. Các hệ số bằng 0 phải được giữ để thấy bậc đó đã xét.

\subsection{Bốn mốc: Gauss vẫn hợp lệ nhưng hai gốc khác nhau}
Cho $x=[-1,0,1,2]$, $y=[0,1,2,9]$. Bảng có $\Delta=[1,1,7]$, $\Delta^2=[0,6]$, $\Delta^3=[6]$. Đa thức là $P=x^3+1$.

Gauss I lấy gốc 0, $t=x$, các hệ số $[1,1,0,1]$:
\[
 Q_{\rm I}(t)=1+t+0\,t(t-1)+t(t-1)(t+1)=1+t^3.
\]
Tại $z=1/2$, các đóng góp $1,1/2,0,-3/8$ cho $9/8$.

Gauss II lấy gốc 1, $t=x-1$, các hệ số $[2,1,3,1]$:
\[
 Q_{\rm II}(t)=2+t+3t(t+1)+t(t+1)(t-1).
\]
Tại $t=-1/2$, các đóng góp $2,-1/2,-3/4,3/8$ cũng cho $9/8$. Không dùng cùng một $t$ cho hai bảng khi gốc đã khác.

\section{Stirling: lẻ lấy trung bình, chẵn lấy một ô giữa}
\label{sec:stirling}
\textbf{Đầu vào}: $2N+1$ mốc đều, đánh chỉ số $-N,\ldots,N$, gốc $a=x_0$, $t=(x-a)/h$. Công thức đầy đủ:
\begin{align}
 Q(t)={}&y_0+\sum_{r=0}^{N-1}
 \frac{\Delta^{2r+1}y_{-r-1}+\Delta^{2r+1}y_{-r}}{2(2r+1)!}
 t\prod_{j=1}^r(t^2-j^2)\notag\\
 &+\sum_{r=1}^{N}\frac{\Delta^{2r}y_{-r}}{(2r)!}
 t^2\prod_{j=1}^{r-1}(t^2-j^2).
\end{align}
Tích rỗng là 1. Các miền như $|t|\le1/4$ là gợi ý dùng thuận lợi; không phải điều kiện làm đa thức tồn tại hay không.

\subsection{Quy trình từng bước}
\begin{enumerate}
\item Chọn đúng số mốc lẻ. Xác định mốc giữa, $h$ và $t_z$.
\item Lập hết bảng $\Delta$ đến cấp $2N$.
\item Với hạng lẻ $k=2r+1$, lấy hai ô đối xứng, tính trung bình $\sigma_k$; hạng chẵn $k=2r$ lấy một ô $\Delta^{2r}y_{-r}$.
\item Chia $\sigma_k$ cho $k!$ đúng một lần để có $\beta_k$; dựng tích $B_k$ tương ứng.
\item Ghép $Q=\sum\beta_kB_k$; thế $t_z$ vào từng hạng, cộng theo bảng. Nếu cần đạo hàm, lấy đạo hàm chính các hạng này.
\end{enumerate}

\subsection{Bài mẫu năm mốc, không bỏ một phép lấy hệ số}
Từ bảng ở mục \ref{sec:difference-tables}, ta có:
\begin{center}\small
\begin{tabular}{@{}rllrr@{}}\toprule
$k$ & Tính $\sigma_k$ từ bảng & $B_k(t)$ & $k!$ & $\beta_k$\\\midrule
0 & $y_0=1$ & 1 & 1 & 1\\
1 & $(0+4)/2=2$ & $t$ & 1 & 2\\
2 & $\Delta^2y_{-1}=4$ & $t^2$ & 2 & 2\\
3 & $(-12+12)/2=0$ & $t(t^2-1)$ & 6 & 0\\
4 & $\Delta^4y_{-2}=24$ & $t^2(t^2-1)$ & 24 & 1\\\bottomrule
\end{tabular}\end{center}
Do đó $Q=1+2t+2t^2+t^2(t^2-1)=1+2t+t^2+t^4$. Tại $t=1/4$:
\begin{center}\small
\begin{tabular}{@{}rrrrr@{}}\toprule
$k$ & $B_k(1/4)$ & đóng góp giá trị & tổng giá trị & đóng góp đạo hàm theo $t$\\\midrule
0 & 1 & 1 & 1 & 0\\
1 & $1/4$ & $1/2$ & $3/2$ & 2\\
2 & $1/16$ & $1/8$ & $13/8$ & 1\\
3 & $-15/64$ & 0 & $13/8$ & 0\\
4 & $-15/256$ & $-15/256$ & $401/256$ & $-7/16$\\\bottomrule
\end{tabular}\end{center}
Ở hạng cuối, $[t^2(t^2-1)]'=4t^3-2t$, tại $1/4$ bằng $-7/16$. Vì $h=1$:
\[
 P(1/4)=\frac{401}{256}=1.56640625,\qquad
 P'(1/4)=2+1-\frac7{16}=\frac{41}{16}=2.5625.
\]

\section{Bessel: chẵn lấy trung bình, lẻ lấy một ô}
\label{sec:bessel}
\textbf{Đầu vào}: $2N+2$ mốc đều, chỉ số $-N,\ldots,N+1$; gốc $a=x_0$ là mốc \emph{bên trái của cặp giữa}. Đặt
\[
 E_0=1,\qquad E_r(t)=\prod_{j=0}^{r-1}(t+j)(t-j-1).
\]
\begin{align}
 Q(t)=\sum_{r=0}^{N}\bigg[
 &\frac{\Delta^{2r}y_{-r}+\Delta^{2r}y_{-r+1}}{2(2r)!}E_r(t)\notag\\
 &+\frac{\Delta^{2r+1}y_{-r}}{(2r+1)!}(t-\tfrac12)E_r(t)\bigg].
\end{align}
\begin{enumerate}
\item Lập bảng $\Delta$ và xác định hai mốc giữa. Không chọn mốc phải làm gốc rồi vẫn dùng công thức trên.
\item Với mỗi $r$, tính hạng chẵn từ trung bình hai ô; chia $(2r)!$ và nhân $E_r$.
\item Tính hạng lẻ từ một ô; chia $(2r+1)!$ và nhân $(t-1/2)E_r$.
\item Chuẩn bị cặp hạng sau bằng $E_{r+1}=E_r(t+r)(t-r-1)$.
\item Thế số và cộng. Với đạo hàm, vẫn phải đạo hàm các hạng lẻ ngay cả khi $t_z=1/2$.
\end{enumerate}

Ví dụ bốn mốc $[-1,0,1,2]$, $y=[0,1,2,9]$, gốc 0, $h=1$:
\begin{center}\small
\begin{tabular}{@{}rllrr@{}}\toprule
$k$ & $\sigma_k$ & $B_k(t)$ & $\beta_k$ & đóng góp tại $t=1/2$\\\midrule
0 & $(1+2)/2=3/2$ & 1 & $3/2$ & $3/2$\\
1 & $\Delta y_0=1$ & $t-1/2$ & 1 & 0\\
2 & $(0+6)/2=3$ & $t(t-1)$ & $3/2$ & $-3/8$\\
3 & $\Delta^3y_{-1}=6$ & $(t-1/2)t(t-1)$ & 1 & 0\\\bottomrule
\end{tabular}\end{center}
Ghép rồi rút gọn:
\[
 Q(t)=\frac32+(t-\tfrac12)+\frac32t(t-1)+(t-\tfrac12)t(t-1)=1+t^3.
\]
Kết quả $P(1/2)=9/8$. Đóng góp đạo hàm lần lượt là $0,1,0,-1/4$, nên $P'(1/2)=3/4$. Hạng cuối có giá trị 0 nhưng đạo hàm $-1/4$, không được bỏ nó trước khi lấy đạo hàm. Với hai mốc, Bessel trở về nội suy tuyến tính qua hai mốc đó.

\section{Lấy đạo hàm từ chính đa thức đã dựng}
\label{sec:interp-derivative}
Nếu $Q(t)=\sum_k\beta_kB_k(t)$ và $x=a+ht$, thì
\[
 P^{(s)}(z)=\frac{Q^{(s)}(t_z)}{h^s}
          =\sum_k\frac{\beta_kB_k^{(s)}(t_z)}{h^s},\qquad t_z=\frac{z-a}{h}.
\]
\textbf{Bốn việc riêng biệt}: đổi điểm $z$ sang $t_z$; đạo hàm $Q$ theo $t$; thế $t_z$; chia cho $h^s$. Ví dụ $x=4+2t$ và $Q=16+16t+4t^2$: tại $x=4$, $t=0$, $Q'(0)=16$, $P'(4)=16/2=8$; $Q''(0)=8$, $P''(4)=8/2^2=2$.

Với tích $B(t)=\prod_j(t-r_j)$, công thức $B'=B\sum1/(t-r_j)$ có dạng không xác định khi $t=r_j$. Tool dùng truy hồi đạo hàm của tích, không chia cho các nhân tử:
\begin{enumerate}
\item Lưu $g_0=1$, $g_1=\cdots=g_s=0$.
\item Khi nhân thêm $(t-r)$, cập nhật từ $l=s$ xuống 1: $g_l\leftarrow(t_z-r)g_l+l\,g_{l-1}$, sau đó $g_0\leftarrow(t_z-r)g_0$.
\item Sau nhân tử cuối, $g_l=B^{(l)}(t_z)$; nhân với $\beta_k$ rồi chia $h^s$ nếu dùng biến chuẩn hóa.
\end{enumerate}
Các giá trị cũ phải được dùng trong cùng lần cập nhật; vì vậy đi từ cấp cao xuống thấp. Nếu $s$ lớn hơn bậc đa thức, đạo hàm của mô hình bằng 0; không suy ra đạo hàm của hàm thật cũng bằng 0.

\section{Horner: hạ hệ số đầu, nhân sang phải, cộng xuống}
\label{sec:horner}
\textbf{Đầu vào}: hệ số của đa thức theo bậc \emph{giảm}, điểm $c$, cấp đạo hàm lớn nhất. Ví dụ
\[
 P(x)=3x^3-2x-1=3x^3+0x^2-2x-1
\]
có dãy \texttt{3 0 -2 -1}. Nếu bỏ 0 của $x^2$ thì đã nhập một đa thức khác.

\subsection{Một lượt tính $P(c)$ và thương của phép chia}
Với $P(x)=a_0x^n+\cdots+a_n$:
\[
 b_0=a_0,\qquad b_j=a_j+c\,b_{j-1}.
\]
\textbf{Đọc bằng số}: với $c=2$, làm bảng ngang ba hàng sau. Vạch đứng tách ô dư cuối khỏi các hệ số thương ở hàng cuối.
\begin{center}\small
\begin{tabular}{@{}lrrr|r@{}}\toprule
$c=2$ & Cột 1 & Cột 2 & Cột 3 & Cột 4: dư\\\midrule
Hệ số đầu vào & 3 & 0 & $-2$ & $-1$\\
Nhân $2$, chuyển phải & --- & 6 & 12 & 20\\\midrule
Cộng xuống & \textbf{3} & \textbf{6} & \textbf{10} & \textbf{19}\\\bottomrule
\end{tabular}\end{center}
\begin{enumerate}
\item Hạ 3 xuống hàng Cộng; ô đầu hàng Nhân để trống.
\item Lấy 3 vừa hạ nhân 2: $2\cdot3=6$; đặt sang cột 2 rồi cộng $0+6=6$.
\item Lấy 6 vừa cộng nhân 2: $2\cdot6=12$; đặt sang cột 3 rồi cộng $-2+12=10$.
\item Lấy 10 vừa cộng nhân 2: $2\cdot10=20$; đặt sang cột 4 rồi cộng $-1+20=19$.
\end{enumerate}
Ô cuối là số dư $P(2)=19$. Các ô trước nó, $3,6,10$, là hệ số của thương $Q_1=3x^2+6x+10$:
\[
 P(x)=(x-2)Q_1(x)+19.
\]
Tool hiện có thể trình bày cùng vết tính bằng bảng dọc nhân--cộng; sơ đồ ngang ở sổ tay biểu diễn đúng các phép tính đó. Giao diện Horner chính nhận đa thức lũy thừa; API chia riêng \texttt{horner\_division} nhận hệ số \emph{bậc tăng}, khác thứ tự nhập GUI.

\subsection{Lấy đạo hàm bằng các lượt chia tiếp}
Chia liên tiếp cho cùng $(x-c)$:
\[
 Q_r(x)=(x-c)Q_{r+1}(x)+s_r,\qquad Q_0=P.
\]
Sau mỗi lượt, ghi riêng số dư; đưa \emph{chỉ các hệ số thương} lên hàng đầu của lượt sau. Không đưa số dư vào dãy hệ số mới.

\begin{center}\small
\begin{tabular}{@{}lrr|r@{}}\toprule
Lượt 2, vẫn $c=2$ & 1 & 2 & 3: dư\\\midrule
Hệ số thương trước & 3 & 6 & 10\\
Nhân 2 & --- & 6 & 24\\\midrule
Cộng & \textbf{3} & \textbf{12} & \textbf{34}\\\bottomrule
\end{tabular}\qquad
\begin{tabular}{@{}lr|r@{}}\toprule
Lượt 3 & 1 & 2: dư\\\midrule
Hệ số & 3 & 12\\
Nhân 2 & --- & 6\\\midrule
Cộng & \textbf{3} & \textbf{18}\\\bottomrule
\end{tabular}\end{center}
Lượt 2: $6+2\cdot3=12$, $10+2\cdot12=34$. Lượt 3: $12+2\cdot3=18$. Lượt 4 chỉ còn hằng số 3 nên số dư là 3.

\begin{center}\begin{tabular}{@{}rrrrr@{}}\toprule
Lượt & $r$ & Số dư $s_r$ & Nhân $r!$ & $P^{(r)}(2)$\\\midrule
1 & 0 & 19 & 1 & 19\\
2 & 1 & 34 & 1 & 34\\
3 & 2 & 18 & 2 & 36\\
4 & 3 & 3 & 6 & 18\\\bottomrule
\end{tabular}\end{center}
Vì thế $P''(2)=36$, không phải 18. Lý do: thay các đẳng thức chia vào nhau cho
\[
 P(x)=s_0+s_1(x-c)+s_2(x-c)^2+\cdots+s_n(x-c)^n,
 \qquad P^{(r)}(c)=r!s_r.
\]
Cấp $r\le n$ cần $r+1$ lượt, tính cả lượt đầu của $P(c)$. Các cấp lớn hơn bậc thực bằng 0, không cần dựng hàng nghìn lượt 0. Các số 0 đầu dãy không làm tăng bậc thực; số 0 giữa/cuối phải giữ.

\subsection{Mẫu đề 20241: nhìn được từng dãy chuyển tiếp}
Đa thức được in trong ảnh đề, không lấy từ workbook:
\[
P(x)=3.5x^5-2.78x^4+5.63x^3+0x^2-15.28x+1.54,\qquad c=2.57.
\]
Muốn cấp 2 và 3, thực hiện bốn lượt. Các dãy của đúng phép tính trong tool là:
\begin{center}\small
\begin{tabular}{@{}rl@{}}\toprule
Lượt & Hệ số đầu vào $\longrightarrow$ toàn bộ hàng tích lũy\\\midrule
1 & $(3.5,-2.78,5.63,0,-15.28,1.54)$\\
 & $\longrightarrow(3.5,6.215,21.60255,55.5185535,127.402682495,328.96489401215)$\\
2 & $(3.5,6.215,21.60255,55.5185535,127.402682495)$\\
 & $\longrightarrow(3.5,15.21,60.69225,211.497636,670.951607015)$\\
3 & $(3.5,15.21,60.69225,211.497636)$\\
 & $\longrightarrow(3.5,24.205,122.8991,527.348323)$\\
4 & $(3.5,24.205,122.8991)$\\
 & $\longrightarrow(3.5,33.2,208.2231)$\\\bottomrule
\end{tabular}\end{center}
Chẳng hạn lượt 3: hạ $3.5$; $15.21+2.57\cdot3.5=24.205$; $60.69225+2.57\cdot24.205=122.8991$; $211.497636+2.57\cdot122.8991=527.348323$.
\[
 \boxed{P''(2.57)=2!\cdot527.348323=1054.696646},\qquad
 \boxed{P'''(2.57)=3!\cdot208.2231=1249.3386}.
\]
Bảng tóm tắt này chỉ gom các lượt của cùng thuật toán; khi cần mọi phép nhân--cộng, mở bài JSON Horner đề 20241 để tool xuất đầy đủ.

\section{Sai số, kiểm tra và các tình huống dễ hiểu sai}
\label{sec:interp-errors}
Nếu dữ liệu chính xác $y_i=f(x_i)$ và $f\in C^m$ trên đoạn chứa toàn bộ mốc cùng điểm $z$:
\[
 f(z)-P(z)=\frac{f^{(m)}(\xi_z)}{m!}\prod_{i=0}^{m-1}(z-x_i),\qquad
 |f(z)-P(z)|\le\frac{M_m}{m!}\left|\prod_i(z-x_i)\right|.
\]
Phải có căn cứ cho $M_m$. Không có cận đạo hàm thì chưa có cận sai số số. Hạng đầu bỏ đi có thể là ước lượng khi có giả thiết phù hợp; không mặc định là cận bảo đảm.

Nếu $|\delta y_i|\le\eta_i$, phần sai lệch do dữ liệu bị chặn bởi
\[
 |\delta P(z)|\le\sum_i|L_i(z)|\eta_i.
\]
Bậc cao có thể khuếch đại nhiễu hoặc dao động mạnh giữa các mốc. Tính nhiều chữ số trung gian không loại bỏ điều này. Sai lệch $P(x_i)-y_i$ nhỏ kiểm tra việc dựng đa thức; không là sai số của hàm thật giữa mốc. Không đạo hàm phần dư có $\xi_z$ như thể $\xi_z$ là hằng số.

\begin{center}\small
\begin{tabular}{@{}p{.30\linewidth}p{.60\linewidth}@{}}\toprule
Tình huống & Cách xử lý đúng\\\midrule
Hoành độ trùng & Dừng kiểm tra dữ liệu; không tự gộp hay lấy trung bình\\
Không đều nhưng chọn Newton đều/Gauss/Stirling/Bessel & Đổi lựa chọn theo đề hoặc sửa dữ liệu nguồn; không tự nắn các $x$\\
Stirling số mốc chẵn/Bessel số mốc lẻ & Chọn lại tập mốc công khai; không tự bỏ một mốc\\
Điểm trùng mốc & Đánh giá bằng tích/lồng hợp lệ; không tạo $0/0$\\
Ngoài dải mốc & Ghi ngoại suy; cận đạo hàm phải đúng cả đoạn mở rộng\\
Mốc có gốc lớn, bước rất nhỏ & Kiểm tra hiệu mốc; đổi gốc/đơn vị khi cần\\
Khai triển theo $x$ có hệ số lớn & Giữ dạng tích hoặc biến chuẩn hóa; đối chiếu Horner Newton ổn định\\
Đạo hàm vượt bậc $P$ & Bằng 0 cho mô hình; không kết luận $f^{(s)}=0$\\
\bottomrule
\end{tabular}\end{center}

\paragraph{Số học trong tool.} Bảng lớn được giữ dạng tích; phép đánh giá nội suy dùng Newton lồng với độ chính xác trung gian tăng theo số mốc/thang dữ liệu. Lời giải theo công thức đã chọn được đối chiếu với cách đánh giá này. Các lưới nhận là gần đều trong dung sai có thể làm hai cách tính lệch; đọc phần đối chiếu và cảnh báo trước khi chép số. Đây là kiểm tra số học, không là chứng nhận sai số của hàm thật.

\paragraph{Nguồn đối chiếu trong repo.} Các đường tính chính: \texttt{pps/interpolation.py}, \texttt{pps/polynomial\_model.py}, \texttt{pps/polynomial\_worked.py}. Hỗ trợ: \texttt{chapter1\_tools.horner\_division}, \texttt{newton\_horner}, \texttt{newton\_append}, các hàm cận sai số và điền dữ liệu theo giả thiết bậc. Tài liệu nền: phần Lagrange, Newton, trung tâm và Horner của PDF Chương 1 trong \texttt{kienthuc/real}. Ví dụ số nhỏ ở chương này là tự dựng; đa thức Horner 20241 được ghi đúng từ ảnh đề trong repo.

\chapter{Nội suy ngược: chọn đúng mô hình và theo dõi từng vòng}
\label{ch:inverse}
\section{Input, quá trình và output của sáu lựa chọn nội suy ngược}
\subsection*{Input: các ô đúng tên trên giao diện}
Bảng này được lấy từ biểu mẫu của tool. Giá trị mặc định là ví dụ nhập, cần thay bằng dữ liệu đề. Hai dãy $x,y$ là bảng dữ liệu riêng bên dưới nếu chức năng dùng bảng.
\subsubsection*{Ô chung của các lựa chọn trong nhóm}
\begin{center}\small\begin{tabular}{@{}p{.32\linewidth}p{.17\linewidth}p{.43\linewidth}@{}}\toprule
Ô nhập & Khóa JSON & Cách nhập / mặc định\\\midrule
Mốc dùng (JSON,\allowbreak{} trống: toàn bộ bảng) & \texttt{selected\_indices} & Để trống
\par Ví dụ [9,\allowbreak{}10,\allowbreak{}11,\allowbreak{}12] dùng đúng bốn dòng đó,\allowbreak{} đánh số từ 0;\allowbreak{} dữ liệu nguồn vẫn được giữ.\\[3pt]
Mức trình bày & \texttt{detail\_level} & \texttt{Tổng hợp}
\par Bảng tổng hợp luôn hiện;\allowbreak{} có thể mở từng nhóm phép tính hoặc từng vòng. Bản xuất vẫn giữ đủ chi tiết.\\[3pt]
\bottomrule\end{tabular}\end{center}
\subsubsection*{Ô riêng: Lặp tiến}
\begin{center}\small\begin{tabular}{@{}p{.32\linewidth}p{.17\linewidth}p{.43\linewidth}@{}}\toprule
Ô nhập & Khóa JSON & Cách nhập / mặc định\\\midrule
Giá trị y cần tìm ngược & \texttt{target} & \texttt{4}
\par Tập cặp nhập là các mốc đã chọn;\allowbreak{} bậc bằng số cặp trừ 1. Không tự đổi nhánh hoặc mốc. Lặp thực hành mặc định kiểm tra thêm phần dư 1e-6,\allowbreak{} tối đa 100 vòng.\\[3pt]
Cận x trái (trống: khoảng đổi dấu duy nhất) & \texttt{a} & Để trống\\[3pt]
Cận x phải (trống: khoảng đổi dấu duy nhất) & \texttt{b} & Để trống\\[3pt]
Dung sai vị trí/\allowbreak{}bước theo x & \texttt{tol} & \texttt{1e-06}
\par Giữ toàn bộ tập mốc và gốc đầu/\allowbreak{}cuối tập;\allowbreak{} khoảng nghiệm có thể ở giữa bảng.\\[3pt]
Dung sai phần dư & \texttt{residual\_tol} & \texttt{1e-06}\\[3pt]
Số vòng tối đa & \texttt{max\_iter} & \texttt{100}\\[3pt]
Quy tắc khởi tạo & \texttt{initial\_rule} & \texttt{Tuyến tính tại gốc}\\[3pt]
t₀ khi khởi tạo thủ công & \texttt{initial\_t} & Để trống\\[3pt]
\bottomrule\end{tabular}\end{center}
\subsubsection*{Ô riêng: Lặp lùi}
\begin{center}\small\begin{tabular}{@{}p{.32\linewidth}p{.17\linewidth}p{.43\linewidth}@{}}\toprule
Ô nhập & Khóa JSON & Cách nhập / mặc định\\\midrule
Giá trị y cần tìm ngược & \texttt{target} & \texttt{2}
\par Tập cặp nhập là các mốc đã chọn;\allowbreak{} bậc bằng số cặp trừ 1. Không tự đổi nhánh hoặc mốc. Lặp thực hành mặc định kiểm tra thêm phần dư 1e-6,\allowbreak{} tối đa 100 vòng.\\[3pt]
Cận x trái (trống: khoảng đổi dấu duy nhất) & \texttt{a} & Để trống\\[3pt]
Cận x phải (trống: khoảng đổi dấu duy nhất) & \texttt{b} & Để trống\\[3pt]
Dung sai vị trí/\allowbreak{}bước theo x & \texttt{tol} & \texttt{1e-06}
\par Giữ toàn bộ tập mốc và gốc đầu/\allowbreak{}cuối tập;\allowbreak{} khoảng nghiệm có thể ở giữa bảng.\\[3pt]
Dung sai phần dư & \texttt{residual\_tol} & \texttt{1e-06}\\[3pt]
Số vòng tối đa & \texttt{max\_iter} & \texttt{100}\\[3pt]
Quy tắc khởi tạo & \texttt{initial\_rule} & \texttt{Tuyến tính tại gốc}\\[3pt]
t₀ khi khởi tạo thủ công & \texttt{initial\_t} & Để trống\\[3pt]
\bottomrule\end{tabular}\end{center}
\subsubsection*{Ô riêng: Lặp có cận co}
\begin{center}\small\begin{tabular}{@{}p{.32\linewidth}p{.17\linewidth}p{.43\linewidth}@{}}\toprule
Ô nhập & Khóa JSON & Cách nhập / mặc định\\\midrule
Giá trị y cần tìm ngược & \texttt{target} & \texttt{4}
\par Tập cặp nhập là các mốc đã chọn;\allowbreak{} bậc bằng số cặp trừ 1. Không tự đổi nhánh hoặc mốc. Lặp thực hành mặc định kiểm tra thêm phần dư 1e-6,\allowbreak{} tối đa 100 vòng.\\[3pt]
Cận x trái của khoảng nghiệm & \texttt{a} & \texttt{1.6}\\[3pt]
Cận x phải của khoảng nghiệm & \texttt{b} & \texttt{1.8}\\[3pt]
Dung sai vị trí & \texttt{tol} & \texttt{1e-06}\\[3pt]
Dung sai phần dư & \texttt{residual\_tol} & \texttt{1e-06}\\[3pt]
Hướng nội suy & \texttt{direction} & \texttt{Tiến}\\[3pt]
Số vòng tối đa & \texttt{max\_iter} & \texttt{100}\\[3pt]
Quy tắc khởi tạo & \texttt{initial\_rule} & \texttt{Tuyến tính tại gốc}\\[3pt]
t₀ khi khởi tạo thủ công & \texttt{initial\_t} & Để trống\\[3pt]
\bottomrule\end{tabular}\end{center}
\subsubsection*{Ô riêng: Hàm ngược — Newton /\allowbreak{} Hàm ngược — Lagrange}
\begin{center}\small\begin{tabular}{@{}p{.32\linewidth}p{.17\linewidth}p{.43\linewidth}@{}}\toprule
Ô nhập & Khóa JSON & Cách nhập / mặc định\\\midrule
Giá trị y cần tìm ngược & \texttt{target} & \texttt{6}
\par Tập cặp nhập là các mốc đã chọn;\allowbreak{} bậc bằng số cặp trừ 1. Không tự đổi nhánh hoặc mốc. Lặp thực hành mặc định kiểm tra thêm phần dư 1e-6,\allowbreak{} tối đa 100 vòng.\\[3pt]
Số mốc đã chọn (để trống: toàn tập nhập) & \texttt{count} & Để trống
\par Không tự lấy mốc gần mức. Muốn dùng ít mốc hơn,\allowbreak{} giữ đúng các cặp cần dùng trong bảng nhập.\\[3pt]
\bottomrule\end{tabular}\end{center}
\subsubsection*{Ô riêng: Chia đôi đa thức}
\begin{center}\small\begin{tabular}{@{}p{.32\linewidth}p{.17\linewidth}p{.43\linewidth}@{}}\toprule
Ô nhập & Khóa JSON & Cách nhập / mặc định\\\midrule
Giá trị y cần tìm ngược & \texttt{target} & \texttt{2}
\par Tập cặp nhập là các mốc đã chọn;\allowbreak{} bậc bằng số cặp trừ 1. Không tự đổi nhánh hoặc mốc. Lặp thực hành mặc định kiểm tra thêm phần dư 1e-6,\allowbreak{} tối đa 100 vòng.\\[3pt]
Cận trái & \texttt{a} & \texttt{1}\\[3pt]
Cận phải & \texttt{b} & \texttt{2}\\[3pt]
Dung sai vị trí & \texttt{tol} & \texttt{0.001}\\[3pt]
\bottomrule\end{tabular}\end{center}

\paragraph{Input chung.} $x,y$ là bảng nguồn; $\bar y$ là mức cần tìm $x$. Ô \emph{Mốc dùng} chọn chỉ số từ 0, trống dùng toàn bảng. Khoảng $[a,b]$ là nơi muốn tìm nghiệm, không thay gốc của đa thức. Không lấy số mốc làm số vòng lặp.

\subsection{Bước đầu chung: bảng khảo sát tổng quát}
\begin{center}\small\begin{tabular}{cllll}\toprule
Dòng&$x_i$&$F_i=y_i-\bar y$&So với dòng sau&Thông tin chuyển tiếp\\\midrule
0&$x_0$&$y_0-\bar y$&Dấu $F_0,F_1$&Mốc đúng mức/khoảng đổi dấu\\
$i$&$x_i$&$y_i-\bar y$&Dấu $F_i,F_{i+1}$&Ứng viên $[x_i,x_{i+1}]$\\
$n$&$x_n$&$y_n-\bar y$&Không có dòng sau&Kiểm riêng $F_n=0$\\\bottomrule
\end{tabular}\end{center}
Sau khảo sát mới chọn nhánh, tập mốc và thuật toán. Nhiều khoảng thì phải giải từng nhánh cần tìm. Bảng chỉ phát hiện đổi dấu giữa các mẫu; có thể bỏ sót nghiệm tiếp xúc.

\subsection{Hàm ngược -- Newton và Hàm ngược -- Lagrange}
\textbf{Input}: tập mốc thuộc nhánh tăng/giảm nghiêm ngặt, tung độ phân biệt; mức trong dải tung độ nếu không bật ngoại suy. \textbf{Quá trình}: đảo cặp $(u_i,v_i)=(y_i,x_i)$ $\to$ dựng $Q(y)$ $\to$ thế $\bar y$.
\begin{center}\small\begin{tabular}{cllll}\toprule
Dòng&$u_i$&$v_i$&Tỉ sai phân cấp 1&Cấp 2\\\midrule
0&$y_0$&$x_0$&$(x_1-x_0)/(y_1-y_0)$&$D_{0,2}$\\
$i$&$y_i$&$x_i$&$(x_{i+1}-x_i)/(y_{i+1}-y_i)$&$D_{i,2}$ nếu đủ mốc\\
$n$&$y_n$&$x_n$&&\\\bottomrule
\end{tabular}\end{center}
Newton lấy $D_{0,k}$ rồi cộng
\[
Q(\bar y)=x_0+D_{0,1}(\bar y-y_0)+D_{0,2}(\bar y-y_0)(\bar y-y_1)+\cdots.
\]
Lagrange tính $\ell_i(\bar y)=\prod_{j\ne i}(\bar y-y_j)/(y_i-y_j)$ rồi cộng
\[
\widehat x=x_0\ell_0(\bar y)+x_1\ell_1(\bar y)+\cdots+x_n\ell_n(\bar y).
\]
\textbf{Output}: đa thức hàm ngược $Q$, giá trị $Q(\bar y)$ và tập mốc; không có vòng hội tụ. Tung độ thường không đều dù hoành độ gốc đều.

\subsection{Lặp tiến, Lặp lùi, Lặp có cận co}
\textbf{Input}: lưới đều, tập mốc cố định, mức, khoảng, ngưỡng vị trí $\varepsilon_x$, ngưỡng phần dư $\varepsilon_y$, khởi tạo và số vòng tối đa. Lặp có cận co thêm lựa chọn hướng; hai mục tiến/lùi mặc định dùng tiêu chuẩn thực hành.

\textbf{Bước 1}: dựng đa thức thuận $P_t(t)=c_0+dt+H(t)$ một lần:
\begin{align*}
\text{Tiến: }&t=(x-x_0)/h,\ c_0=y_0,\ d=\Delta y_0,\\
&H(t)=\tfrac{\Delta^2y_0}2t(t-1)+\tfrac{\Delta^3y_0}6t(t-1)(t-2)+\cdots;\\
\text{Lùi: }&t=(x-x_n)/h,\ c_0=y_n,\ d=\Delta y_{n-1},\\
&H(t)=\tfrac{\Delta^2y_{n-2}}2t(t+1)+\tfrac{\Delta^3y_{n-3}}6t(t+1)(t+2)+\cdots.
\end{align*}
\textbf{Bước 2}: lập $\varphi(t)=(\bar y-c_0-H(t))/d$, chọn $t_0$ theo ô khởi tạo. $d$ phải khác 0. Nếu có chứng nhận, kiểm $t_0\in J$, $\varphi(J)\subseteq J$, $q\ge\sup_J|\varphi'|$, $q<1$.

\textbf{Bước 3}: điền bảng vòng lặp tổng quát theo hai nửa để đọc rõ:
\begin{center}\small\begin{tabular}{cllll}\toprule
Vòng&Giá trị nhận&Phần bậc cao&Tử số&Giá trị mới\\\midrule
1&$t_0$&$H(t_0)$&$\bar y-c_0-H(t_0)$&$t_1=(\bar y-c_0-H(t_0))/d$\\
$k+1$&$t_k$&$H(t_k)$&$\bar y-c_0-H(t_k)$&$t_{k+1}=(\bar y-c_0-H(t_k))/d$\\\bottomrule
\end{tabular}\end{center}
\begin{center}\small\begin{tabular}{cllll}\toprule
Vòng&Đổi lại hoành độ&Bước hoặc cận&Phần dư tại điểm mới&Quyết định\\\midrule
$k+1$&$x_{k+1}=a_0+ht_{k+1}$&$D_{k+1}$ hoặc $B_{k+1}$&$|P_t(t_{k+1})-\bar y|$&Đạt cả hai ngưỡng?\\\bottomrule
\end{tabular}\end{center}
Ở đây $a_0$ là gốc đầu/cuối tập mốc, không phải cận trái khoảng nghiệm.
$D_{k+1}=|h||t_{k+1}-t_k|$ dùng ở thực hành;
$B_{k+1}=|h|q|t_{k+1}-t_k|/(1-q)$ dùng khi có chứng nhận hợp lệ.
Vòng chưa đạt chuyển \emph{giá trị mới} thành đầu vào vòng sau. Không dùng $H(t_k)$ để tính phần dư tại $t_{k+1}$.

\textbf{Output}: nghiệm mô hình trong khoảng yêu cầu, trạng thái dừng, gốc/hướng/tập mốc, bảng mọi vòng, bước/cận và phần dư. Hết vòng hoặc tràn số là trạng thái chưa đạt. Đúng một mốc hoặc đa thức thực tế tuyến tính có thể cho kết quả không cần vòng lặp.

\subsection{Chia đôi đa thức}
\textbf{Input}: cùng bảng và mức nhưng cần khoảng $[a,b]$, ngưỡng vị trí. \textbf{Quá trình}: dựng $P$ $\to$ đặt $F=P-\bar y$ $\to$ kiểm hai đầu $\to$ thu khoảng:
\begin{center}\small\begin{tabular}{cllll}\toprule
Lần&Khoảng nhận&Trung điểm&Giá trị kiểm&Khoảng chuyển tiếp\\\midrule
$k+1$&$[a_k,b_k]$&$c_k=(a_k+b_k)/2$&$F(c_k)$&Giữ nửa có đổi dấu\\\bottomrule
\end{tabular}\end{center}
Khi nửa độ rộng đạt ngưỡng, \textbf{output} là trung điểm của khoảng cuối, khoảng cuối và cận nửa độ rộng. Trung điểm vừa thử ở hàng cuối không nhất thiết là trung điểm của khoảng sau thu. Phần giải thích và ví dụ đầy đủ theo từng mô hình ở các mục tiếp theo.


\section{Cùng câu hỏi tìm x, nhưng có hai mô hình khác nhau}
Cho bảng $(x_i,y_i)$ và một mức $\bar y$. Bài toán hỏi hoành độ $x$ sao cho $f(x)=\bar y$. Trước khi tính, phải quyết định đang làm việc gì:
\begin{center}\small
\begin{tabular}{@{}p{.24\linewidth}p{.34\linewidth}p{.31\linewidth}@{}}\toprule
Mô hình & Phép tính thực hiện & Kết quả mang nghĩa gì?\\\midrule
Nội suy hàm ngược & Đảo cặp thành $(y_i,x_i)$, dựng $Q(y)$ rồi tính $Q(\bar y)$ & Giá trị của đa thức xấp xỉ $g=f^{-1}$ trên nhánh đã chọn\\
Giải đa thức thuận & Dựng $P(x)$ từ $(x_i,y_i)$ rồi giải $P(x)=\bar y$ bằng lặp hoặc chia đôi & Nghiệm gần đúng của mô hình thuận $P$\\\bottomrule
\end{tabular}\end{center}

Thông thường $Q$ \textbf{không phải nghịch đảo chính xác} của $P$. Vì vậy $Q(\bar y)$ và nghiệm của $P(x)=\bar y$ có thể khác nhau; đây không tự là lỗi của một trong hai phương pháp.

\subsection{Ví dụ mở đầu: vì sao hai đáp số khác nhau?}
Cho $x=[0,1,2]$, $y=[0,1,4]$, tìm tại mức $\bar y=2$.
\begin{itemize}
\item Nội suy thuận qua ba cặp cho $P(x)=x^2$. Giải $P(x)=2$ trên $[1,2]$ cho $x=\sqrt2\approx1.41421356$.
\item Đảo bảng thành $(0,0),(1,1),(4,2)$, nội suy bậc hai theo $y$ cho
\[
 Q(y)=y-\frac{y(y-1)}6,\qquad Q(2)=\frac53\approx1.66666667.
\]
\end{itemize}
Nếu thế $5/3$ vào mô hình thuận, $P(5/3)-2=7/9$, khác 0. Điều đó chỉ đối chiếu hai mô hình khác nhau; không phải tiêu chuẩn dừng của nội suy hàm ngược.

\paragraph{Dấu hiệu chọn phương pháp.} Đề nói ``nội suy hàm ngược'', ``đảo bảng'', ``Lagrange/Newton theo $y$'' thì dùng $Q$. Đề yêu cầu ``lập công thức lặp tiến/lùi'', ``kiểm tra co'', hoặc ``chia đôi đa thức'' thì dựng $P$ rồi giải phương trình mô hình. Nếu đề không chỉ định, phải công bố lựa chọn, tập mốc và nhánh.

\section{Đầu vào và những lựa chọn phải tách riêng}
\label{sec:inverse-input}
\begin{enumerate}
\item \textbf{Bảng nguồn}: hai dãy hữu hạn, cùng độ dài, ít nhất hai cặp. Tool giữ ID dòng gốc, sắp theo $x$ và không gộp hoành độ trùng.
\item \textbf{Tập mốc tính}: tập thực sự dùng dựng một đa thức. Nếu dùng $m$ mốc thì bậc tối đa $m-1$. Mọi vòng sau dùng lại đa thức này; không cần mốc mới cho mỗi vòng.
\item \textbf{Nhánh mục tiêu}: phần của bảng muốn xét tính một-một. Có thể có nhiều nhánh đi qua cùng mức $\bar y$; không tự lấy nhánh đầu tiên.
\item \textbf{Khoảng nghiệm $[a,b]$}: đoạn hoành độ đang yêu cầu tìm nghiệm, nằm trong dải mốc được chọn. Nó độc lập với hai đầu tập mốc tính.
\item \textbf{Gốc khai triển}: tiến lấy mốc đầu của \emph{tập tính}; lùi lấy mốc cuối. Không dời gốc sang sát khoảng nghiệm rồi đòi thêm mốc.
\item \textbf{Khởi tạo}: tuyến tính tại gốc, tuyến tính giữa hai đầu khoảng, hoặc nhập thủ công. Phải ghi rõ cách đã chọn.
\item \textbf{Ngưỡng và chế độ}: $\varepsilon_x$, $\varepsilon_y$, số vòng tối đa; thực hành hay yêu cầu chứng nhận co.
\end{enumerate}

\subsection{Cách nhập trong giao diện}
\begin{center}\small
\begin{tabular}{@{}p{.30\linewidth}p{.59\linewidth}@{}}\toprule
Ô/lựa chọn & Nghĩa chính xác\\\midrule
Giá trị $y$ cần tìm ngược & Mức $\bar y$, không phải điểm $x$ cần tính của nội suy thuận\\
Mốc dùng, JSON & Ví dụ \texttt{[0,1,2]} là ID dòng nguồn từ 0; trống là toàn bảng\\
Số mốc ở Hàm ngược & Kiểm tra tập đưa vào đã có đúng số mốc; không tự lấy các mốc gần mức\\
Cận $a,b$ & Khoảng hoành độ độc lập với gốc tiến/lùi; phải nhập đủ cả hai nếu đã nhập một đầu\\
Dung sai vị trí & $\varepsilon_x$, đơn vị của hoành độ\\
Dung sai phần dư & $\varepsilon_y$, đơn vị của tung độ\\
Khởi tạo tại gốc & Bỏ phần bậc cao để tính $t_0=(\bar y-y_{\rm base})/d$\\
Khởi tạo hai đầu khoảng & Nội suy tuyến tính giữa hai đầu của khoảng mục tiêu rồi đổi sang $t$\\
Thủ công & Nhập $t_0$ theo gốc và $h$ của đa thức đang dùng\\
Lặp có cận co & Chọn hướng tiến/lùi; chỉ báo đạt khi đủ chứng nhận trên $J$\\\bottomrule
\end{tabular}\end{center}

Nếu không nhập khoảng cho lặp thực hành, tool chỉ tự xác định khi có một khoảng đổi dấu duy nhất trên tập đã nhập (hoặc trường hợp mốc đúng mức duy nhất thích hợp). Có nhiều ứng viên thì yêu cầu chọn. Chia đôi luôn cần khoảng được chỉ định. Các cấu hình như đoạn chứng nhận $J$, cận sai số hàm thật hoặc nguồn chứng nhận có thể được cung cấp qua bài JSON/API; không giả định mọi trường kỹ thuật đều là ô riêng đang hiện trên màn hình.

\section{Khảo sát nhánh: bắt đầu từ bảng trước khi lập công thức}
\label{sec:inverse-scan}
Tính $F_i=y_i-\bar y$ và dấu các hiệu $y_{i+1}-y_i$.
\begin{enumerate}
\item Nếu $F_i=0$, đánh dấu mốc đúng mức. ``Gần bằng'' trong ngưỡng cảnh báo không tự biến thành bằng đúng.
\item Nếu $F_iF_{i+1}<0$, đánh dấu một khoảng đổi dấu. Khi hàm thật liên tục, khoảng đó có ít nhất một nghiệm.
\item Theo dấu sai phân cấp một để chia các nhánh mẫu tăng/giảm. Một đoạn bằng phẳng phải được xét riêng.
\item Chọn nhánh/khoảng phù hợp với câu hỏi; giữ nguyên lựa chọn trong toàn bộ lượt chạy.
\end{enumerate}
Với $x=[0,1,2,3,4]$, $y=[0,1,2,1,0]$, mức $1.5$, bảng quét là
\begin{center}\begin{tabular}{@{}rrrr@{}}\toprule
$i$ & $x_i$ & $y_i$ & $y_i-1.5$\\\midrule
0 & 0 & 0 & $-1.5$\\
1 & 1 & 1 & $-0.5$\\
2 & 2 & 2 & $0.5$\\
3 & 3 & 1 & $-0.5$\\
4 & 4 & 0 & $-1.5$\\\bottomrule
\end{tabular}\end{center}
Có khoảng $[1,2]$ và $[2,3]$. Nhánh trái tăng, nhánh phải giảm. Nếu đảo toàn bảng sẽ có các tung độ 0 và 1 lặp lại, nên không xác định một hàm ngược đơn trị trên toàn bảng.

Đối với phép lặp từ đa thức thuận, có thể giữ tập mốc đi qua nhiều nhánh; phải cảnh báo khả năng nhiều nghiệm và đối chiếu khoảng đã yêu cầu. Đối với \emph{hàm ngược}, tool yêu cầu nhánh mẫu tăng/giảm nghiêm ngặt và các tung độ phân biệt.

\section{Hàm ngược bằng Lagrange}
\label{sec:inverse-lagrange}
Đặt $u_i=y_i$, $v_i=x_i$ trên nhánh được chọn. Lập
\[
 Q(y)=\sum_{i=0}^{m-1}v_i\ell_i(y),\qquad
 \ell_i(y)=\prod_{j\ne i}\frac{y-u_j}{u_i-u_j},\qquad
 \widehat x=Q(\bar y).
\]
\textbf{Điều kiện}: các $u_i$ phân biệt, dữ liệu thuộc nhánh một-một thích hợp. Hoành độ gốc có đều hay không không ảnh hưởng tính hợp lệ của Lagrange trên bảng đảo. Mức phải nằm trong miền tung độ của \emph{tập chọn} và nhánh; ngoại suy ngược cần được bật rõ.

\subsection{Từng bước cho ví dụ $x=[0,1,2]$, $y=[0,1,4]$}
\begin{enumerate}
\item Chọn cả ba mốc của nhánh tăng; mức 2 nằm giữa 0 và 4.
\item Đảo thành $u=[0,1,4]$, $v=[0,1,2]$.
\item Với mỗi $i$, lập tích tử tại $\bar y=2$ và tích mẫu tại $u_i$.
\item Chia để được $\ell_i(2)$, nhân $v_i$, rồi cộng.
\end{enumerate}
\begin{center}\small
\begin{tabular}{@{}rrrrrr@{}}\toprule
$i$ & $u_i$ & $v_i$ & mẫu $\prod_{j\ne i}(u_i-u_j)$ & $\ell_i(2)$ & $v_i\ell_i(2)$\\\midrule
0 & 0 & 0 & 4 & $-1/2$ & 0\\
1 & 1 & 1 & $-3$ & $4/3$ & $4/3$\\
2 & 4 & 2 & 12 & $1/6$ & $1/3$\\\bottomrule
\end{tabular}\end{center}
Kết luận $\widehat x=0+4/3+1/3=5/3$. Kiểm tra $\sum_i\ell_i(2)=1$. Không có vòng hội tụ; chỉ có ba hạng hữu hạn.

\section{Hàm ngược bằng Newton tỉ sai phân}
\label{sec:inverse-newton}
Trên cùng bảng đảo $(u_i,v_i)$, tính
\[
 D_{i,0}=v_i,\qquad D_{i,k}=\frac{D_{i+1,k-1}-D_{i,k-1}}{u_{i+k}-u_i},
\]
rồi
\[
 Q(y)=D_{0,0}+D_{0,1}(y-u_0)+D_{0,2}(y-u_0)(y-u_1)+\cdots.
\]
\textbf{Điểm dễ nhầm}: dù $x_i$ cách đều, các $u_i=y_i$ thường không đều. Phải dùng \emph{tỉ sai phân}, không dùng sai phân tiến đều trên bảng đảo.

\begin{center}\begin{tabular}{@{}rrrr@{}}\toprule
$u_i$ & $v_i$ & cấp 1 & cấp 2\\\midrule
0 & 0 & 1 & $-1/6$\\
1 & 1 & $1/3$ & \\
4 & 2 & & \\\bottomrule
\end{tabular}\end{center}
Các phép chia: $(1-0)/(1-0)=1$; $(2-1)/(4-1)=1/3$; $(1/3-1)/(4-0)=-1/6$. Vì thế
\[
 Q(y)=y-\frac16y(y-1).
\]
Đánh giá Horner tại 2: bắt đầu $v=-1/6$; $v\leftarrow1+(2-1)(-1/6)=5/6$; $v\leftarrow0+(2-0)(5/6)=5/3$. Đây là hai lượt đánh giá lồng, không phải hai vòng hội tụ.

\paragraph{Kết luận của hai dạng hàm ngược.} Trên cùng tập mốc, Lagrange và Newton cho đúng một đa thức $Q$ và cùng $5/3$. Việc thêm cột $P(\widehat x)-\bar y$ vào báo cáo chỉ là thông tin đối chiếu với mô hình thuận khác, không dùng làm điều kiện bác bỏ kết quả $Q$.

\section{Lặp từ Newton tiến: lập công thức trước, sau đó mới lặp}
\label{sec:inverse-forward}
Chọn các mốc đều $x_0<\cdots<x_n$, $h>0$, gốc đầu tập tính. Đặt
\[
 t=\frac{x-x_0}{h},\quad
 P_t(t)=y_0+d\,t+H(t),\quad d=\Delta y_0,
\]
\[
 H(t)=\sum_{k=2}^{n}\frac{\Delta^ky_0}{k!}
             t(t-1)\cdots(t-k+1).
\]
Từ $P_t(t)=\bar y$ suy ra phép lặp
\[
 \boxed{t_{j+1}=\varphi(t_j)=\frac{\bar y-y_0-H(t_j)}d},\qquad
 x_{j+1}=x_0+ht_{j+1}.
\]
Mẫu $d$ phải khác 0. Không thay một mẫu bằng 0 bằng một số nhỏ tùy ý.

\subsection{Một vòng thực hiện những việc nào?}
\begin{enumerate}
\item Lấy $t_j$ của vòng trước; không lấy $x_j$ thế trực tiếp vào biểu thức theo $t$.
\item Với mọi $k\ge2$, tính tích cơ sở tại $t_j$, chia $k!$, nhân sai phân đúng hàng gốc; ghi từng đóng góp.
\item Cộng các đóng góp thành $H(t_j)$.
\item Tính tử $\bar y-y_0-H(t_j)$, chia $d$ để có $t_{j+1}$.
\item Đổi lại $x_{j+1}=x_0+ht_{j+1}$; tính bước $|h|\,|t_{j+1}-t_j|$.
\item Đánh giá lại đa thức tại \emph{điểm mới} để có phần dư $\rho_{j+1}=|P_t(t_{j+1})-\bar y|$.
\item Kiểm tra vị trí/đoạn và tiêu chuẩn dừng của chế độ đã chọn. Nếu chưa đạt, chuyển $t_{j+1}$ làm đầu vào vòng sau.
\end{enumerate}

\subsection{Ba quy tắc khởi tạo}
\begin{center}\small
\begin{tabular}{@{}p{.24\linewidth}p{.65\linewidth}@{}}\toprule
Quy tắc & Phép tính\\\midrule
Tuyến tính tại gốc & $t_0=(\bar y-y_0)/d$; bỏ phần $H$ chỉ để khởi tạo\\
Tuyến tính hai đầu $a,b$ & $x^{(0)}=a+(b-a)(\bar y-y_a)/(y_b-y_a)$, rồi $t_0=(x^{(0)}-\mathrm{origin})/h$\\
Thủ công & Nhập $t_0$; qua API có thể nhập $x^{(0)}$, nhưng phải đổi theo đúng gốc và $h$\\\bottomrule
\end{tabular}\end{center}
Nếu đầu $a,b$ không trùng mốc dùng, quy tắc hai đầu lấy giá trị \emph{mô hình thuận} tại đầu đó và ghi rõ nguồn; không giả sử biết thêm giá trị hàm thật. Khi $y_b=y_a$, quy tắc này không xác định. Không tự đổi quy tắc sau một lần thất bại.

\section{Khoảng mục tiêu T khác đoạn chứng nhận J}
\label{sec:inverse-contraction}
Từ khoảng hoành độ đang yêu cầu $[a,b]$, tính
\[
 T=\left[\frac{a-\mathrm{origin}}h,\frac{b-\mathrm{origin}}h\right].
\]
Đây là nơi muốn nhận nghiệm. Một đoạn $J$ trong biến $t$ dùng để chứng minh phép lặp hội tụ; nó có thể khác $T$.

\textbf{Chứng nhận co trên $J$ cần đủ ba điều kiện}:
\[
 t_0\in J,\qquad \varphi(J)\subseteq J,\qquad
 \sup_{t\in J}|\varphi'(t)|\le q<1.
\]
Khi đó phép lặp hội tụ đến nghiệm duy nhất của $t=\varphi(t)$ trong $J$. Vì $\varphi'=-H'/d$, sai phân bậc cao nhỏ so với $d$ có thể gợi ý co, nhưng chưa thay cho kiểm tra trên toàn đoạn.

\subsection{Cận hậu nghiệm và hai tiêu chuẩn dừng}
Với chứng nhận hợp lệ:
\[
 |x_{j+1}-x_*|\le B_{j+1}
 =|h|\frac q{1-q}|t_{j+1}-t_j|.
\]
Tool chỉ nhận khi đồng thời
\[
 B_{j+1}\le\varepsilon_x,\qquad
 \rho_{j+1}\le\varepsilon_y.
\]
Nếu $q=0$, ánh xạ là hằng trên $J$ và cận lý thuyết bằng 0 sau một cập nhật hợp lệ. Chứng nhận trên mô hình bằng số hữu tỷ không có nghĩa mọi sai số máy và sai số dữ liệu của hàm thật đã được chặn.

Ở chế độ \textbf{thực hành}, không dùng $q$ để dừng. Thay $B_{j+1}$ bằng bước $|h||t_{j+1}-t_j|$ rồi vẫn yêu cầu phần dư nhỏ. Bước nhỏ là dấu hiệu thực hành, không phải chứng minh sai số nghiệm nhỏ nếu chưa có thêm căn cứ.

\subsection{Rời $T$ và rời $J$ có nghĩa khác nhau}
Trong thực hành, một vòng rời $T$ được ghi cảnh báo rồi tiếp tục: đa thức vẫn tính được. Khi cuối cùng đạt ngưỡng, nghiệm phải thuộc khoảng yêu cầu $[a,b]$; nếu ở ngoài thì chưa giải đúng câu hỏi trong khoảng đó. Trong chế độ có chứng nhận, rời $J$ vi phạm giả thiết giữ đoạn nên phải dừng kiểm tra. Không kẹp điểm về biên để tạo cảm giác hội tụ.

\paragraph{Vì sao khoảng nghiệm ở giữa bảng không làm thiếu mốc?} Nếu đã chọn 36 mốc, gốc tiến vẫn là mốc đầu trong 36 mốc đó; khoảng nghiệm có thể nằm giữa bảng và $T$ có thể như $[10,11]$. Công thức được lập một lần từ 36 mốc rồi dùng lại. Ba mươi sáu mốc không tương ứng với 36 vòng, và không cần thêm 36 mốc tính từ đầu khoảng nghiệm.

\section{Ví dụ lặp tiến có chứng nhận: đủ tất cả ba vòng}
\label{sec:inverse-forward-example}
\textbf{Dữ liệu tự dựng}:
\[
 x=[0,1,2],\quad y=[1,3.1,5.4],\quad \bar y=2,\quad
 [a,b]=[0,1],\quad\varepsilon_x=\varepsilon_y=10^{-6}.
\]
Dữ liệu này lấy từ đa thức $1+2x+0.1x^2$, nhưng thuật toán chỉ cần bảng. Sai phân:
\begin{center}\begin{tabular}{@{}rrrr@{}}\toprule
$x_i$ & $y_i$ & $\Delta y_i$ & $\Delta^2y_i$\\\midrule
0 & 1 & 2.1 & 0.2\\
1 & 3.1 & 2.3 & \\
2 & 5.4 & & \\\bottomrule
\end{tabular}\end{center}
Gốc 0, $h=1$, $t=x$, nên
\[
 P_t(t)=1+2.1t+0.1t(t-1),\qquad
 H(t)=0.1t(t-1),\qquad
 \varphi(t)=\frac{1-0.1t(t-1)}{2.1}.
\]
Khởi tạo $t_0=1/2.1=10/21$. Chọn $J=T=[0,1]$:
\[
 \varphi(J)\subseteq\left[\frac{10}{21},\frac{41}{84}\right]\subseteq[0,1],
 \qquad \varphi'(t)=\frac{1-2t}{21},\quad q=\frac1{21}.
\]
Cả ba điều kiện co đều đạt. Cận đổi sang $x$ là $B_{j+1}=\frac1{20}|t_{j+1}-t_j|$.

\subsection{Vòng 1: viết ra từng phép thế}
\begin{align*}
H(t_0)&=0.1\frac{10}{21}\left(\frac{10}{21}-1\right)
       =-0.0249433106575964,\\
\text{tử}&=2-1-H(t_0)=1.0249433106575964,\\
t_1&=1.0249433106575964/2.1=0.488068243170284,\\
x_1&=0+1t_1=0.488068243170284,\\
B_1&=|t_1-t_0|/20=5.93888348990\cdot10^{-4},\\
\rho_1&=|1+2t_1+0.1t_1^2-2|=4.24526602993\cdot10^{-5}.
\end{align*}
Hai ngưỡng chưa đạt nên dùng $t_1$ cho vòng 2. Không dùng $H(t_0)$ tính phần dư ở điểm mới.

\subsection{Bảng đủ các vòng, không bỏ vòng chưa đạt}
\begin{center}\small
\begin{tabular}{@{}rrrrr@{}}\toprule
Vòng & $t_j$ mới & $H(t_{j-1})$ & $B_j$ & $\rho_j$\\\midrule
1 & 0.488068243170 & $-0.024943310658$ & $5.93888349\cdot10^{-4}$ & $4.24526603\cdot10^{-5}$\\
2 & 0.488088458723 & $-0.024985763318$ & $1.01077763\cdot10^{-6}$ & $4.82005445\cdot10^{-8}$\\
3 & 0.488088481675 & $-0.024985811518$ & $1.14763201\cdot10^{-9}$ & $5.46802117\cdot10^{-11}$\\\bottomrule
\end{tabular}\end{center}
Vòng 2 có phần dư đạt nhưng $B_2>10^{-6}$ nên \emph{chưa dừng}. Vòng 3 đạt cả hai. Tool trả
\[
 \widehat x=0.4880884816754477.
\]
Để kiểm riêng ví dụ tự dựng, nghiệm dương của đa thức đã biết là $-10+\sqrt{110}\approx0.4880884817015$, nằm trong cận của vòng cuối. So sánh này không phải bước bắt buộc của thuật toán.

\section{Lặp từ Newton lùi: cùng dữ liệu, khác gốc và sai phân}
\label{sec:inverse-backward}
Gốc là $x_n$, $t=(x-x_n)/h$. Đặt
\[
 P_t(t)=y_n+d\,t+H(t),\quad d=\Delta y_{n-1},\quad
 H(t)=\sum_{k=2}^n\frac{\Delta^ky_{n-k}}{k!}t(t+1)\cdots(t+k-1).
\]
Vẫn dùng $t_{j+1}=(\bar y-y_n-H(t_j))/d$, nhưng cả gốc, sai phân và các dấu trong tích phải đổi đồng bộ.

\subsection{Ví dụ lặp lùi, đủ năm vòng}
Dùng cùng ba cặp và mức 2 của mục trước. Gốc 2, $h=1$, $t=x-2$, $d=2.3$, $y_n=5.4$:
\[
 P_t(t)=5.4+2.3t+0.1t(t+1),\qquad
 \varphi(t)=\frac{-3.4-0.1t(t+1)}{2.3}.
\]
Khoảng $x\in[0,1]$ thành $T=[-2,-1]$. Khởi tạo $t_0=-3.4/2.3=-34/23$. Chọn $J=T$:
\[
 \varphi(J)\subseteq[-36/23,-34/23]\subseteq[-2,-1],\qquad
 |\varphi'(t)|=\frac{|2t+1|}{23}\le\frac3{23}=q.
\]
Cận $B_j=\frac3{20}|t_j-t_{j-1}|$. Vòng đầu:
\begin{align*}
H(t_0)&=0.1t_0(t_0+1)=0.0706994328922495,\\
t_1&=(-3.4-0.0706994328922495)/2.3=-1.508999753431413,\\
x_1&=2+t_1=0.491000246568587.
\end{align*}
\begin{center}\small
\begin{tabular}{@{}rrrrr@{}}\toprule
Vòng & $t_j$ mới & $x_j=2+t_j$ & $B_j$ & $\rho_j$\\\midrule
1 & $-1.508999753431$ & 0.491000246569 & $4.61083258\cdot10^{-3}$ & $6.10861735\cdot10^{-3}$\\
2 & $-1.511655674018$ & 0.488344325982 & $3.98388088\cdot10^{-4}$ & $5.36670035\cdot10^{-4}$\\
3 & $-1.511889008816$ & 0.488110991184 & $3.50002197\cdot10^{-5}$ & $4.72163389\cdot10^{-5}$\\
4 & $-1.511909537659$ & 0.488090462341 & $3.07932645\cdot10^{-6}$ & $4.15462427\cdot10^{-6}$\\
5 & $-1.511911344018$ & 0.488088655982 & $2.70953756\cdot10^{-7}$ & $3.65574580\cdot10^{-7}$\\\bottomrule
\end{tabular}\end{center}
Vòng 5 đạt cả hai ngưỡng $10^{-6}$. Đáp số khác các chữ số cuối của lặp tiến vì cả hai dừng ở những sai số mô hình được chứng nhận khác nhau; không phải hai đa thức thuận khác nhau.

\paragraph{Tại sao không tự đổi hướng?} Trong ví dụ này tiến có $q=1/21$, lùi có $q=3/23$, nên tiến nhanh hơn. Nhưng nếu đề yêu cầu lùi phải giữ lùi; việc chọn hướng khác là một lượt tính khác có công bố. Có dữ liệu mà khởi tạo của một hướng nằm ngoài $J$ trong khi hướng còn lại hợp lệ. Thiếu chứng nhận của một hướng không đồng nghĩa thiếu mốc.

\section{Thực hành, bài tuyến tính và các trạng thái dừng}
\label{sec:inverse-practical}
\subsection{Chế độ thực hành dùng lại cùng công thức}
Với ví dụ tiến, các bước $x$ là $1.18777670\cdot10^{-2}$, $2.02155525\cdot10^{-5}$, $2.29526402\cdot10^{-8}$. Nếu dùng ngưỡng $10^{-6}$, chế độ thực hành cũng dừng ở vòng 3, nhưng căn cứ là bước và phần dư; báo cáo không gọi bước nhỏ này là cận Banach khi chưa có chứng nhận.

Các tình huống đáng đọc:
\begin{itemize}
\item Một vài vòng có bước tăng chưa đủ kết luận phân kỳ toàn cục; tiếp tục đến giới hạn hoặc lỗi thực sự, đồng thời lưu quỹ đạo.
\item Giá trị mới bằng giá trị cũ nhưng phần dư chưa đạt là đình trệ số học, không được gọi là hội tụ.
\item Hết số vòng là chưa đạt yêu cầu; số cuối chỉ dùng xem lại lựa chọn.
\item Tràn số khi đánh giá đa thức bậc cao là lỗi số học của quỹ đạo, không phải thiếu mốc.
\end{itemize}

\subsection{Dữ liệu tuyến tính không cần tạo vòng lặp thừa}
Ví dụ $x=[0,1,2]$, $y=[1,3,5]$, mức 4. Các sai phân cấp 2 trở lên bằng 0; $H=0$ nên $t=(4-1)/2=1.5$ cho nghiệm mô hình ngay. Tool có thể trả trạng thái mô hình tuyến tính với không có vòng lặp. Tương tự, nếu mức bằng đúng một mốc được chọn trong khoảng mục tiêu, trả mốc đó mà không tạo lịch sử giả.

\subsection{Không có chứng nhận có thể là ba việc khác nhau}
\begin{enumerate}
\item \emph{Không đủ căn cứ}: bao khoảng dùng để chặn còn rộng nên chưa chứng minh được co.
\item \emph{Một điều kiện bị bác bỏ}: chẳng hạn có điểm với $|\varphi'|\ge1$ hoặc ảnh đi ra ngoài $J$.
\item \emph{Khởi tạo không thuộc $J$}: chứng nhận trên đoạn đó không áp dụng cho quỹ đạo đang bắt đầu.
\end{enumerate}
Cả ba đều phải nêu đúng nguyên nhân. Chế độ chỉ nhận chứng nhận không tự đổi sang thực hành, không tự chuyển chia đôi và không tự đổi khởi tạo.

\section{Chia đôi đa thức thuận: mỗi vòng giữ lại một nửa khoảng}
\label{sec:inverse-bisection}
Lập đa thức thuận $P$ bằng tỉ sai phân từ tập mốc đã chọn. Đặt $F(x)=P(x)-\bar y$. Đầu vào thêm khoảng $a<b$, ngưỡng $\varepsilon_x>0$ và số vòng tối đa.

\textbf{Điều kiện}: $F$ liên tục trên khoảng và hai đầu trái dấu, hoặc có đầu bằng 0. Đa thức luôn liên tục. Với hàm thật chưa biết, kết luận đang xét $F$ của \emph{mô hình}, không tự mở rộng bảo đảm sang $f$.
\begin{enumerate}
\item Tính $F(a),F(b)$. Có đầu bằng 0 thì trả đầu đó. Nếu cùng dấu, chưa đủ điều kiện chia đôi trên khoảng này.
\item Khi $(b-a)/2>\varepsilon_x$, tính $c=(a+b)/2$ và $F(c)$.
\item Nếu $F(c)=0$, thu khoảng về $[c,c]$.
\item Nếu $F(a)F(c)<0$, thay $b\leftarrow c$; ngược lại thay $a\leftarrow c$. Lưu cả khoảng trước và sau cập nhật.
\item Sau khi nửa độ rộng đạt ngưỡng, trả \emph{trung điểm của khoảng cuối}, kèm khoảng và cận nửa độ rộng.
\end{enumerate}
Sau $N$ lần giữ nửa khoảng, độ rộng còn $(b_0-a_0)/2^N$. Để nửa độ rộng không quá $\varepsilon_x$ cần
\[
 N\ge\max\left(0,\left\lceil\log_2\frac{b_0-a_0}{2\varepsilon_x}\right\rceil\right).
\]
Đây là cận vị trí nghiệm của mô hình trong khoảng, không yêu cầu phần dư tung độ bằng $\varepsilon_x$. Đơn vị hai đại lượng khác nhau.

\subsection{Ví dụ đủ năm lần: giải $x^2=2$ bằng dữ liệu bảng}
Nhập $x=[0,1,2]$, $y=[0,1,4]$; tool dựng đúng $P=x^2$. Mức 2, khoảng $[1,2]$, ngưỡng $\varepsilon_x=0.02$. Ban đầu $F(1)=-1$, $F(2)=2$.
\begin{center}\small
\begin{tabular}{@{}rrrrrr@{}}\toprule
Vòng & $a$ trước & $b$ trước & $c$ đã thử & $F(c)$ & khoảng sau cập nhật\\\midrule
1 & 1 & 2 & 1.5 & 0.25 & $[1,1.5]$\\
2 & 1 & 1.5 & 1.25 & $-0.4375$ & $[1.25,1.5]$\\
3 & 1.25 & 1.5 & 1.375 & $-0.109375$ & $[1.375,1.5]$\\
4 & 1.375 & 1.5 & 1.4375 & 0.06640625 & $[1.375,1.4375]$\\
5 & 1.375 & 1.4375 & 1.40625 & $-0.0224609375$ & $[1.40625,1.4375]$\\\bottomrule
\end{tabular}\end{center}
Khoảng cuối có nửa độ rộng $0.015625\le0.02$. Đáp số tool là trung điểm cuối
\[
 \widehat x=\frac{1.40625+1.4375}{2}=1.421875,
 \qquad |\widehat x-\sqrt2|\le0.015625.
\]
Không lấy $c=1.40625$ của vòng vừa thử làm kết quả cuối: mã hiện trả trung điểm khoảng sau khi đã thu. Nếu float không còn phân biệt trung điểm với đầu khoảng trong khi chưa đạt ngưỡng, tool báo giới hạn độ phân giải thay vì lặp vô hạn.

\section{Sai số mô hình không phải sai số của hàm thật}
\label{sec:inverse-errors}
\subsection{Với nội suy hàm ngược}
Nếu hàm thật có nghịch đảo $g$ đủ trơn và $|g^{(m)}|\le M_m$ trên đoạn chứa các tung độ dùng cùng mức $\bar y$:
\[
 |g(\bar y)-Q(\bar y)|\le\frac{M_m}{m!}
            \left|\prod_i(\bar y-y_i)\right|.
\]
Đây là cận theo hàm ngược, không phải cận $|f^{(m)}|$. Khi $f'$ rất nhỏ, hàm ngược có thể nhạy số vì $g'=1/f'$. Tung độ gần trùng được cảnh báo nhưng không tự gộp thành một mốc.

\subsection{Với nghiệm của đa thức thuận}
Giả sử trên đoạn chứa nghiệm thật và điểm tính có $|f-P|\le E$ và $|f'|\ge\mu>0$ trên cùng nhánh thích hợp. Nếu $\rho=|P(\widehat x)-\bar y|$ thì
\[
 |\widehat x-x_{\rm thật}|\le\frac{E+\rho}{\mu}.
\]
Phải nêu nguồn và miền áp dụng của $E,\mu$. Không lấy cận co của phép lặp thay cho $E$; cận co chỉ chặn sai số \emph{giải phương trình mô hình}. Việc sai số tại các mốc bằng 0 cũng không suy ra $E=0$ giữa mốc.

\paragraph{Những số nên chép cùng đáp án.} Tập mốc/nhánh, khoảng $[a,b]$, gốc và $h$, công thức $P_t$ hoặc $Q$, quy tắc khởi tạo, $\varepsilon_x,\varepsilon_y$, trạng thái dừng, số vòng, phần dư cuối và phạm vi của cận được công bố. Số chữ số hiển thị là lựa chọn trình bày, không là số chữ số được bảo đảm đúng.

\section{Bảng lỗi để quyết định thao tác tiếp theo}
\begin{center}\small
\begin{longtable}{@{}p{.30\linewidth}p{.59\linewidth}@{}}
\toprule
Thông tin tool báo & Diễn giải và việc cần làm\\\midrule\endhead
Hoành độ trùng & Dữ liệu thuận không xác định nội suy duy nhất theo chính sách giữ cặp; xử lý nguồn\\
Tung độ trùng khi đảo bảng & Mẫu tỉ sai phân ngược bằng 0 hoặc nhiều $x$ cho một $y$; tách nhánh phù hợp\\
Nhánh không đơn điệu nghiêm ngặt trên mẫu & Không có một hàm ngược đơn trị trên tập đang xét; không sắp toàn bảng theo $y$ để che đổi chiều\\
Mức ngoài miền tung độ & Đang cần ngoại suy ngược; chỉ chạy khi đã chọn rõ và hiểu giới hạn\\
Không đủ/chưa chọn đúng mốc cho bậc & Bậc $p$ cần đúng $p+1$ cặp đã chọn; không liên quan số vòng\\
Sai phân đầu $d=0$ & Không chia được khi lập phép lặp này; phải chọn lượt khác theo đề, không thay mẫu bằng số nhỏ\\
Chưa đủ chứng nhận & Kiểm $t_0\in J$, ảnh giữ $J$, $q<1$; chưa chứng minh không đồng nghĩa chắc chắn phân kỳ\\
Rời $T$ trong thực hành & Ghi cảnh báo và tiếp tục; kết quả cuối vẫn phải thuộc khoảng yêu cầu\\
Rời $J$ đã chứng nhận & Dừng do vi phạm chứng nhận/số học; không kẹp về biên\\
Đạt ngưỡng ngoài khoảng yêu cầu & Chưa phải nghiệm của câu hỏi trong khoảng đã chọn\\
Hết vòng hoặc đình trệ & Giữ lịch sử; chưa trả kết quả đạt yêu cầu\\
Hai đầu chia đôi cùng dấu & Khoảng này chưa đủ điều kiện; không suy phương trình vô nghiệm ở mọi nơi\\
\bottomrule
\end{longtable}\end{center}

\paragraph{Nguồn đối chiếu.} \texttt{pps/inverse\_solution.py} giữ lựa chọn và điều khiển từng bước; \texttt{inverse\_model.py} dựng hệ số, các đóng góp và chứng nhận trên đoạn; \texttt{inverse\_presentation.py} trình bày từ vết tính đó. Đặc tả nội suy ngược bản 2 trong \texttt{kienthuc/real} là nguồn yêu cầu về giữ tập mốc, tách $T/J$ và phân loại trạng thái. Các ví dụ trong chương này đều là dữ liệu tự dựng, đã chạy qua đúng bộ giải của app; không nhận là dữ liệu đề thi thật.

\chapter{Spline: dựng từng mảnh và giữ đúng điều kiện biên}
\label{ch:spline}
\section{Từ ô nhập đến các mảnh spline: khuôn làm cho mọi dữ liệu}
\subsection*{Input: các ô đúng tên trên giao diện}
Bảng này được lấy từ biểu mẫu của tool. Giá trị mặc định là ví dụ nhập, cần thay bằng dữ liệu đề. Hai dãy $x,y$ là bảng dữ liệu riêng bên dưới nếu chức năng dùng bảng.
\subsubsection*{Ô chung của các lựa chọn trong nhóm}
\begin{center}\small\begin{tabular}{@{}p{.32\linewidth}p{.17\linewidth}p{.43\linewidth}@{}}\toprule
Ô nhập & Khóa JSON & Cách nhập / mặc định\\\midrule
Điểm cần tính x (có thể để trống) & \texttt{target} & \texttt{1.5}\\[3pt]
Cấp đạo hàm (0 = giá trị) & \texttt{order} & \texttt{0}\\[3pt]
Cận tích phân a (không bắt buộc) & \texttt{integral\_a} & Để trống\\[3pt]
Cận tích phân b (không bắt buộc) & \texttt{integral\_b} & Để trống\\[3pt]
Cách giải hệ biên & \texttt{solver} & \texttt{Tự chọn}\\[3pt]
Mức trình bày & \texttt{detail\_level} & \texttt{Tổng hợp}\\[3pt]
\bottomrule\end{tabular}\end{center}
\subsubsection*{Ô riêng: Bậc hai — biên trái}
\begin{center}\small\begin{tabular}{@{}p{.32\linewidth}p{.17\linewidth}p{.43\linewidth}@{}}\toprule
Ô nhập & Khóa JSON & Cách nhập / mặc định\\\midrule
Giá trị đạo hàm tại biên trái & \texttt{left} & \texttt{0}\\[3pt]
\bottomrule\end{tabular}\end{center}
\subsubsection*{Ô riêng: Bậc hai — biên phải}
\begin{center}\small\begin{tabular}{@{}p{.32\linewidth}p{.17\linewidth}p{.43\linewidth}@{}}\toprule
Ô nhập & Khóa JSON & Cách nhập / mặc định\\\midrule
Giá trị đạo hàm tại biên phải & \texttt{right} & \texttt{0}\\[3pt]
\bottomrule\end{tabular}\end{center}
\subsubsection*{Ô riêng: Cubic kẹp /\allowbreak{} Cubic đạo hàm bậc hai}
\begin{center}\small\begin{tabular}{@{}p{.32\linewidth}p{.17\linewidth}p{.43\linewidth}@{}}\toprule
Ô nhập & Khóa JSON & Cách nhập / mặc định\\\midrule
Giá trị đạo hàm tại biên trái & \texttt{left} & \texttt{0}\\[3pt]
Giá trị đạo hàm tại biên phải & \texttt{right} & \texttt{0}\\[3pt]
\bottomrule\end{tabular}\end{center}
\subsubsection*{Ô riêng: Cubic hỗn hợp}
\begin{center}\small\begin{tabular}{@{}p{.32\linewidth}p{.17\linewidth}p{.43\linewidth}@{}}\toprule
Ô nhập & Khóa JSON & Cách nhập / mặc định\\\midrule
Cấp đạo hàm biên trái & \texttt{left\_order} & \texttt{1}\\[3pt]
Cấp đạo hàm biên phải & \texttt{right\_order} & \texttt{2}\\[3pt]
Giá trị tại biên trái & \texttt{left} & \texttt{0}\\[3pt]
Giá trị tại biên phải & \texttt{right} & \texttt{2}\\[3pt]
\bottomrule\end{tabular}\end{center}
\subsubsection*{Ô riêng: Bậc 4 — đạo hàm trái}
\begin{center}\small\begin{tabular}{@{}p{.32\linewidth}p{.17\linewidth}p{.43\linewidth}@{}}\toprule
Ô nhập & Khóa JSON & Cách nhập / mặc định\\\midrule
S′ S′′ S′′′ ở đầu trái & \texttt{derivatives} & \texttt{0 2 0}\\[3pt]
\bottomrule\end{tabular}\end{center}
\subsubsection*{Ô riêng: Biên tổng quát}
\begin{center}\small\begin{tabular}{@{}p{.32\linewidth}p{.17\linewidth}p{.43\linewidth}@{}}\toprule
Ô nhập & Khóa JSON & Cách nhập / mặc định\\\midrule
Bậc spline & \texttt{degree} & \texttt{4}\\[3pt]
Điều kiện biên (JSON) & \texttt{conditions} & \texttt{[["trái",\allowbreak{}1,\allowbreak{}0],\allowbreak{}["trái",\allowbreak{}2,\allowbreak{}2],\allowbreak{}["trái",\allowbreak{}3,\allowbreak{}0]]}
\par Mỗi điều kiện ghi phía,\allowbreak{} cấp đạo hàm,\allowbreak{} giá trị;\allowbreak{} giữ đúng các điều kiện đề cho.\\[3pt]
\bottomrule\end{tabular}\end{center}

\paragraph{Input chung.} Bảng $n+1$ cặp $(x_i,y_i)$ phân biệt, sắp tăng theo $x$; không cần cách đều. Chọn đúng tên biên, nhập giá trị đạo hàm mà đề cho. Điểm đánh giá và cặp cận tích phân là yêu cầu sau khi đã dựng spline, không phải điều kiện biên bổ sung. Muốn tích phân cần nhập cả hai cận.

\subsection{Bước 1 -- bảng khoảng mốc tổng quát}
\begin{center}\small\begin{tabular}{cllll}\toprule
Đoạn&Hai đầu&Độ dài $h_i$&Hiệu tung độ&Độ dốc $\delta_i$\\\midrule
0&$x_0,x_1$&$x_1-x_0$&$y_1-y_0$&$(y_1-y_0)/h_0$\\
$i$&$x_i,x_{i+1}$&$x_{i+1}-x_i$&$y_{i+1}-y_i$&$(y_{i+1}-y_i)/h_i$\\
$n-1$&$x_{n-1},x_n$&$x_n-x_{n-1}$&$y_n-y_{n-1}$&$(y_n-y_{n-1})/h_{n-1}$\\\bottomrule
\end{tabular}\end{center}
Các mảnh dùng $u=x-x_i$, không phải $t=u/h_i$. Những bước tiếp theo nhận bảng $h_i,\delta_i$ này.

\subsection{Bước 2 -- chọn quy trình theo đúng loại spline}
\begin{center}\small\begin{tabular}{@{}p{.28\linewidth}p{.30\linewidth}p{.33\linewidth}@{}}\toprule
Tool được chọn&Input biên cần thêm&Các đại lượng phải tìm\\\midrule
Tuyến tính&Không thêm biên&$S_i=y_i+\delta_i u$\\
Bậc hai -- biên trái&$m_0=S'(x_0)$&Truyền $m_{i+1}=2\delta_i-m_i$ từ trái\\
Bậc hai -- biên phải&$m_n=S'(x_n)$&Truyền $m_i=2\delta_i-m_{i+1}$ từ phải\\
Cubic tự nhiên&Không nhập số biên&$M_0=M_n=0$, giải các moment còn lại\\
Cubic kẹp&$S'(x_0),S'(x_n)$&Giải $M_i=S''(x_i)$ với hàng biên độ dốc\\
Cubic đạo hàm bậc hai&$S''(x_0),S''(x_n)$&Ấn định hai moment biên\\
Cubic hỗn hợp&Cấp 1/2 riêng ở từng đầu và giá trị&Chọn hàng biên theo từng đầu\\
Cubic tuần hoàn&$y_0=y_n$&Khớp thêm $S'(x_0)=S'(x_n)$, $M_0=M_n$\\
Cubic not-a-knot&Không nhập số biên&Ghép đạo hàm cấp ba tại hai nút gần biên\\
Bậc 4 -- đạo hàm trái&$S',S'',S'''$ ở đầu trái&Truyền ba đạo hàm qua từng đoạn\\
Biên tổng quát&Bậc $p$, danh sách điều kiện JSON&Lập hệ hệ số mảnh, kiểm hạng/tương thích\\\bottomrule
\end{tabular}\end{center}

\paragraph{Khuôn bảng bậc hai.}
\begin{center}\small\begin{tabular}{cllll}\toprule
Đoạn&$\delta_i$&$m_i$&$c_i=(\delta_i-m_i)/h_i$&Mảnh trả ra\\\midrule
$i$&Từ bước 1&Nhận/truy ngược&Sau khi biết $m_i$&$S_i=y_i+m_iu+c_iu^2$\\\bottomrule
\end{tabular}\end{center}
Biên trái đi $i=0,1,\ldots,n-1$; biên phải đi ngược $i=n-1,\ldots,0$. Sau khi đủ $m_i$, công thức mảnh giống nhau.

\paragraph{Khuôn bảng cubic.} Tại nút trong lập
\[
h_{i-1}M_{i-1}+2(h_{i-1}+h_i)M_i+h_iM_{i+1}=6(\delta_i-\delta_{i-1}).
\]
Thêm hai hàng biên đúng loại rồi giải. Từ nghiệm $M$ lập bảng tổng quát:
\begin{center}\small\begin{tabular}{cllll}\toprule
Đoạn&$a_i$&$b_i$&$c_i$&$d_i$\\\midrule
$i$&$y_i$&$\delta_i-h_i(2M_i+M_{i+1})/6$&$M_i/2$&$(M_{i+1}-M_i)/(6h_i)$\\\bottomrule
\end{tabular}\end{center}
Mảnh là $S_i(x_i+u)=a_i+b_iu+c_iu^2+d_iu^3$. Hệ moment, hàng biên và các hệ số là ba bảng khác nhau; không lấy một hàng hệ phương trình làm hệ số mảnh.

\paragraph{Khuôn bảng bậc p biết đạo hàm trái.} Nhận $v_r=S^{(r)}(x_i)$, $r=1,\ldots,p-1$:
\[
a_{i,0}=y_i,\quad a_{i,1}=v_1,\quad a_{i,2}=v_2/2!,\quad\ldots,\quad a_{i,p-1}=v_{p-1}/(p-1)!.
\]
Tính $K_i=a_{i,0}+a_{i,1}h_i+\cdots+a_{i,p-1}h_i^{p-1}$, rồi $a_{i,p}=(y_{i+1}-K_i)/h_i^p$. Tính tất cả đạo hàm đầu phải từ cùng bộ hệ số trước khi truyền sang đoạn kế.
\begin{center}\small\begin{tabular}{clll}\toprule
Đoạn&Đạo hàm nhận&Hệ số mảnh tìm được&Đạo hàm truyền tiếp\\\midrule
$i$&$(v_1,\ldots,v_{p-1})$&$(a_{i,0},\ldots,a_{i,p})$&$(S_i'(x_{i+1}),\ldots,S_i^{(p-1)}(x_{i+1}))$\\\bottomrule
\end{tabular}\end{center}
Biên phải/trộn của bậc 4 nhập qua \emph{Biên tổng quát}; không dùng công thức truyền trái với các số biên phải. Điều kiện tổng quát được chuyển thành hệ trên mọi $a_{i,r}$; đủ số điều kiện chưa thay việc kiểm hạng.

\subsection{Bước 3 -- đánh giá và output}
\begin{enumerate}
\item Điểm $z$ thuộc đoạn $i$: lấy $u=z-x_i$, tính
\[
S(z)=a_{i,0}+a_{i,1}u+\cdots+a_{i,p}u^p.
\]
\item Đạo hàm: $S_i'(z)=a_{i,1}+2a_{i,2}u+\cdots+pa_{i,p}u^{p-1}$; cấp khác lấy đúng giai thừa. Tại nút trong phải kiểm hai phía trước khi trả đạo hàm chung.
\item Tích phân qua nhiều đoạn: cắt tại các nút; đoạn con từ $u$ đến $v$ đóng góp
\[
a_{i,0}(v-u)+\frac{a_{i,1}}2(v^2-u^2)+\cdots+\frac{a_{i,p}}{p+1}(v^{p+1}-u^{p+1}).
\]
\end{enumerate}
\textbf{Output}: loại spline/biên, bảng hệ và nghiệm nếu có, hệ số từng mảnh theo biến cục bộ, kiểm khớp và độ trơn, giá trị/đạo hàm/tích phân được hỏi, trạng thái lỗi nếu thiếu/thừa mâu thuẫn điều kiện. Đây là phép dựng hữu hạn; không có ngưỡng hội tụ spline. Các mục bên dưới viết rõ từng hàng biên và ví dụ số.


\section{Nhìn đúng bài toán trước khi bấm tính}
\label{sec:spline-map}
Cho các cặp $(x_i,y_i)$, $i=0,\ldots,n$, với $x_0<\cdots<x_n$.
Spline là một hàm gồm $n$ mảnh: mảnh $S_i$ áp dụng trên $[x_i,x_{i+1}]$.
Mỗi mảnh đi qua hai đầu đoạn; tại nút trong, các mảnh được ghép theo cấp trơn đã chọn.
\emph{Không cần các mốc cách đều.} Spline nội suy đi qua dữ liệu; nó không tự làm trơn nhiễu đo.

Trong chương này, spline bậc $p$ có mỗi mảnh bậc \emph{không quá} $p$ và ghép trơn
$C^{p-1}$. Một mảnh có thể thực tế chỉ là đường thẳng dù đang dựng spline bậc 4.
Đó là kết quả hợp lệ khi những hệ số cao bằng 0.

\begin{center}\small
\begin{tabular}{p{.28\linewidth}p{.32\linewidth}p{.31\linewidth}}
\toprule
Đề cho/yêu cầu & Chọn trong tool & Cần thêm gì?\\ \midrule
Nối các điểm bằng đoạn thẳng & Tuyến tính & Không cần biên bổ sung\\
Parabol, biết độ dốc ở một đầu & Bậc hai -- biên trái/phải & Một đạo hàm cấp 1 ở đúng đầu đó\\
Cubic, hai đầu có độ cong 0 & Cubic tự nhiên & $S''(x_0)=S''(x_n)=0$ là giả thiết được chọn\\
Cubic, cho hai độ dốc & Cubic kẹp & $S'(x_0),S'(x_n)$\\
Cubic, cho hai độ cong & Cubic đạo hàm bậc hai & $S''(x_0),S''(x_n)$\\
Cubic, mỗi đầu cho cấp khác & Cubic hỗn hợp & Mỗi cấp là 1 hoặc 2, cùng giá trị tương ứng\\
Cubic tuần hoàn/not-a-knot & Đúng tên loại biên & Điều kiện nối chu kỳ/điều kiện cấp 3 đặc biệt\\
Bậc 4, biết ba đạo hàm trái & Bậc 4 -- đạo hàm trái & $S',S'',S'''$ tại $x_0$\\
Bậc khác, hoặc biên phải/trộn & Biên tổng quát & Bậc $p$ và các điều kiện độc lập\\
\bottomrule
\end{tabular}\end{center}

\subsection{Những ô nhập có ý nghĩa gì?}
\begin{enumerate}
\item \textbf{Hai dãy $x,y$:} các số hữu hạn, cùng độ dài, ít nhất hai cặp. Tool sắp cả cặp theo $x$; không tự gộp mốc trùng. Hai tung độ khác nhau tại cùng $x$ không thể là một giá trị của hàm nội suy.
\item \textbf{Loại spline và biên:} xác định một bài toán cụ thể. Không lấy kết quả tự nhiên để trả lời bài có biên kẹp.
\item \textbf{Điểm tính và cấp đạo hàm:} để trống điểm nếu chỉ dựng spline. Cấp 0 là giá trị, cấp 1 là đạo hàm một, v.v.
\item \textbf{Hai cận tích phân:} nhập đồng thời nếu cần tích phân spline; tool chia tích phân qua đúng các mảnh.
\item \textbf{Bộ giải:} Thomas phù hợp hệ ba đường chéo thông thường; Gauss dùng cho hệ đầy đủ. Tự chọn giữ cấu trúc thuận tiện khi có thể. Những ngoại lệ thực tế được nêu ở mục~\ref{sec:spline-limits}.
\item \textbf{Số chữ số:} phục vụ hiển thị. Nó không phải cam kết số chữ số đúng của hàm thật.
\end{enumerate}

\section{Bước chuẩn bị chung: tính khoảng mốc và độ dốc}
\label{sec:spline-prepare}
Đặt, cho từng đoạn $i=0,\ldots,n-1$,
\[
h_i=x_{i+1}-x_i>0,\qquad
\Delta y_i=y_{i+1}-y_i,\qquad \delta_i=\frac{\Delta y_i}{h_i}.
\]
Không lấy một $h$ chung khi lưới không đều. Đặt biến cục bộ
\[
u=x-x_i,\qquad 0\leq u\leq h_i.
\]
Mọi hệ số mảnh trong chương đều nhân các lũy thừa của \emph{$u$ này}, không phải $x$ toàn cục.

\paragraph{Bảng dùng xuyên suốt nhiều ví dụ.}
Với $(x,y)=(0,0),(1,1),(2,0)$:
\begin{center}
\begin{tabular}{rrrrrr}
\toprule
$i$ & $x_i$ & $x_{i+1}$ & $h_i$ & $\Delta y_i$ & $\delta_i$\\ \midrule
0&0&1&1&1&1\\
1&1&2&1&$-1$&$-1$\\ \bottomrule
\end{tabular}\end{center}
Từ đúng cùng bảng này, các điều kiện biên khác nhau cho spline khác nhau.
Đó là thay đổi bài toán, không phải bằng chứng rằng một trong các đáp số bị sai.

\section{Spline tuyến tính: không có ẩn biên}
\label{sec:spline-linear}
\textbf{Đầu vào:} chỉ bảng điểm. \textbf{Đầu ra:} các dây cung
\[
S_i(x_i+u)=y_i+\delta_i u.
\]
\begin{enumerate}
\item Tính $h_i,\delta_i$ theo mục~\ref{sec:spline-prepare}.
\item Lấy hệ số tự do bằng $y_i$, hệ số $u$ bằng $\delta_i$.
\item Tìm đoạn chứa điểm hỏi $z$, đặt $u=z-x_i$ rồi thế số.
\end{enumerate}
Trong ví dụ chung,
\[
S_0(x)=x\quad(0\leq x\leq1),\qquad
S_1(1+u)=1-u\quad(0\leq u\leq1).
\]
Tại $z=1/2$, dùng mảnh 0: $S(1/2)=1/2$.
Tại nút $x=1$, đạo hàm trái bằng $1$, đạo hàm phải bằng $-1$; $S'(1)$ không tồn tại.
Không lấy riêng đạo hàm của mảnh phải rồi gọi đó là đạo hàm chung của spline.
Đây là phép dựng hữu hạn, không có tiêu chuẩn hội tụ theo vòng lặp.

\section{Spline bậc hai: truyền độ dốc từ đúng đầu đã biết}
\label{sec:spline-quadratic}
Đặt $m_i=S'(x_i)$ và viết $S_i(x_i+u)=y_i+m_i u+c_i u^2$.
Điều kiện đi qua đầu phải cho $c_i=(\delta_i-m_i)/h_i$.
Đạo hàm tại đầu phải cho
\[
m_{i+1}=m_i+2c_i h_i=2\delta_i-m_i,
\qquad
c_i=\frac{m_{i+1}-m_i}{2h_i}.
\]
Quan hệ này vừa dựng mảnh vừa ghép đạo hàm cấp 1 giữa hai mảnh.

\subsection{Biên trái: biết $m_0$}
\label{sec:spline-quadratic-left}
\begin{enumerate}
\item Ghi $m_0=s_L$ từ đề, không tự thay bằng 0.
\item Với $i=0,1,\ldots,n-1$, tính $m_{i+1}=2\delta_i-m_i$.
\item Tính $c_i$ rồi viết mảnh. Mỗi vòng dùng $m_i$ đã có để tạo $m_{i+1}$ cho vòng sau.
\end{enumerate}
Ví dụ chung, cho $S'(0)=0$:
\begin{center}
\begin{tabular}{rrrrr}
\toprule
Đoạn & $\delta_i$ & Độ dốc đã biết & Phép tính độ dốc mới & $c_i$\\ \midrule
0&1&$m_0=0$&$m_1=2(1)-0=2$&1\\
1&$-1$&$m_1=2$&$m_2=2(-1)-2=-4$&$-3$\\ \bottomrule
\end{tabular}\end{center}
Do đó
\[
S_0(x)=x^2,\qquad S_1(1+u)=1+2u-3u^2.
\]
Tại $1.5$, $u=0.5$, nên $S(1.5)=1+1-3/4=5/4$.
Kiểm tra tại nút 1: cả hai giá trị bằng 1 và cả hai đạo hàm bằng 2.

\subsection{Biên phải: biết $m_n$}
\label{sec:spline-quadratic-right}
Lần này ghi $m_n=s_R$ rồi duyệt \emph{ngược}:
\[
\boxed{m_i=2\delta_i-m_{i+1},\qquad i=n-1,n-2,\ldots,0.}
\]
Sau khi có mọi $m_i$, vẫn dựng mảnh theo $u=x-x_i$ từ đầu trái như trên.
Với cùng dữ liệu và $m_2=-4$:
\[
m_1=2(-1)-(-4)=2,\qquad m_0=2(1)-2=0.
\]
Nhận lại đúng hai mảnh vừa tính. Nếu thay biên phải bằng số khác, kết quả thường đổi.
Không gán độ dốc phải vào vị trí $m_0$.

\section{Cubic: từ các moment đến bốn hệ số mỗi mảnh}
\label{sec:spline-cubic}
Đặt $M_i=S''(x_i)$, gọi là moment. Khi biết hai moment của đoạn,
\begin{align}
S_i(x_i+u)&=a_i+b_i u+c_i u^2+d_i u^3,\\
a_i&=y_i,&
b_i&=\delta_i-\frac{h_i}{6}(2M_i+M_{i+1}),\nonumber\\
c_i&=\frac{M_i}{2},&
d_i&=\frac{M_{i+1}-M_i}{6h_i}.\label{eq:spline-cubic-coef}
\end{align}
Nội suy và liên tục cấp 2 đã được xây vào công thức; ghép đạo hàm cấp 1 ở nút trong tạo hệ
\[
h_{i-1}M_{i-1}+2(h_{i-1}+h_i)M_i+h_iM_{i+1}
=6(\delta_i-\delta_{i-1}),\quad i=1,\ldots,n-1.
\]
\textbf{Thứ tự thực hiện:}
\begin{enumerate}
\item Tính $h_i,\delta_i$.
\item Viết $n-1$ hàng nút trong theo công thức trên.
\item Thêm đúng hai hàng biên. Phần tiếp theo cho từng loại hàng.
\item Giải hệ tìm $M_0,\ldots,M_n$.
\item Trên từng đoạn, tính lần lượt $a_i,b_i,c_i,d_i$ bằng~\eqref{eq:spline-cubic-coef}.
\item Kiểm các giá trị và đạo hàm ở hai phía nút; cuối cùng mới thế điểm hỏi.
\end{enumerate}
Không cần vòng hội tụ để dựng cubic. Độ tin cậy của phép giải phụ thuộc hệ, thang số liệu và sai số làm tròn; phần dư nhỏ không tự là cận sai số với hàm thật.

\subsection{Biên tự nhiên: hai độ cong bằng 0}
\label{sec:spline-natural}
Hai hàng biên là $M_0=0,M_n=0$. Đây là giả thiết về \emph{độ cong}, không phải hai độ dốc bằng 0.
Ví dụ ba mốc chung có hàng trong
\[
M_0+4M_1+M_2=6(-1-1)=-12.
\]
Suy ra $M=(0,-3,0)$, rồi
\[
S_0(x)=\frac32x-\frac12x^3,\qquad
S_1(1+u)=1-\frac32u^2+\frac12u^3.
\]
Vì vậy $S(1.5)=1-3/8+1/16=11/16$.
Đạo hàm trái tại đầu bảng bằng $3/2$, không bằng 0.

\subsection{Thomas: khử xuôi và thế ngược trên hệ ba đường chéo}
\label{sec:spline-thomas}
Viết hàng $i$ là $l_iM_{i-1}+b_iM_i+q_iM_{i+1}=r_i$.
Khởi tạo $\widetilde b_0=b_0$, $\widetilde r_0=r_0$. Tại mỗi hàng $i\geq1$:
\[
\lambda_i=\frac{l_i}{\widetilde b_{i-1}},\quad
\widetilde b_i=b_i-\lambda_i q_{i-1},\quad
\widetilde r_i=r_i-\lambda_i\widetilde r_{i-1}.
\]
Thế ngược:
\[
M_n=\frac{\widetilde r_n}{\widetilde b_n},\qquad
M_i=\frac{\widetilde r_i-q_iM_{i+1}}{\widetilde b_i}.
\]
Phải kiểm pivot khác 0. Không dùng các công thức ba đường chéo này cho một ma trận còn có hệ số ngoài băng mà không xử lý chúng.

\paragraph{Ví dụ bốn mốc, đủ mọi vòng.}
Cho $x=(0,1,2,3)$, $y=(0,1,0,1)$ và biên tự nhiên. Hệ là
\[
\begin{bmatrix}1&0&0&0\\1&4&1&0\\0&1&4&1\\0&0&0&1\end{bmatrix}
\begin{bmatrix}M_0\\M_1\\M_2\\M_3\end{bmatrix}
=\begin{bmatrix}0\\-12\\12\\0\end{bmatrix}.
\]
\begin{center}
\begin{tabular}{rrrrrr}
\toprule
$i$&$l_i$&$\lambda_i$&$\widetilde b_i$&$\widetilde r_i$&$q_i/\widetilde b_i$\\ \midrule
0&0&---&1&0&0\\
1&1&1&4&$-12$&$1/4$\\
2&1&$1/4$&$15/4$&15&$4/15$\\
3&0&0&1&0&0\\ \bottomrule
\end{tabular}\end{center}
Các phép thế ngược, theo đúng thứ tự, là
\[
M_3=0/1=0,\quad M_2=(15-0)/(15/4)=4,
\quad M_1=(-12-4)/4=-4,\quad M_0=0.
\]
Không đổi dấu $M_1$ vì nó là một độ cong âm hợp lệ.
Từ $M=(0,-4,4,0)$, bảng bốn hệ số là
\begin{center}
\begin{tabular}{rrrrrr}
\toprule
Đoạn & Miền $x$ & $a_i$ & $b_i$ & $c_i$ & $d_i$\\ \midrule
0&$[0,1]$&0&$5/3$&0&$-2/3$\\
1&$[1,2]$&1&$-1/3$&$-2$&$4/3$\\
2&$[2,3]$&0&$-1/3$&2&$-2/3$\\ \bottomrule
\end{tabular}\end{center}
Chẳng hạn đoạn 1:
$b_1=-1-(2(-4)+4)/6=-1/3$, $c_1=-4/2=-2$, $d_1=(4-(-4))/6=4/3$.
Tại $z=1.5$, $u=0.5$, nên $S(z)=1-1/6-1/2+1/6=1/2$.

\subsection{Biên kẹp: cho độ dốc ở hai đầu}
\label{sec:spline-clamped}
Đầu vào thêm $s_L=S'(x_0)$, $s_R=S'(x_n)$. Hai hàng biên là
\[
2h_0M_0+h_0M_1=6(\delta_0-s_L),\qquad
h_{n-1}M_{n-1}+2h_{n-1}M_n=6(s_R-\delta_{n-1}).
\]
Ví dụ ba mốc chung, $s_L=2,s_R=-2$:
\[
\begin{bmatrix}2&1&0\\1&4&1\\0&1&2\end{bmatrix}M
=\begin{bmatrix}-6\\-12\\-6\end{bmatrix},\qquad M=(-2,-2,-2)^T.
\]
Thomas cho $\lambda_1=1/2$, pivot mới $7/2$, vế phải $-9$;
$\lambda_2=2/7$, pivot mới $12/7$, vế phải $-24/7$.
Thế ngược cho ba moment $-2$.
Các mảnh là $S_0(x)=2x-x^2$, $S_1(1+u)=1-u^2$.
Tại $1/2$, $S=3/4$; tại hai đầu, đạo hàm lần lượt bằng $2,-2$, đúng đề.

\subsection{Cho đạo hàm cấp hai ở hai đầu}
\label{sec:spline-second}
Nếu đề cho $L=S''(x_0)$, $R=S''(x_n)$ thì dùng trực tiếp $M_0=L,M_n=R$.
Các hàng trong và công thức mảnh giữ nguyên.
Vẫn bảng ba mốc, $L=R=-2$:
\[
M_0=-2,\quad -2+4M_1-2=-12\ \Longrightarrow\ M_1=-2,\quad M_2=-2.
\]
Nhận lại $S(x)=2x-x^2$. Hai loại biên khác tên có thể cho cùng spline nếu cùng mô tả đúng hàm này; không mặc định chúng luôn tương đương.

\subsection{Biên hỗn hợp: mỗi đầu được chọn riêng}
\label{sec:spline-mixed}
Tại mỗi đầu, chọn cấp 1 hoặc cấp 2; ghép hàng tương ứng từ hai mục trên.
Nếu trái cho cấp 1, phải cho cấp 2:
\[
2h_0M_0+h_0M_1=6(\delta_0-s_L),\qquad M_n=R.
\]
Nếu trái cho cấp 2, phải cho cấp 1:
\[
M_0=L,\qquad h_{n-1}M_{n-1}+2h_{n-1}M_n=6(s_R-\delta_{n-1}).
\]
Trường hợp $(1,1)$ chính là kẹp; $(2,2)$ là cho hai độ cong.
Ví dụ $s_L=2,R=-2$ trên bảng chung:
$2M_0+M_1=-6$, $M_0+4M_1+M_2=-12$, $M_2=-2$.
Suy ra $M=(-2,-2,-2)$, $S(0.5)=3/4$.
Đổi sang trái cấp 2 giá trị $-2$, phải cấp 1 giá trị $-2$ cũng nhận đúng spline này.
Phải nhập \emph{cấp} và \emph{giá trị} riêng; số $-2$ tự nó không cho biết là độ dốc hay độ cong.

\subsection{Lưới không đều vẫn làm cùng quy trình}
\label{sec:spline-nonuniform}
Cho $(0,0),(1,1),(3,0)$, biên tự nhiên.
Tính $h=(1,2)$, $\delta=(1,-1/2)$. Hàng trong:
\[
1M_0+2(1+2)M_1+2M_2=6(-1/2-1)=-9.
\]
Do đó $M=(0,-3/2,0)$.
Mảnh thứ hai là $S_1(1+u)=1+\frac12u-\frac34u^2+\frac18u^3$, $0\leq u\leq2$.
Ở $x=2$, $u=1$ nên $S(2)=1+1/2-3/4+1/8=7/8$.
Không thay hai bước $1,2$ bằng bước trung bình để dựng cubic.

\section{Cubic tuần hoàn: hai hàng nối đầu và cuối}
\label{sec:spline-periodic}
Đầu vào phải có $y_0=y_n$ chính xác theo dữ liệu đã đọc. Ngoài ghép trong, yêu cầu
$S'(x_0)=S'(x_n)$ và $S''(x_0)=S''(x_n)$.
Giữ các hàng nút trong, thay hai hàng biên bằng
\begin{align*}
M_0-M_n&=0,\\
2h_0M_0+h_0M_1+h_{n-1}M_{n-1}+2h_{n-1}M_n
&=6(\delta_0-\delta_{n-1}).
\end{align*}
Các hệ số ở hai đầu có thể nằm ngoài ba đường chéo thông thường.
Khi ít mốc, một ô có thể nhận hai đóng góp: phải \emph{cộng dồn}, không ghi đè.

Ví dụ $(0,0),(1,1),(2,0)$ có hệ
\[
M_0-M_2=0,\quad M_0+4M_1+M_2=-12,\quad
2M_0+2M_1+2M_2=12.
\]
Từ $M_2=M_0$, hai hàng còn lại cho $M_0+2M_1=-6$ và $2M_0+M_1=6$.
Suy ra $M=(6,-6,6)$.
Các mảnh là
\[
S_0(x)=3x^2-2x^3,\qquad S_1(1+u)=1-3u^2+2u^3.
\]
Tại hai đầu, giá trị đều 0, độ dốc đều 0 và độ cong đều 6.
Tool yêu cầu ít nhất 3 mốc cho quy ước tuần hoàn này; không tự tạo trường hợp hai mốc.

\section{Cubic not-a-knot: nối cả đạo hàm cấp ba ở hai nút gần biên}
\label{sec:spline-notaknot}
Với ít nhất 4 mốc, yêu cầu mảnh 0 và 1 có cùng đạo hàm cấp 3; tương tự hai mảnh cuối.
Vì $S_i'''=(M_{i+1}-M_i)/h_i$, hai hàng là
\begin{align*}
-\frac{M_0}{h_0}+\left(\frac1{h_0}+\frac1{h_1}\right)M_1-\frac{M_2}{h_1}&=0,\\
-\frac{M_{n-2}}{h_{n-2}}+\left(\frac1{h_{n-2}}+\frac1{h_{n-1}}\right)M_{n-1}-\frac{M_n}{h_{n-1}}&=0.
\end{align*}
Giữ các hàng nút trong và giải hệ. Đây không phải tự nhiên và cũng không đòi $y_0=y_n$.

Ví dụ $x=(0,1,2,3)$, $y=(0,1,8,27)$:
\[
-M_0+2M_1-M_2=0,\quad M_0+4M_1+M_2=36,
\]
\[
M_1+4M_2+M_3=72,\quad -M_1+2M_2-M_3=0.
\]
Nghiệm $M=(0,6,12,18)$ cho hệ số mảnh lần lượt
$(0,0,0,1)$, $(1,3,3,1)$, $(8,12,6,1)$ theo bậc $u$ tăng.
Chúng là các khai triển của $x^3$ quanh $0,1,2$; $S(1.5)=27/8$.
Hai mảnh gần mỗi biên thực chất thuộc cùng một đa thức cubic.

\section{Spline bậc p: vì sao cần p−1 điều kiện bổ sung?}
\label{sec:spline-degree-p}
Có $n$ đoạn, mỗi đoạn $p+1$ hệ số, tổng $n(p+1)$ ẩn.
Nội suy hai đầu mỗi đoạn cho $2n$ phương trình. Ghép đạo hàm cấp $1,\ldots,p-1$
tại $n-1$ nút trong cho $(n-1)(p-1)$ phương trình.
Số bậc tự do còn lại là
\[
n(p+1)-2n-(n-1)(p-1)=p-1.
\]
Vì vậy cần thêm $p-1$ điều kiện \emph{độc lập và tương thích} để xác định duy nhất.
Đếm đủ số hàng không chứng minh độc lập. Một điều kiện về giá trị ở đầu mút có thể chỉ lặp lại dữ liệu nội suy và không giảm bậc tự do.
Hai mốc cũng có thể đủ nếu các điều kiện thêm xác định một mảnh duy nhất; không bắt số mốc bằng $p+1$ như nội suy toàn cục.

\subsection{Truy hồi khi biết $p-1$ đạo hàm ở đầu trái}
\label{sec:spline-left-general}
Trên đoạn $i$, biết $v_r=S^{(r)}(x_i)$, $r=1,\ldots,p-1$.
Viết $S_i(x_i+u)=\sum_{r=0}^p a_{i,r}u^r$.
\begin{enumerate}
\item Gán $a_{i,0}=y_i$ và $a_{i,r}=v_r/r!$ cho $1\leq r<p$.
\item Tính giá trị đầu phải từ các hệ số đã biết:
$K_i=\sum_{r=0}^{p-1}a_{i,r}h_i^r$.
\item Dùng nội suy đầu phải tìm hệ số cuối:
\[
a_{i,p}=\frac{y_{i+1}-K_i}{h_i^p}.
\]
\item Tính \emph{toàn bộ} đạo hàm đầu phải vào mảng mới:
\[
w_r=\sum_{j=r}^{p}a_{i,j}\frac{j!}{(j-r)!}h_i^{j-r},\qquad r=1,\ldots,p-1.
\]
\item Lưu mảnh, rồi gán đồng thời $v\leftarrow w$ để bắt đầu đoạn $i+1$.
\end{enumerate}
Không cập nhật từng phần tử $v_r$ khi những đạo hàm khác vẫn cần bộ cũ.
Đây là truyền hữu hạn qua các đoạn, không phải một phép lặp đến epsilon.

\subsection{Bậc 4, ví dụ đủ hai lần truyền}
\label{sec:spline-quartic-left}
Nhập bảng chung và $S'(0)=1,S''(0)=S'''(0)=0$.
\begin{center}\small
\begin{tabular}{rp{.25\linewidth}p{.25\linewidth}p{.25\linewidth}}
\toprule
Đoạn & Đạo hàm đầu vào & Hệ số tìm được & Đạo hàm chuyển sang phải\\ \midrule
0 & $(1,0,0)$ & $(0,1,0,0,0)$ & $(1,0,0)$\\
1 & $(1,0,0)$ & $(1,1,0,0,-2)$ & $(-7,-24,-48)$\\ \bottomrule
\end{tabular}\end{center}
Đoạn 0: $K_0=0+1=1$, $a_{0,4}=(1-1)/1^4=0$.
Đoạn 1: $K_1=1+1=2$, $a_{1,4}=(0-2)/1^4=-2$.
Do đó
\[
S_0(x)=x,\qquad S_1(1+u)=1+u-2u^4.
\]
Các đạo hàm chuyển sau đoạn 1 lần lượt là
$1-8=-7$, $-24$, $-48$. Tại nút 1, ba đạo hàm hai phía là $(1,0,0)$.
Tại $1.5$: $u=0.5$ nên $S=1+1/2-2/16=11/8$.
Không nhầm bộ $(-7,-24,-48)$ ở đầu bảng cuối với điều kiện đầu trái ban đầu.

\section{Biên tổng quát: đầu phải, hai đầu hoặc tổ hợp tuyến tính}
\label{sec:spline-general-boundary}
Trong tool chọn \textbf{Biên tổng quát}, nhập bậc $p$.
Một điều kiện đơn có dạng \texttt{["trái", cấp, giá trị]} hoặc
\texttt{["phải", cấp, giá trị]}. Một tổ hợp có dạng
\texttt{[[[phía,cấp,hệ số],...], giá trị]}.
Các phía đều là \emph{hai đầu toàn bộ bảng}, không phải một nút trong tùy ý.

\subsection{Lập hệ từ định nghĩa}
\label{sec:spline-general-system}
\begin{enumerate}
\item Đặt toàn bộ $a_{i,0},\ldots,a_{i,p}$ làm ẩn, theo từng đoạn và bậc tăng.
\item Với đoạn $i$, thêm $a_{i,0}=y_i$ và $\sum_{r=0}^p a_{i,r}h_i^r=y_{i+1}$.
\item Với mỗi nút trong và mỗi $k=1,\ldots,p-1$, thêm
\[
\sum_{r=k}^p a_{i,r}\frac{r!}{(r-k)!}h_i^{r-k}=k!a_{i+1,k}.
\]
\item Chuyển từng điều kiện biên thành một hàng tuyến tính trên cùng dãy ẩn.
\item Khử Gauss có chọn trụ. Kiểm hạng và tính tương thích; nếu duy nhất thì thế ngược.
\item Dựng các mảnh rồi kiểm lại \emph{mọi} điều kiện đã nhập, kể cả điều kiện dư tương thích.
\end{enumerate}
Nếu hạng hệ số nhỏ hơn số ẩn thì chưa duy nhất. Nếu có hàng $0=\beta\ne0$ sau khử thì mâu thuẫn.
Một điều kiện dư đúng có thể được giữ; một điều kiện dư sai không được bỏ qua.

\subsection{Bậc 4 chỉ cho đầu phải: một ví dụ hoàn chỉnh}
\label{sec:spline-quartic-right}
Cho $x=(0,1,2)$, $y=(0,1,16)$ và
$S'(2)=32,S''(2)=48,S'''(2)=48$.
Gọi mảnh 0 là $A(x)=\sum_{r=0}^4a_rx^r$, mảnh 1 là $B(1+u)=\sum_{r=0}^4b_ru^r$.
Mười phương trình theo đúng thứ tự lập hệ là
\begin{align*}
a_0&=0,&a_0+a_1+a_2+a_3+a_4&=1,\\
b_0&=1,&b_0+b_1+b_2+b_3+b_4&=16,\\
a_1+2a_2+3a_3+4a_4-b_1&=0,\\
2a_2+6a_3+12a_4-2b_2&=0,\\
6a_3+24a_4-6b_3&=0,\\
b_1+2b_2+3b_3+4b_4&=32,\\
2b_2+6b_3+12b_4&=48,\\
6b_3+24b_4&=48.
\end{align*}
Vết chọn trụ thực tế của tool cho thứ tự ẩn $(a_0,\ldots,a_4,b_0,\ldots,b_4)$:
\begin{center}\small
\begin{tabular}{rrrrl}
\toprule
Bước & Cột & Hàng được chọn & Pivot & Hàng dưới bị khử: hệ số nhân\\ \midrule
0&0&0&1&$1:1$\\
1&1&1&1&$4:1$\\
2&2&5&2&$4:1/2$\\
3&3&6&6&$4:-1/6$\\
4&4&4&1&không có\\
5&5&5&1&$6:1$\\
6&6&6&1&$7:1$\\
7&7&8&2&$8:1/2$\\
8&8&9&6&$9:-1/6$\\
9&9&9&1&không có\\ \bottomrule
\end{tabular}\end{center}
Chỉ số hàng/cột bắt đầu từ 0; hàng được chọn là vị trí tại \emph{bước đang xét}, sau các lần đổi hàng trước.
Khử dùng $R_i\leftarrow R_i-\lambda R_k$, đồng thời trên vế phải.
Thế ngược cho
\[
a=(0,0,0,0,1),\qquad b=(1,4,6,4,1).
\]
Hai mảnh chính là $x^4$ ở hai gốc khác nhau; $S(1.5)=81/16=5.0625$.
JSON nhập vào ô điều kiện là
\begin{verbatim}
[["phải",1,32],["phải",2,48],["phải",3,48]]
\end{verbatim}

\subsection{Biên chia ở hai đầu hoặc tổ hợp}
\label{sec:spline-quartic-mixed}
Trên cùng bảng, thay ba biên phải bằng
$S'(0)=0,S'(2)=32,S''(2)=48$. Hàng thêm là
\[
a_1=0,\quad b_1+2b_2+3b_3+4b_4=32,\quad
2b_2+6b_3+12b_4=48.
\]
Giữ bảy hàng nội suy/ghép; hệ vẫn cho hai dãy hệ số trên. JSON:
\begin{verbatim}
[["trái",1,0],["phải",1,32],["phải",2,48]]
\end{verbatim}
Hoặc dùng $S'(0)+S'(2)=32$, $S''(2)=48$, $S'''(2)=48$:
\begin{verbatim}
[[[["trái",1,1],["phải",1,1]],32],
 ["phải",2,48],["phải",3,48]]
\end{verbatim}
Hàng đầu khi ấy là $a_1+b_1+2b_2+3b_3+4b_4=32$.
Đây cũng cho $S=x^4$, đã được chạy đối chiếu bằng tool.
Sự trùng kết quả trong ví dụ này do các biên đều đúng với $x^4$; không được tự đổi các bộ biên trong một đề khác.

\section{Đánh giá, đạo hàm ở nút và tích phân qua nhiều mảnh}
\label{sec:spline-evaluation}
\begin{enumerate}
\item \textbf{Tìm đoạn:} với $z$ trong miền, lấy đoạn chứa $z$. Tại đầu phải cuối cùng, dùng mảnh cuối.
\item \textbf{Thế biến đúng:} $u=z-x_i$. Giá trị là $\sum a_{i,r}u^r$.
\item \textbf{Đạo hàm:} $S_i^{(s)}(z)=\sum_{r=s}^p a_{i,r}r!u^{r-s}/(r-s)!$.
Vì $du/dx=1$, không chia thêm $h_i^s$ như khi dùng biến chuẩn hóa $t$ của nội suy đều.
\item \textbf{Ở nút trong:} kiểm hai phía cho mọi cấp đến cấp đang hỏi. Bậc $p$ chỉ được đảm bảo trơn đến $p-1$.
Ngay cả khi hai mảnh đều có đạo hàm cấp $p+1$ bằng 0, spline toàn cục vẫn có thể không có đạo hàm cấp $p$ tại nút, nên không thể suy ra đạo hàm cấp $p+1$ tại đó.
\end{enumerate}
Ví dụ cubic tự nhiên ba mốc có $S_0'''=-3,S_1'''=3$ tại nút 1.
Đạo hàm cấp 3 tại nút không tồn tại; hỏi cấp 4 ở đúng nút cũng không được trả 0 chỉ từ hai mảnh.

Tích phân từ $a$ đến $b$: chèn mọi nút nằm giữa hai cận, rồi tính trên từng đoạn con
\[
\int_{x_i+u}^{x_i+v}S_i(x)\,dx
=\sum_{r=0}^p a_{i,r}\frac{v^{r+1}-u^{r+1}}{r+1}.
\]
Nếu đảo cận thì đổi dấu. Ví dụ cubic tự nhiên ba mốc:
$\int_0^1(3x/2-x^3/2)dx=5/8$ và $\int_1^2S_1dx=5/8$,
nên $\int_0^2S=5/4$.

Ngoài miền, mặc định không đánh giá. Khi được bật ngoại suy tường minh, tool kéo dài mảnh biên và ghi cảnh báo; không kẹp điểm hỏi về đầu mút.
Đây là giá trị ngoại suy của mảnh, chưa có bảo đảm về hàm thật.

\section{Tính duy nhất, sai số và những giới hạn thực tế}
\label{sec:spline-limits}
\begin{itemize}
\item Với các $h_i>0$, hệ cubic tự nhiên/kẹp/cấp hai/hỗn hợp cấp 1--2 chuẩn được giải trực tiếp bằng hệ moment. Cần giữ đúng các hàng biên và kiểm pivot/phần dư.
\item Spline tổng quát cần kiểm hạng; không có định lý rằng \emph{mọi} ba số biên tùy ý của spline bậc 4 đều xác định một nghiệm duy nhất.
\item Spline có thể vượt ngoài khoảng các tung độ hoặc dao động; spline tự nhiên không tự bảo toàn đơn điệu. Sai lệch đúng ở các mốc không đo sai số giữa mốc.
\item Cận sai số cần giả thiết đủ trơn và đúng kiểu biên. Không áp một hằng số sai số của spline kẹp cho mọi spline tự nhiên hoặc cho mọi bậc.
\item Các mảnh và vết giải trong ví dụ nhỏ của tool dùng hữu tỷ của số thập phân đã đọc; sau đó chuyển sang biểu diễn số máy để vẽ/đánh giá. Báo cáo còn kiểm sự chuyển đổi này. Đây không làm dữ liệu đo ban đầu trở thành chính xác tuyệt đối.
\item Cubic thường dùng Thomas; chọn Gauss thì có vết khử Gauss thật. Tuần hoàn/not-a-knot không dùng Thomas thường. Với hơn 80 mốc và chế độ tự chọn cho hai loại này, tool gọi bộ giải thư viện, giữ hệ thưa và kiểm điều kiện; không tạo một vết khử Gauss/Thomas giả để trang trí.
\item Menu chỉ có lựa chọn nhanh \emph{bậc 4 đầu trái}. Biên phải và biên trộn đi qua \emph{Biên tổng quát}; không phải giới hạn toán học của spline.
\item Trong phiên bản được đối chiếu, nhánh Biên tổng quát lập hệ đặc cho mọi hệ số và dùng Gauss. Lựa chọn Thomas ở nhánh này đang bị bỏ qua; nên chọn Gauss/tự chọn và đọc bộ giải thực tế trong báo cáo. Đây là giới hạn triển khai, không phải một phép biến đổi Thomas hợp lệ.
\item Lõi \texttt{solve\_spline} nhận bậc nguyên $p\geq1$; hai API phụ cũ \texttt{spline\_degree}/\texttt{spline\_boundary} giới hạn bậc 1--4 và có quy tắc số biên riêng. Không đồng nhất hai đường gọi.
\item Biên tổng quát hiện là các điều kiện tuyến tính ở hai đầu; nó không bao gồm mọi họ B-spline, spline làm trơn, knot bội hay mức trơn tùy ý.
\end{itemize}

\paragraph{Khi chép bài thi.}
Ghi bảng mốc, $h_i,\delta_i$; loại biên và dữ kiện biên; ẩn và hệ/truy hồi;
các bước giải; bốn hoặc $p+1$ hệ số trên mỗi đoạn; miền mảnh và biến $u$;
phép thế tại điểm hỏi; kiểm nội suy/ghép/biên; cuối cùng ghi loại sai số biết được.
Không gọi khử hệ hoặc truyền qua hữu hạn các đoạn là ``lặp đến hội tụ''.

\paragraph{Nguồn trong repo và dữ liệu kiểm chứng.}
Nội dung đối chiếu \texttt{pps/spline\_worked.py}, các trường Spline trong
\texttt{pps/gui.py}, tài liệu spline 54 trang ở \texttt{kienthuc/real}, và các bài kiểm tra spline.
Các ví dụ số trong chương là ví dụ tự dựng/bài mẫu đã có trong repo, không gán là dữ liệu của một đề thi khác.
File \texttt{verification/spline\_ls\_examples.py} chạy lại các ví dụ bằng bộ tính thật;
JSON cùng tên giữ hệ số và vết giải đầy đủ để đối chiếu.

\chapter[Bình phương tối thiểu: từng mô hình]{Tool Bình phương tối thiểu: input, bảng tính và output từng mô hình}
\label{ch:ls}
\examplebreak
\section{Chọn mô hình rồi mới chọn phép biến đổi}
Cho $N$ cặp $(x_0,y_0),\ldots,(x_{N-1},y_{N-1})$. Các $x$ có thể lặp: đó là các quan sát riêng. Tool tìm hàm khớp $g$, không ép hàm đi qua mọi điểm.
\subsection*{Input: các ô đúng tên trên giao diện}
Bảng này được lấy từ biểu mẫu của tool. Giá trị mặc định là ví dụ nhập, cần thay bằng dữ liệu đề. Hai dãy $x,y$ là bảng dữ liệu riêng bên dưới nếu chức năng dùng bảng.
\subsubsection*{Ô chung của các lựa chọn trong nhóm}
\begin{center}\small\begin{tabular}{@{}p{.32\linewidth}p{.17\linewidth}p{.43\linewidth}@{}}\toprule
Ô nhập & Khóa JSON & Cách nhập / mặc định\\\midrule
Trọng số w\_i (để trống: tất cả bằng 1) & \texttt{weights} & Để trống
\par Không âm,\allowbreak{} ít nhất một trọng số dương;\allowbreak{} giữ cùng thứ tự dữ liệu.\\[3pt]
Điểm dự đoán x (không bắt buộc) & \texttt{target} & Để trống\\[3pt]
Mức trình bày & \texttt{detail\_level} & \texttt{Tổng hợp}\\[3pt]
\bottomrule\end{tabular}\end{center}
\subsubsection*{Ô riêng: Tuyến tính theo cơ sở /\allowbreak{} Log của tổ hợp cơ sở}
\begin{center}\small\begin{tabular}{@{}p{.32\linewidth}p{.17\linewidth}p{.43\linewidth}@{}}\toprule
Ô nhập & Khóa JSON & Cách nhập / mặc định\\\midrule
Các hàm cơ sở (ngăn bằng dấu ;\allowbreak{}) & \texttt{basis} & \texttt{1;\allowbreak{} x}
\par Nhập đầy đủ các cơ sở,\allowbreak{} kể cả 1 nếu mô hình có hằng;\allowbreak{} lựa chọn này không tự thêm cột 1.\\[3pt]
Cách giải bình phương tối thiểu & \texttt{solver} & \texttt{QR}\\[3pt]
Quy tắc khi thiếu hạng & \texttt{rank\_policy} & \texttt{Nghiệm duy nhất}\\[3pt]
\bottomrule\end{tabular}\end{center}
\subsubsection*{Ô riêng: Lũy thừa /\allowbreak{} 1/\allowbreak{}(a+b*x) /\allowbreak{} x/\allowbreak{}(a+b*x) /\allowbreak{} a+b*ln(x) /\allowbreak{} A*b\textasciicircum{}x}
\begin{center}\small\begin{tabular}{@{}p{.32\linewidth}p{.17\linewidth}p{.43\linewidth}@{}}\toprule
Ô nhập & Khóa JSON & Cách nhập / mặc định\\\midrule
Cách giải bình phương tối thiểu & \texttt{solver} & \texttt{QR}\\[3pt]
Quy tắc khi thiếu hạng & \texttt{rank\_policy} & \texttt{Nghiệm duy nhất}\\[3pt]
\bottomrule\end{tabular}\end{center}
\subsubsection*{Ô riêng: Hàm mũ}
\begin{center}\small\begin{tabular}{@{}p{.32\linewidth}p{.17\linewidth}p{.43\linewidth}@{}}\toprule
Ô nhập & Khóa JSON & Cách nhập / mặc định\\\midrule
Các hàm cơ sở (ngăn bằng dấu ;\allowbreak{}) & \texttt{basis} & \texttt{x}
\par Chỉ nhập các cơ sở biến thiên trong số mũ;\allowbreak{} tool tự thêm cột hằng cho biên độ. Không nhập thêm 1.\\[3pt]
Cách giải bình phương tối thiểu & \texttt{solver} & \texttt{QR}\\[3pt]
Quy tắc khi thiếu hạng & \texttt{rank\_policy} & \texttt{Nghiệm duy nhất}\\[3pt]
\bottomrule\end{tabular}\end{center}
\subsubsection*{Ô riêng: Lũy thừa — SSE thang y}
\begin{center}\small\begin{tabular}{@{}p{.32\linewidth}p{.17\linewidth}p{.43\linewidth}@{}}\toprule
Ô nhập & Khóa JSON & Cách nhập / mặc định\\\midrule
Tham số đầu A b1 b2... (để trống: từ log) & \texttt{initial} & Để trống
\par Lời giải ghi tham số và SSE từng vòng.\\[3pt]
Ngưỡng dừng phi tuyến & \texttt{tol} & \texttt{1e-10}\\[3pt]
Số vòng tối đa & \texttt{max\_iter} & \texttt{100}\\[3pt]
Bộ giải phi tuyến & \texttt{solver} & \texttt{Levenberg–Marquardt}\\[3pt]
\bottomrule\end{tabular}\end{center}
\subsubsection*{Ô riêng: Hàm mũ — SSE thang y}
\begin{center}\small\begin{tabular}{@{}p{.32\linewidth}p{.17\linewidth}p{.43\linewidth}@{}}\toprule
Ô nhập & Khóa JSON & Cách nhập / mặc định\\\midrule
Các hàm cơ sở (ngăn bằng dấu ;\allowbreak{}) & \texttt{basis} & \texttt{x}
\par Chỉ nhập các cơ sở biến thiên trong số mũ;\allowbreak{} tool tự thêm cột hằng cho biên độ. Không nhập thêm 1.\\[3pt]
Tham số đầu A b1 b2... (để trống: từ log) & \texttt{initial} & Để trống
\par Lời giải ghi tham số và SSE từng vòng.\\[3pt]
Ngưỡng dừng phi tuyến & \texttt{tol} & \texttt{1e-10}\\[3pt]
Số vòng tối đa & \texttt{max\_iter} & \texttt{100}\\[3pt]
Bộ giải phi tuyến & \texttt{solver} & \texttt{Levenberg–Marquardt}\\[3pt]
\bottomrule\end{tabular}\end{center}

\subsubsection*{Các bước bấm.} Chọn nhóm Bình phương tối thiểu $\to$ chọn đúng mô hình bên dưới $\to$ nhập $x,y$ $\to$ nhập cơ sở nếu có ô đó $\to$ trọng số nếu đề cho $\to$ điểm dự đoán nếu cần $\to$ Tính kết quả. Mục \emph{Chi tiết} mở các bước của bộ giải thực tế. Chương này tập trung việc lập bài toán và đọc kết quả; không giảng lại Gauss, QR hay SVD.

\examplebreak
\section{Mẫu chung cho bài hai tham số sau biến đổi}
\label{sec:ls-general-table}
Mỗi mô hình sẽ quy định cụ thể $U_i,V_i$. Khi có dạng $V\approx c_0+c_1U$, điền bảng dưới đây. Nếu không nhập trọng số thì $w_i=1$ cho mọi dòng.
\begin{center}\small\begin{tabular}{@{}crrrrr@{}}\toprule
Dòng&$w_i$&$w_iU_i$&$w_iV_i$&$w_iU_i^2$&$w_iU_iV_i$\\\midrule
0&$w_0$&$w_0U_0$&$w_0V_0$&$w_0U_0^2$&$w_0U_0V_0$\\
$i$&$w_i$&$w_iU_i$&$w_iV_i$&$w_iU_i^2$&$w_iU_iV_i$\\
$N-1$&$w_{N-1}$&$w_{N-1}U_{N-1}$&$w_{N-1}V_{N-1}$&$w_{N-1}U_{N-1}^2$&$w_{N-1}U_{N-1}V_{N-1}$\\\midrule
Tổng&$S_1$&$S_U$&$S_V$&$S_{UU}$&$S_{UV}$\\\bottomrule
\end{tabular}\end{center}
Hàng $i$ là khuôn để điền mọi hàng giữa, không phải một quan sát bổ sung. Thứ tự: biến đổi $U,V$ $\to$ nhân trọng số và tính các tích $\to$ cộng từng cột:
\begin{align*}
S_1&=\sum_{i=0}^{N-1}w_i=w_0+w_1+\cdots+w_{N-1},\\
S_U&=\sum_{i=0}^{N-1}w_iU_i=w_0U_0+w_1U_1+\cdots+w_{N-1}U_{N-1},\\
S_V&=\sum_{i=0}^{N-1}w_iV_i=w_0V_0+w_1V_1+\cdots+w_{N-1}V_{N-1},\\
S_{UU}&=\sum_{i=0}^{N-1}w_iU_i^2=w_0U_0^2+w_1U_1^2+\cdots+w_{N-1}U_{N-1}^2,\\
S_{UV}&=\sum_{i=0}^{N-1}w_iU_iV_i=w_0U_0V_0+w_1U_1V_1+\cdots+w_{N-1}U_{N-1}V_{N-1}.
\end{align*}
Từ hàng tổng lập hệ
\[
\boxed{\begin{cases}S_1c_0+S_Uc_1=S_V,\\S_Uc_0+S_{UU}c_1=S_{UV}.\end{cases}}
\]
Không trọng số thì $S_1=N$. Nếu $D=S_1S_{UU}-S_U^2\ne0$, kiểm nhanh bằng
\[
c_0=(S_VS_{UU}-S_US_{UV})/D,\quad c_1=(S_1S_{UV}-S_US_V)/D.
\]
Nếu $D=0$ thì tham số chưa duy nhất theo dữ liệu hiệu lực; không gọi là vô nghiệm bình phương tối thiểu. Dòng trọng số 0 vẫn cần dữ liệu và phép biến đổi hợp lệ.

\examplebreak
\section{Tuyến tính theo cơ sở: tự chỉ định mọi cột}
\label{sec:ls-linear-tool}
\subsubsection*{Input.} Chọn \textbf{Tuyến tính theo cơ sở}. Ô cơ sở nhập đầy đủ, ngăn bằng \texttt{;}: đường thẳng \texttt{1;x}, bậc hai \texttt{1;x;x\^{}2}, đi qua gốc \texttt{x}, hoặc \texttt{1;x;cos(x)}. Ở lựa chọn này không tự thêm cột 1.
\subsubsection*{Bước 1 -- mỗi dòng tạo một hàng.}
\[
g(x)=\sum_{j=0}^{p-1}c_j\phi_j(x)=c_0\phi_0(x)+c_1\phi_1(x)+\cdots+c_{p-1}\phi_{p-1}(x).
\]
\begin{center}\small\begin{tabular}{cllll}\toprule
Dòng&$y_i$&$\phi_0(x_i)$&$\phi_1(x_i)$&$\ldots,\phi_{p-1}(x_i)$\\\midrule
0&$y_0$&$\phi_0(x_0)$&$\phi_1(x_0)$&$\ldots,\phi_{p-1}(x_0)$\\
$i$&$y_i$&$\phi_0(x_i)$&$\phi_1(x_i)$&$\ldots,\phi_{p-1}(x_i)$\\
$N-1$&$y_{N-1}$&$\phi_0(x_{N-1})$&$\phi_1(x_{N-1})$&$\ldots,\phi_{p-1}(x_{N-1})$\\\bottomrule
\end{tabular}\end{center}
Các cột cơ sở tạo $\Phi$; đích $z=(y_0,\ldots,y_{N-1})^T$. Viết $\Phi c\approx z$ vì đây là bài khớp.
\subsubsection*{Bước 2 -- tính các tích và tổng.}
\begin{align*}
G_{jk}&=\sum_iw_i\phi_j(x_i)\phi_k(x_i)\\
 &=w_0\phi_j(x_0)\phi_k(x_0)+\cdots+w_{N-1}\phi_j(x_{N-1})\phi_k(x_{N-1}),\\
b_j&=\sum_iw_i\phi_j(x_i)y_i=w_0\phi_j(x_0)y_0+\cdots+w_{N-1}\phi_j(x_{N-1})y_{N-1}.
\end{align*}
\subsubsection*{Bước 3 -- giải bài khớp.} Hệ chuẩn $Gc=b$, trong đó hàng $j$ là
\[
G_{j0}c_0+G_{j1}c_1+\cdots+G_{j,p-1}c_{p-1}=b_j.
\]
Tool dùng đúng bộ giải đã chọn; không giả sử mọi lựa chọn đều khử hệ chuẩn. Đây là quá trình hữu hạn, không có ngưỡng hội tụ tham số cho bài tuyến tính.
\subsubsection*{Bước 4 -- output.} Vector $(c_0,\ldots,c_{p-1})$ theo thứ tự cơ sở, hàm $g$, dự đoán trên dữ liệu, phần dư, sai số và $g(x_*)$ nếu nhập điểm dự đoán.

\subsection{Đường thẳng và đa thức: viết rõ hệ từ những tổng nào}
Với \texttt{1;x}, dùng bảng hai tham số với $U_i=x_i,V_i=y_i$. Ví dụ $x=(0,1,2,3),y=(1,2,2,4)$:
\[
S_1=4,\quad S_U=0+1+2+3=6,\quad S_V=1+2+2+4=9,
\]
\[
S_{UU}=0^2+1^2+2^2+3^2=14,\quad S_{UV}=0(1)+1(2)+2(2)+3(4)=18.
\]
Hệ $4c_0+6c_1=9$, $6c_0+14c_1=18$ cho $g(x)=0.9+0.9x$.

Với \texttt{1;x;x\^{}2}, bảng cần thêm các tích $w_ix_i^3,w_ix_i^4,w_ix_i^2y_i$. Đặt
\[
M_r=w_0x_0^r+\cdots+w_{N-1}x_{N-1}^r,\quad
T_r=w_0x_0^ry_0+\cdots+w_{N-1}x_{N-1}^ry_{N-1}.
\]
Riêng $M_0=w_0+\cdots+w_{N-1}$, $T_0=w_0y_0+\cdots+w_{N-1}y_{N-1}$, kể cả khi có $x_i=0$. Hệ là
\[
\begin{bmatrix}M_0&M_1&M_2\\M_1&M_2&M_3\\M_2&M_3&M_4\end{bmatrix}
\begin{bmatrix}c_0\\c_1\\c_2\end{bmatrix}=\begin{bmatrix}T_0\\T_1\\T_2\end{bmatrix}.
\]
Khôi phục $g=c_0+c_1x+c_2x^2$. Bậc cao hơn tạo hệ tương tự với đủ các lũy thừa đến $2p-2$.


\examplebreak
\section{Lũy thừa: Ax mũ b}

\label{sec:ls-power}

\subsubsection*{Input trong tool.} Chọn \textbf{Lũy thừa}; nhập $x,y$, trọng số nếu có và điểm dự đoán nếu cần. Không có ô cơ sở. Cần mọi $x_i>0$, $y_i\ne0$ và cùng dấu; điểm dự đoán cũng cần $x_*>0$.

\subsubsection*{Bước 1 -- biến đổi từng dòng.}\[y\approx Ax^b\quad\Longrightarrow\quad\ln|y|\approx C+b\ln x,\quad C=\ln|A|.\]

\begin{center}\small\begin{tabular}{cllll}\toprule

Dòng&$U_i$&$V_i$&$w_iU_i^2$&$w_iU_iV_i$\\\midrule

$0$&$\ln x_0$&$\ln|y_0|$&$w_{0}(\ln x_0)^2$&$w_{0}(\ln x_0)(\ln|y_0|)$\\

$i$&$\ln x_{i}$&$\ln|y_{i}|$&$w_{i}(\ln x_{i})^2$&$w_{i}(\ln x_{i})(\ln|y_{i}|)$\\

$N-1$&$\ln x_{N-1}$&$\ln|y_{N-1}|$&$w_{N-1}(\ln x_{N-1})^2$&$w_{N-1}(\ln x_{N-1})(\ln|y_{N-1}|)$\\

\bottomrule\end{tabular}\end{center}

\subsubsection*{Bước 2 -- cộng cột, lập hệ.} Các tổng viết riêng cho loại này là

\begin{align*}S_U&=w_0\ln x_0+\cdots+w_{N-1}\ln x_{N-1},\\S_V&=w_0\ln|y_0|+\cdots+w_{N-1}\ln|y_{N-1}|,\\S_{UV}&=w_0\ln x_0\ln|y_0|+\cdots+w_{N-1}\ln x_{N-1}\ln|y_{N-1}|.\end{align*}

$S_1=w_0+\cdots+w_{N-1}$; $S_{UU}$ là tổng cột bình phương ở bảng trên, tức $w_0U_0^2+\cdots+w_{N-1}U_{N-1}^2$. Hệ cần giải:

\[\boxed{\begin{cases}S_1C+S_Ub=S_V,\\S_UC+S_{UU}b=S_{UV}.\end{cases}}\]

\subsubsection*{Bước 3 -- output và cách đọc.} Đặt $s=\operatorname{sign}(y_0)$; trả $A=se^C$, giữ $b$ rồi viết $g=Ax^b$. Hệ số $(C,b)$ của bài tuyến tính chưa phải $(A,b)$ của đề.

Tool cực tiểu hóa tổng bình phương có trọng số trên thang $\ln|y|$; dữ liệu $y$ âm cùng dấu vẫn khôi phục biên độ $A$ âm. Sai số trên thang $y$ còn được báo thêm sau khôi phục. Đây là phép giải hữu hạn; không có bảng vòng hội tụ tham số ở nhánh này.

\examplebreak
\section{Hàm mũ: Ae mũ bx}

\label{sec:ls-exp-log}

\subsubsection*{Input trong tool.} Chọn \textbf{Hàm mũ}; nhập $x,y$, trọng số nếu có và điểm dự đoán nếu cần. Ô cơ sở nhập \texttt{x} cho $Ae^{bx}$. Tool tự thêm cột 1 để tìm $\ln|A|$; không nhập thêm 1. Cần mọi $y_i$ cùng dấu khác 0; $x_i$ hữu hạn.

\subsubsection*{Bước 1 -- biến đổi từng dòng.}\[y\approx Ae^{bx}\quad\Longrightarrow\quad\ln|y|\approx C+bx,\quad C=\ln|A|.\]

\begin{center}\small\begin{tabular}{cllll}\toprule

Dòng&$U_i$&$V_i$&$w_iU_i^2$&$w_iU_iV_i$\\\midrule

$0$&$x_0$&$\ln|y_0|$&$w_{0}(x_0)^2$&$w_{0}(x_0)(\ln|y_0|)$\\

$i$&$x_{i}$&$\ln|y_{i}|$&$w_{i}(x_{i})^2$&$w_{i}(x_{i})(\ln|y_{i}|)$\\

$N-1$&$x_{N-1}$&$\ln|y_{N-1}|$&$w_{N-1}(x_{N-1})^2$&$w_{N-1}(x_{N-1})(\ln|y_{N-1}|)$\\

\bottomrule\end{tabular}\end{center}

\subsubsection*{Bước 2 -- cộng cột, lập hệ.} Các tổng viết riêng cho loại này là

\begin{align*}S_U&=w_0x_0+\cdots+w_{N-1}x_{N-1},\\S_V&=w_0\ln|y_0|+\cdots+w_{N-1}\ln|y_{N-1}|,\\S_{UV}&=w_0x_0\ln|y_0|+\cdots+w_{N-1}x_{N-1}\ln|y_{N-1}|.\end{align*}

$S_1=w_0+\cdots+w_{N-1}$; $S_{UU}$ là tổng cột bình phương ở bảng trên, tức $w_0U_0^2+\cdots+w_{N-1}U_{N-1}^2$. Hệ cần giải:

\[\boxed{\begin{cases}S_1C+S_Ub=S_V,\\S_UC+S_{UU}b=S_{UV}.\end{cases}}\]

\subsubsection*{Bước 3 -- output và cách đọc.} Trả $A=se^C$ với $s=\operatorname{sign}(y_0)$; viết $g=Ae^{bx}$. $b$ là hệ số của $x$ trong số mũ, không phải cơ số.

Tool cực tiểu hóa tổng bình phương có trọng số trên thang $\ln|y|$. Sai số trên thang $y$ còn được báo thêm sau khôi phục. Đây là phép giải hữu hạn; không có bảng vòng hội tụ tham số ở nhánh này.

\subsection{Nhiều cơ sở trong số mũ: Ae mũ (bx bình phương + c/x)}
Input cơ sở \texttt{x\^{}2;1/x}, không thêm 1. Hàng $i$ là $(1,u_i,v_i)$ với $u_i=x_i^2,v_i=1/x_i$; đích $z_i=\ln|y_i|$. Lập đủ các cột $w_iu_i,w_iv_i,w_iu_i^2,w_iu_iv_i,w_iv_i^2,w_iz_i,w_iu_iz_i,w_iv_iz_i$.
Đặt $\langle f\rangle=w_0f(x_0,y_0)+\cdots+w_{N-1}f(x_{N-1},y_{N-1})$. Khi ấy
\[
\begin{bmatrix}\langle1\rangle&\langle x^2\rangle&\langle1/x\rangle\\
\langle x^2\rangle&\langle x^4\rangle&\langle x\rangle\\
\langle1/x\rangle&\langle x\rangle&\langle1/x^2\rangle\end{bmatrix}
\begin{bmatrix}C\\b\\c\end{bmatrix}
=\begin{bmatrix}\langle\ln|y|\rangle\\\langle x^2\ln|y|\rangle\\\langle\ln|y|/x\rangle\end{bmatrix}.
\]
Ví dụ $\langle x\rangle=w_0x_0+\cdots+w_{N-1}x_{N-1}$ do $u_iv_i=x_i$;
$\langle x^2\ln|y|\rangle=w_0x_0^2\ln|y_0|+\cdots+w_{N-1}x_{N-1}^2\ln|y_{N-1}|$.
Output $A=se^C$, $g=A\exp(bx^2+c/x)$. Đổi cơ sở thành \texttt{x;1/x} thì phải lập lại các tích từ hàng $(1,x_i,1/x_i)$.


\examplebreak
\section{Log của tổ hợp cơ sở: ln(a+bx) và dạng nhiều cơ sở}

\label{sec:ls-log-combination}

\subsubsection*{Input trong tool.} Chọn \textbf{Log của tổ hợp cơ sở}; nhập $x,y$, trọng số nếu có và điểm dự đoán nếu cần. Nhập đầy đủ cơ sở trong log: \texttt{1;x} cho $\ln(a+bx)$; không tự thêm cột 1. $y_i$ hữu hạn bất kỳ, nhưng $e^{y_i}$ phải không tràn/hụt số.

\subsubsection*{Bước 1 -- biến đổi từng dòng.}\[y\approx\ln(a+bx)\quad\Longrightarrow\quad e^y\approx a+bx.\]

\begin{center}\small\begin{tabular}{cllll}\toprule

Dòng&$U_i$&$V_i$&$w_iU_i^2$&$w_iU_iV_i$\\\midrule

$0$&$x_0$&$e^{y_0}$&$w_{0}(x_0)^2$&$w_{0}(x_0)(e^{y_0})$\\

$i$&$x_{i}$&$e^{y_{i}}$&$w_{i}(x_{i})^2$&$w_{i}(x_{i})(e^{y_{i}})$\\

$N-1$&$x_{N-1}$&$e^{y_{N-1}}$&$w_{N-1}(x_{N-1})^2$&$w_{N-1}(x_{N-1})(e^{y_{N-1}})$\\

\bottomrule\end{tabular}\end{center}

\subsubsection*{Bước 2 -- cộng cột, lập hệ.} Các tổng viết riêng cho loại này là

\begin{align*}S_U&=w_0x_0+\cdots+w_{N-1}x_{N-1},\\S_V&=w_0e^{y_0}+\cdots+w_{N-1}e^{y_{N-1}},\\S_{UV}&=w_0x_0e^{y_0}+\cdots+w_{N-1}x_{N-1}e^{y_{N-1}}.\end{align*}

$S_1=w_0+\cdots+w_{N-1}$; $S_{UU}$ là tổng cột bình phương ở bảng trên, tức $w_0U_0^2+\cdots+w_{N-1}U_{N-1}^2$. Hệ cần giải:

\[\boxed{\begin{cases}S_1a+S_Ub=S_V,\\S_Ua+S_{UU}b=S_{UV}.\end{cases}}\]

\subsubsection*{Bước 3 -- output và cách đọc.} Giữ $a,b$ và viết $g=\ln(a+bx)$; không đổi $a$ thành $e^a$. Kiểm $a+bx>0$ tại các điểm cần đánh giá. Với nhiều cơ sở, giữ $V_i=e^{y_i}$ và lập thiết kế từ mọi cơ sở bên trong log.

Tool cực tiểu hóa tổng bình phương có trọng số trên thang $e^y$. Sai số trên thang $y$ còn được báo thêm sau khôi phục. Đây là phép giải hữu hạn; không có bảng vòng hội tụ tham số ở nhánh này.

\examplebreak
\section{1/(a+b*x): biến đổi nghịch đảo tung độ}

\label{sec:ls-reciprocal}

\subsubsection*{Input trong tool.} Chọn \textbf{1/\allowbreak{}(a+b*x)}; nhập $x,y$, trọng số nếu có và điểm dự đoán nếu cần. Không nhập cơ sở. Cần mọi $y_i\ne0$; có thể khác dấu. Điểm dự đoán cần $a+bx_*\ne0$.

\subsubsection*{Bước 1 -- biến đổi từng dòng.}\[y\approx\frac1{a+bx}\quad\Longrightarrow\quad\frac1y\approx a+bx.\]

\begin{center}\small\begin{tabular}{cllll}\toprule

Dòng&$U_i$&$V_i$&$w_iU_i^2$&$w_iU_iV_i$\\\midrule

$0$&$x_0$&$1/y_0$&$w_{0}(x_0)^2$&$w_{0}(x_0)(1/y_0)$\\

$i$&$x_{i}$&$1/y_{i}$&$w_{i}(x_{i})^2$&$w_{i}(x_{i})(1/y_{i})$\\

$N-1$&$x_{N-1}$&$1/y_{N-1}$&$w_{N-1}(x_{N-1})^2$&$w_{N-1}(x_{N-1})(1/y_{N-1})$\\

\bottomrule\end{tabular}\end{center}

\subsubsection*{Bước 2 -- cộng cột, lập hệ.} Các tổng viết riêng cho loại này là

\begin{align*}S_U&=w_0x_0+\cdots+w_{N-1}x_{N-1},\\S_V&=w_0/y_0+\cdots+w_{N-1}/y_{N-1},\\S_{UV}&=w_0x_0/y_0+\cdots+w_{N-1}x_{N-1}/y_{N-1}.\end{align*}

$S_1=w_0+\cdots+w_{N-1}$; $S_{UU}$ là tổng cột bình phương ở bảng trên, tức $w_0U_0^2+\cdots+w_{N-1}U_{N-1}^2$. Hệ cần giải:

\[\boxed{\begin{cases}S_1a+S_Ub=S_V,\\S_Ua+S_{UU}b=S_{UV}.\end{cases}}\]

\subsubsection*{Bước 3 -- output và cách đọc.} Hệ số đã là $a,b$; viết $g=1/(a+bx)$, tính lại $g(x_i)$ để đánh giá sai số gốc. Mẫu gần 0 làm dự đoán nhạy.

Tool cực tiểu hóa tổng bình phương có trọng số trên thang $1/y$. Sai số trên thang $y$ còn được báo thêm sau khôi phục. Đây là phép giải hữu hạn; không có bảng vòng hội tụ tham số ở nhánh này.

\examplebreak
\section{x/(a+b*x): trả lại đúng thứ tự a,b}

\label{sec:ls-rational}

\subsubsection*{Input trong tool.} Chọn \textbf{x/\allowbreak{}(a+b*x)}; nhập $x,y$, trọng số nếu có và điểm dự đoán nếu cần. Không nhập cơ sở. Khi khớp cần mọi $x_i\ne0,y_i\ne0$. Tool tự dùng hai cơ sở $(1,1/x)$.

\subsubsection*{Bước 1 -- biến đổi từng dòng.}\[y\approx\frac{x}{a+bx}\quad\Longrightarrow\quad\frac1y\approx b+a\frac1x.\]

\begin{center}\small\begin{tabular}{cllll}\toprule

Dòng&$U_i$&$V_i$&$w_iU_i^2$&$w_iU_iV_i$\\\midrule

$0$&$1/x_0$&$1/y_0$&$w_{0}(1/x_0)^2$&$w_{0}(1/x_0)(1/y_0)$\\

$i$&$1/x_{i}$&$1/y_{i}$&$w_{i}(1/x_{i})^2$&$w_{i}(1/x_{i})(1/y_{i})$\\

$N-1$&$1/x_{N-1}$&$1/y_{N-1}$&$w_{N-1}(1/x_{N-1})^2$&$w_{N-1}(1/x_{N-1})(1/y_{N-1})$\\

\bottomrule\end{tabular}\end{center}

\subsubsection*{Bước 2 -- cộng cột, lập hệ.} Các tổng viết riêng cho loại này là

\begin{align*}S_U&=w_0/x_0+\cdots+w_{N-1}/x_{N-1},\\S_V&=w_0/y_0+\cdots+w_{N-1}/y_{N-1},\\S_{UV}&=w_0/(x_0y_0)+\cdots+w_{N-1}/(x_{N-1}y_{N-1}).\end{align*}

$S_1=w_0+\cdots+w_{N-1}$; $S_{UU}$ là tổng cột bình phương ở bảng trên, tức $w_0U_0^2+\cdots+w_{N-1}U_{N-1}^2$. Hệ cần giải:

\[\boxed{\begin{cases}S_1b+S_Ua=S_V,\\S_Ub+S_{UU}a=S_{UV}.\end{cases}}\]

\subsubsection*{Bước 3 -- output và cách đọc.} Bảng tuyến tính cho $(c_0,c_1)=(b,a)$, phải đổi lại $a=c_1,b=c_0$ rồi viết $g=x/(a+bx)$. Không thay bước đổi biến bằng $V=x/y$ khi mô tả tool: mục tiêu biến đổi sẽ khác. Tại điểm mới $x_*=0$, mô hình gốc có thể tính được nếu $a\ne0$.

Tool cực tiểu hóa tổng bình phương có trọng số trên thang $1/y$ với cơ sở $1/x$. Sai số trên thang $y$ còn được báo thêm sau khôi phục. Đây là phép giải hữu hạn; không có bảng vòng hội tụ tham số ở nhánh này.

\examplebreak
\section{a+b*ln(x): biến đổi hoành độ}

\label{sec:ls-plus-log}

\subsubsection*{Input trong tool.} Chọn \textbf{a+b*ln(x)}; nhập $x,y$, trọng số nếu có và điểm dự đoán nếu cần. Không nhập cơ sở. Cần mọi $x_i>0$; $y_i$ hữu hạn bất kỳ. Điểm dự đoán cũng cần $x_*>0$.

\subsubsection*{Bước 1 -- biến đổi từng dòng.}\[y\approx a+b\ln x\quad\Longrightarrow\quad V\approx a+bU,\quad U=\ln x,\ V=y.\]

\begin{center}\small\begin{tabular}{cllll}\toprule

Dòng&$U_i$&$V_i$&$w_iU_i^2$&$w_iU_iV_i$\\\midrule

$0$&$\ln x_0$&$y_0$&$w_{0}(\ln x_0)^2$&$w_{0}(\ln x_0)(y_0)$\\

$i$&$\ln x_{i}$&$y_{i}$&$w_{i}(\ln x_{i})^2$&$w_{i}(\ln x_{i})(y_{i})$\\

$N-1$&$\ln x_{N-1}$&$y_{N-1}$&$w_{N-1}(\ln x_{N-1})^2$&$w_{N-1}(\ln x_{N-1})(y_{N-1})$\\

\bottomrule\end{tabular}\end{center}

\subsubsection*{Bước 2 -- cộng cột, lập hệ.} Các tổng viết riêng cho loại này là

\begin{align*}S_U&=w_0\ln x_0+\cdots+w_{N-1}\ln x_{N-1},\\S_V&=w_0y_0+\cdots+w_{N-1}y_{N-1},\\S_{UV}&=w_0y_0\ln x_0+\cdots+w_{N-1}y_{N-1}\ln x_{N-1}.\end{align*}

$S_1=w_0+\cdots+w_{N-1}$; $S_{UU}$ là tổng cột bình phương ở bảng trên, tức $w_0U_0^2+\cdots+w_{N-1}U_{N-1}^2$. Hệ cần giải:

\[\boxed{\begin{cases}S_1a+S_Ub=S_V,\\S_Ua+S_{UU}b=S_{UV}.\end{cases}}\]

\subsubsection*{Bước 3 -- output và cách đọc.} Giữ trực tiếp $a,b$, viết $g=a+b\ln x$; không mũ hóa $a$. Chỉ đổi hoành độ nên mục tiêu vẫn ở thang tung độ gốc.

Tool cực tiểu hóa tổng bình phương có trọng số trên thang $y$ gốc. Sai số trên thang $y$ còn được báo thêm sau khôi phục. Đây là phép giải hữu hạn; không có bảng vòng hội tụ tham số ở nhánh này.

\examplebreak
\section{A*b mũ x: khôi phục biên độ và cơ số}

\label{sec:ls-geometric}

\subsubsection*{Input trong tool.} Chọn \textbf{A*b\textasciicircum{}x}; nhập $x,y$, trọng số nếu có và điểm dự đoán nếu cần. Không nhập cơ sở. Cần $y_i$ cùng dấu, khác 0; $x_i$ hữu hạn bất kỳ. Ký hiệu cơ số bằng $B$ để phân biệt hệ số mũ của $Ae^{bx}$; dạng hỗ trợ có $B>0$.

\subsubsection*{Bước 1 -- biến đổi từng dòng.}\[y\approx AB^x\quad\Longrightarrow\quad\ln|y|\approx C+Dx,\quad C=\ln|A|,\ D=\ln B.\]

\begin{center}\small\begin{tabular}{cllll}\toprule

Dòng&$U_i$&$V_i$&$w_iU_i^2$&$w_iU_iV_i$\\\midrule

$0$&$x_0$&$\ln|y_0|$&$w_{0}(x_0)^2$&$w_{0}(x_0)(\ln|y_0|)$\\

$i$&$x_{i}$&$\ln|y_{i}|$&$w_{i}(x_{i})^2$&$w_{i}(x_{i})(\ln|y_{i}|)$\\

$N-1$&$x_{N-1}$&$\ln|y_{N-1}|$&$w_{N-1}(x_{N-1})^2$&$w_{N-1}(x_{N-1})(\ln|y_{N-1}|)$\\

\bottomrule\end{tabular}\end{center}

\subsubsection*{Bước 2 -- cộng cột, lập hệ.} Các tổng viết riêng cho loại này là

\begin{align*}S_U&=w_0x_0+\cdots+w_{N-1}x_{N-1},\\S_V&=w_0\ln|y_0|+\cdots+w_{N-1}\ln|y_{N-1}|,\\S_{UV}&=w_0x_0\ln|y_0|+\cdots+w_{N-1}x_{N-1}\ln|y_{N-1}|.\end{align*}

$S_1=w_0+\cdots+w_{N-1}$; $S_{UU}$ là tổng cột bình phương ở bảng trên, tức $w_0U_0^2+\cdots+w_{N-1}U_{N-1}^2$. Hệ cần giải:

\[\boxed{\begin{cases}S_1C+S_UD=S_V,\\S_UC+S_{UU}D=S_{UV}.\end{cases}}\]

\subsubsection*{Bước 3 -- output và cách đọc.} Trả $A=se^C$, $B=e^D$, $s=\operatorname{sign}(y_0)$; viết $g=AB^x$. Không lấy $D$ làm cơ số. Ví dụ $y=(2,6,18)$ tại $x=(0,1,2)$ cho $g=2\cdot3^x$.

Tool cực tiểu hóa tổng bình phương có trọng số trên thang $\ln|y|$. Sai số trên thang $y$ còn được báo thêm sau khôi phục. Đây là phép giải hữu hạn; không có bảng vòng hội tụ tham số ở nhánh này.


\examplebreak
\section{Lũy thừa -- SSE thang y: lặp trực tiếp cho Ax mũ b}
\label{sec:ls-power-raw}
\subsubsection*{Input.} Chọn \textbf{Lũy thừa -- SSE thang y}, nhập $x_i>0$, $y_i$ hữu hạn bất kỳ; không cần $y$ cùng dấu. Không có ô cơ sở. Tham số đầu theo thứ tự \texttt{A b}, ngưỡng dừng mặc định $10^{-10}$, tối đa 100 vòng.
\subsubsection*{Quá trình riêng của mô hình.} Tại bộ $(A,b)$ hiện tại, tính từng dòng:
\begin{center}\small\begin{tabular}{cllll}\toprule
Dòng&$g_i$&$g_i-y_i$&$\partial g_i/\partial A$&$\partial g_i/\partial b$\\\midrule
0&$Ax_0^b$&$Ax_0^b-y_0$&$x_0^b$&$Ax_0^b\ln x_0$\\
$i$&$Ax_i^b$&$Ax_i^b-y_i$&$x_i^b$&$Ax_i^b\ln x_i$\\
$N-1$&$Ax_{N-1}^b$&$Ax_{N-1}^b-y_{N-1}$&$x_{N-1}^b$&$Ax_{N-1}^b\ln x_{N-1}$\\\bottomrule
\end{tabular}\end{center}
Mục tiêu $E=w_0(Ax_0^b-y_0)^2+\cdots+w_{N-1}(Ax_{N-1}^b-y_{N-1})^2$. Dùng bảng này tạo phần dư và Jacobian, chạy chu trình LM ở mục \ref{sec:ls-lm-tool}, tính lại bảng sau mỗi bước nhận. Không dùng $\ln y_i$ làm đích của mục tiêu.
\subsubsection*{Output.} $(A,b)$ theo đơn vị gốc, $g=Ax^b$, mọi bước thử nhận/từ chối, SSE và lý do dừng. Khi chưa đạt ngưỡng, tham số cuối chỉ để kiểm tra.

\examplebreak
\section{Hàm mũ -- SSE thang y: một hoặc nhiều cơ sở số mũ}
\label{sec:ls-exp-raw}
\subsubsection*{Input.} Chọn \textbf{Hàm mũ -- SSE thang y}. Cơ sở chỉ nhập phần số mũ, ví dụ \texttt{x} hoặc \texttt{x\^{}2;1/x}, không thêm 1. Với $q$ cơ sở, tham số đầu là \texttt{A b1 ... bq}. $y$ có thể khác dấu/0 nhưng mọi cơ sở phải xác định. Một mô hình mũ có biên độ cố định vẫn có một dấu; cho phép dữ liệu đổi dấu không có nghĩa khớp đúng mọi dòng.
\subsubsection*{Quá trình riêng của mô hình.}
\[
L_i=\sum_{j=1}^qb_j\psi_j(x_i)=b_1\psi_1(x_i)+\cdots+b_q\psi_q(x_i),\quad g_i=Ae^{L_i}.
\]
Jacobian hàng $i$ theo $(A,b_1,\ldots,b_q)$ là
\[
(e^{L_i},\ Ae^{L_i}\psi_1(x_i),\ldots,Ae^{L_i}\psi_q(x_i)).
\]
Ví dụ \texttt{x\^{}2;1/x}, $L_i=bx_i^2+c/x_i$:
\begin{center}\small\begin{tabular}{cllll}\toprule
Dòng&$g_i$&Cột theo $A$&Cột theo $b$&Cột theo $c$\\\midrule
0&$Ae^{L_0}$&$e^{L_0}$&$Ae^{L_0}x_0^2$&$Ae^{L_0}/x_0$\\
$i$&$Ae^{L_i}$&$e^{L_i}$&$Ae^{L_i}x_i^2$&$Ae^{L_i}/x_i$\\
$N-1$&$Ae^{L_{N-1}}$&$e^{L_{N-1}}$&$Ae^{L_{N-1}}x_{N-1}^2$&$Ae^{L_{N-1}}/x_{N-1}$\\\bottomrule
\end{tabular}\end{center}
Mục tiêu gốc $E=w_0(Ae^{L_0}-y_0)^2+\cdots+w_{N-1}(Ae^{L_{N-1}}-y_{N-1})^2$. Tính mọi dòng bằng cùng một bộ tham số rồi mới cộng; chạy chu trình dưới đây.
\subsubsection*{Output.} Vector $(A,b_1,\ldots,b_q)$, công thức $g$, bảng vòng lặp, SSE gốc, lý do dừng, dự đoán tùy chọn. Thứ tự cơ sở quyết định ý nghĩa từng tham số.

\examplebreak
\section{Chu trình LM: một vòng nhận gì, tính gì, chuyển gì sang vòng sau?}
\label{sec:ls-lm-tool}
\subsubsection*{Khởi tạo.} Nhập tham số đầu thì dùng đúng bộ đó. Để trống: nếu mọi $y$ cùng dấu khác 0, khởi tạo từ khớp log; trường hợp còn lại dùng trung bình có trọng số làm biên độ (nếu bằng 0 lấy giá trị khởi tạo nội bộ 1), các hệ số mũ bằng 0. Đây chỉ là khởi tạo, không đổi mục tiêu SSE gốc.

Tool đặt $s_y=\max_{w_i>0}|y_i|$ (bằng 0 thì lấy 1), $w_{\max}=\max w_i$, $\alpha=A/s_y$, $\rho_i=\sqrt{w_i/w_{\max}}$. Tham số nội bộ $p=(\alpha,b_1,\ldots,b_q)$. Tại các dòng có trọng số dương:
\[
r_i=\rho_i(\alpha e^{L_i}-y_i/s_y),\quad
J_i=\rho_i(e^{L_i},\alpha e^{L_i}\psi_1(x_i),\ldots,\alpha e^{L_i}\psi_q(x_i)).
\]
Dòng $w_i=0$ có phần dư nội bộ bằng 0. $F=r_0^2+\cdots+r_{N-1}^2$ là mục tiêu nội bộ; SSE có trọng số gốc là $E=s_y^2w_{\max}F$. Tham số hiển thị đổi lại $A=s_y\alpha$.
\begin{enumerate}
\item Nhận $p_k,r_k,J_k,F_k,\lambda_k$ từ vòng trước; ban đầu $\lambda_0=10^{-3}$.
\item Tính $h=J_k^Tr_k$, với $h_j=J_{0j}r_0+\cdots+J_{N-1,j}r_{N-1}$. Nếu $\max_j|h_j|\le\tau$, dừng trước khi tạo bước thử mới.
\item Tính $d_j=\sqrt{J_{0j}^2+\cdots+J_{N-1,j}^2}$, nếu bằng 0 thay 1; giải bước
\[
\begin{bmatrix}J_k\\\sqrt{\lambda_k}\operatorname{diag}(d_j)\end{bmatrix}\Delta p
\approx\begin{bmatrix}-r_k\\0\end{bmatrix}.
\]
\item Tạo $p^{\rm thử}=p_k+\Delta p$, tính lại mô hình, phần dư, Jacobian và $F^{\rm thử}$ trên toàn bảng.
\item Chỉ nhận nếu bước hữu hạn và $F^{\rm thử}<F_k$. Khi nhận, chuyển toàn bộ bộ mới sang vòng sau, $\lambda\leftarrow\max(\lambda/3,10^{-15})$. Khi từ chối, giữ bộ cũ, $\lambda\leftarrow10\lambda$, vẫn lưu bước thử.
\item Sau bước nhận còn kiểm đồng thời $\|\Delta p\|_2/(1+\|p_{k+1}\|_2)\le\tau$ và $|F_k-F_{k+1}|\le\tau(1+F_k)$. Nếu chưa đạt thì tiếp tục. Hết vòng chưa đạt phải báo chưa hội tụ.
\end{enumerate}
\subsubsection*{Bảng tổng quát.} Chia thành hai bảng cùng khóa vòng $k$:
\begin{center}\small\begin{tabular}{cllll}\toprule
Vòng&Tham số cũ&Tham số thử&$\lambda$ dùng&Quyết định\\\midrule
1&$p_0$&$p_0+\Delta p_0$&$10^{-3}$&Nhận/từ chối\\
$k+1$&$p_k$&$p_k+\Delta p_k$&$\lambda_k$&So mục tiêu mới/cũ\\\bottomrule
\end{tabular}\end{center}
\begin{center}\small\begin{tabular}{cllll}\toprule
Vòng&SSE trước&SSE thử&$\|J^Tr\|_\infty$&$\|\Delta p\|_2$\\\midrule
1&$s_y^2w_{\max}F_0$&$s_y^2w_{\max}F^{\rm thử}$&Gradient trước bước&Bước nội bộ\\
$k+1$&$s_y^2w_{\max}F_k$&$s_y^2w_{\max}F^{\rm thử}$&Gradient trước bước&Bước nội bộ\\\bottomrule
\end{tabular}\end{center}
Mở chi tiết để xem $J,r$ dùng tạo bước. Không đưa tham số thử bị từ chối sang vòng sau. Đạt tiêu chuẩn số chưa chứng minh cực tiểu toàn cục.

\examplebreak
\section{Output: bảng sai số tổng quát và kết luận bài khớp}
\label{sec:ls-output}
Sau khi khôi phục tham số, thế lại mọi $x_i$ bằng mô hình gốc:
\begin{center}\small\begin{tabular}{clllll}\toprule
Dòng&$y_i$&$g_i$&$e_i=y_i-g_i$&$e_i^2$&$w_ie_i^2$\\\midrule
0&$y_0$&$g(x_0)$&$y_0-g(x_0)$&$e_0^2$&$w_0e_0^2$\\
$i$&$y_i$&$g(x_i)$&$y_i-g(x_i)$&$e_i^2$&$w_ie_i^2$\\
$N-1$&$y_{N-1}$&$g(x_{N-1})$&$y_{N-1}-g(x_{N-1})$&$e_{N-1}^2$&$w_{N-1}e_{N-1}^2$\\\midrule
Tổng&&&&$\mathrm{SSE}_y$&$E_y$\\\bottomrule
\end{tabular}\end{center}
\begin{align*}
\mathrm{SSE}_y&=\sum_ie_i^2=e_0^2+e_1^2+\cdots+e_{N-1}^2,\\
E_y&=\sum_iw_ie_i^2=w_0e_0^2+\cdots+w_{N-1}e_{N-1}^2,\\
\mathrm{MSE}_y&=\mathrm{SSE}_y/N,\quad \mathrm{RMSE}_y=\sqrt{\mathrm{SSE}_y/N},\\
\mathrm{RMSE}_w&=\sqrt{E_y/(w_0+w_1+\cdots+w_{N-1})}.
\end{align*}
Bài biến đổi còn có $\widehat V_i$ và phần dư $V_i-\widehat V_i$ trên thang đang khớp. Không coi SSE gốc là mục tiêu đã cực tiểu nếu tool đang khớp log/nghịch đảo/mũ hóa. $R^2$ nếu có không phải cận sai số. Sai số thống kê chia $N-p$ khác RMSE chia $N$.

Nếu cơ sở thiếu hạng, mặc định báo tham số chưa duy nhất; chỉ khi đã chủ động chọn SVD và quy tắc chuẩn nhỏ nhất mới nhận nghiệm kiểu đó. Điểm dự đoán phải thuộc miền mô hình khôi phục. Một điểm dự đoán sai miền không làm dữ liệu khớp trước đó tự thay đổi.

\subsubsection*{Bài thi cần ghi.} Tên mô hình và cơ sở $\to$ miền áp dụng $\to$ bảng biến đổi $\to$ các tích và tổng $\to$ hệ/bài tối ưu $\to$ tham số khôi phục $\to$ hàm $g$ $\to$ sai số đúng thang $\to$ dự đoán. Với SSE phi tuyến thêm mọi bước thử và lý do dừng. Các ví dụ số đã kiểm ở \texttt{verification/spline\_ls\_examples.py}; đề thật ở Phụ lục \ref{app:exam}.

\chapter{Đạo hàm số và tích phân số}
\label{ch:calculus}
\section{Tool Đạo hàm từ bảng: dữ liệu qua bảng trọng số}
\paragraph{Input trên GUI.} Trong nhóm \emph{Đạo hàm \& Horner}, chọn \textbf{Đạo hàm từ bảng}. Nhập hai dãy $x,y$; ô \emph{Điểm x} là $z$; \emph{Bậc p} dùng đúng $p+1$ mốc; \emph{Cấp đạo hàm} là $s$. Các $x$ phải phân biệt, không cần đều; $p$ không phải cấp đạo hàm. Tên khóa JSON là \texttt{target}, \texttt{degree}, \texttt{order}.
\paragraph{Quá trình.} Chọn $p+1$ mốc theo quy tắc của tool $\to$ lập từng cơ sở Lagrange $L_j$ $\to$ tính trọng số $w_j=L_j^{(s)}(z)$ $\to$ nhân với $y_j$ $\to$ cộng. Mẫu bảng tổng quát:
\begin{center}\small\begin{tabular}{clllll}\toprule
Mốc&$x_j$&$y_j$&$w_j$&Đóng góp $w_jy_j$&Tổng dồn\\\midrule
0&$x_0$&$y_0$&$L_0^{(s)}(z)$&$w_0y_0$&$w_0y_0$\\
$j$&$x_j$&$y_j$&$L_j^{(s)}(z)$&$w_jy_j$&$w_0y_0+\cdots+w_jy_j$\\
$p$&$x_p$&$y_p$&$L_p^{(s)}(z)$&$w_py_p$&Kết quả cuối\\\bottomrule
\end{tabular}\end{center}
\[
f^{(s)}(z)\approx\sum_{j=0}^pw_jy_j=w_0y_0+w_1y_1+\cdots+w_py_p.
\]
\textbf{Output}: tập mốc thực sự dùng, bảng tạo trọng số/đóng góp, đạo hàm gần đúng và kiểm các moment. Nếu đã chọn đạo hàm trong tool nội suy thì đọc đạo hàm của chính đa thức đó; đây là lựa chọn riêng có thể dùng tập mốc khác.

\section{Các tool Tích phân: input và bảng tổng quát trước ví dụ}
\subsection*{Input: các ô đúng tên trên giao diện}
Bảng này được lấy từ biểu mẫu của tool. Giá trị mặc định là ví dụ nhập, cần thay bằng dữ liệu đề. Hai dãy $x,y$ là bảng dữ liệu riêng bên dưới nếu chức năng dùng bảng.
\subsubsection*{Ô chung của các lựa chọn trong nhóm}
\begin{center}\small\begin{tabular}{@{}p{.32\linewidth}p{.17\linewidth}p{.43\linewidth}@{}}\toprule
Ô nhập & Khóa JSON & Cách nhập / mặc định\\\midrule
Cận trái a & \texttt{a} & \texttt{0}\\[3pt]
Cận phải b & \texttt{b} & \texttt{1}\\[3pt]
Số khoảng chia n & \texttt{n} & \texttt{20}\\[3pt]
\bottomrule\end{tabular}\end{center}
\subsubsection*{Ô riêng: Hình thang /\allowbreak{} Simpson /\allowbreak{} Trung điểm}
\begin{center}\small\begin{tabular}{@{}p{.32\linewidth}p{.17\linewidth}p{.43\linewidth}@{}}\toprule
Ô nhập & Khóa JSON & Cách nhập / mặc định\\\midrule
Hàm f(x) & \texttt{expr} & \texttt{x\textasciicircum{}2}\\[3pt]
Dung sai tự chia lưới (để trống: dùng n) & \texttt{tol} & Để trống
\par Chỉ dùng với hàm;\allowbreak{} tự tăng n đến khi ước lượng Runge đủ nhỏ.\\[3pt]
Cận đạo hàm M (không bắt buộc) & \texttt{bound} & Để trống
\par M2 cho hình thang/\allowbreak{}trung điểm,\allowbreak{} M4 cho Simpson.\\[3pt]
Nguồn dữ liệu & \texttt{source} & \texttt{Hàm}
\par Nhập Hàm hoặc Bảng;\allowbreak{} bảng x,\allowbreak{}y điền bên dưới.\\[3pt]
\bottomrule\end{tabular}\end{center}
\subsubsection*{Ô riêng: Newton-Cotes}
\begin{center}\small\begin{tabular}{@{}p{.32\linewidth}p{.17\linewidth}p{.43\linewidth}@{}}\toprule
Ô nhập & Khóa JSON & Cách nhập / mặc định\\\midrule
Hàm f(x) & \texttt{expr} & \texttt{x\textasciicircum{}2}\\[3pt]
Bậc Newton-Cotes & \texttt{degree} & \texttt{4}
\par n phải chia hết cho bậc này.\\[3pt]
Dung sai tự chia lưới (để trống: dùng n) & \texttt{tol} & Để trống
\par Chỉ dùng với hàm;\allowbreak{} tự tăng n đến khi ước lượng Runge đủ nhỏ.\\[3pt]
Cận đạo hàm M (không bắt buộc) & \texttt{bound} & Để trống
\par M2 cho hình thang/\allowbreak{}trung điểm,\allowbreak{} M4 cho Simpson.\\[3pt]
Nguồn dữ liệu & \texttt{source} & \texttt{Hàm}
\par Nhập Hàm hoặc Bảng;\allowbreak{} bảng x,\allowbreak{}y điền bên dưới.\\[3pt]
\bottomrule\end{tabular}\end{center}
\subsubsection*{Ô riêng: Tích phân kép}
\begin{center}\small\begin{tabular}{@{}p{.32\linewidth}p{.17\linewidth}p{.43\linewidth}@{}}\toprule
Ô nhập & Khóa JSON & Cách nhập / mặc định\\\midrule
Hàm f(x,\allowbreak{}y) & \texttt{expr} & \texttt{x\textasciicircum{}2+y\textasciicircum{}2}\\[3pt]
Cận dưới y & \texttt{c} & \texttt{0}\\[3pt]
Cận trên y & \texttt{d} & \texttt{1}\\[3pt]
Quy tắc & \texttt{rule} & \texttt{Simpson}
\par Nhập Simpson hoặc Hình thang.\\[3pt]
Dung sai Runge & \texttt{tol} & \texttt{1e-07}\\[3pt]
\bottomrule\end{tabular}\end{center}

\paragraph{Hai kiểu input.} Nguồn \emph{Hàm}: nhập biểu thức, cận, số khoảng $n$; tool tự đánh giá hàm. Nguồn \emph{Bảng}: $x,y$ phải khớp đúng lưới mà quy tắc cần; không tự tạo thêm dữ liệu ở điểm chưa biết. Ô \emph{Dung sai tự chia lưới} chỉ dùng với hàm. Ô $n$ là số khoảng, thường có $n+1$ nút; trung điểm dùng $n$ tâm.

\subsection{Hình thang: từ n+1 nút đến một tổng có trọng số}
\textbf{Input riêng}: $n\ge1$, $h=(b-a)/n$, $x_i=a+ih$.
\begin{center}\small\begin{tabular}{cllll}\toprule
Dòng&Vị trí&Giá trị&Trọng số $W_i$&Đóng góp\\\midrule
0&$a$&$y_0=f(a)$&$1/2$&$y_0/2$\\
$i$, $1\le i<n$&$a+ih$&$y_i$&1&$y_i$\\
$n$&$b$&$y_n=f(b)$&$1/2$&$y_n/2$\\\bottomrule
\end{tabular}\end{center}
Cộng cột đóng góp rồi nhân $h$:
\[
I_n=h\sum_{i=0}^nW_iy_i=h(y_0/2+y_1+y_2+\cdots+y_{n-1}+y_n/2).
\]
\textbf{Output}: lưới, bảng đóng góp, $I_n$; nếu yêu cầu tinh chỉnh còn có bảng Runge/cận và trạng thái.

\subsection{Trung điểm: giá trị phải nằm ở tâm từng khoảng}
\textbf{Input riêng}: $n\ge1$, $h=(b-a)/n$. Tâm $\xi_i=a+(i+1/2)h$, $i=0,\ldots,n-1$.
\begin{center}\small\begin{tabular}{clll}\toprule
Khoảng&Hai đầu&Tâm lấy hàm&Đóng góp\\\midrule
0&$a,a+h$&$a+h/2$&$f(a+h/2)$\\
$i$&$a+ih,a+(i+1)h$&$a+(i+1/2)h$&$f(\xi_i)$\\
$n-1$&$b-h,b$&$b-h/2$&$f(b-h/2)$\\\bottomrule
\end{tabular}\end{center}
\[
I_n=h\sum_{i=0}^{n-1}f(\xi_i)=h[f(a+h/2)+f(a+3h/2)+\cdots+f(b-h/2)].
\]
\textbf{Output}: $n$ tâm và các giá trị tại tâm, tổng đóng góp, kết quả. Bảng chỉ biết hai đầu khoảng chưa tự có tung độ ở tâm.

\subsection{Simpson: kiểm n chẵn trước khi điền trọng số}
\textbf{Input riêng}: $n$ chẵn, $n\ge2$; $n+1$ nút đều kể cả hai đầu.
\begin{center}\small\begin{tabular}{clll}\toprule
Vị trí dòng&Giá trị&Trọng số nguyên&Đóng góp trong ngoặc\\\midrule
Đầu $i=0,n$&$y_i$&1&$y_i$\\
Trong, $i$ lẻ&$y_i$&4&$4y_i$\\
Trong, $i$ chẵn&$y_i$&2&$2y_i$\\\bottomrule
\end{tabular}\end{center}
\[
I_n=\frac h3\big[y_0+4y_1+2y_2+4y_3+\cdots+2y_{n-2}+4y_{n-1}+y_n\big].
\]
Nếu bảng tool đã dùng trọng số $1/3,4/3,2/3,\ldots$, chỉ nhân $h$ ở ngoài, không chia 3 lần nữa. \textbf{Output}: bảng lưới/trọng số, kết quả và thông tin sai số nếu đã yêu cầu.

\subsection{Newton--Cotes: ghép các khối bậc d}
\textbf{Input riêng}: bậc $1\le d\le10$, $n$ chia hết cho $d$. Một khối gồm $d$ khoảng, $d+1$ nút; trọng số $w_0,\ldots,w_d$ trong bảng bậc đã bao gồm hệ số chung.
\[
I_{\text{khối }k}=h(w_0y_{kd}+w_1y_{kd+1}+\cdots+w_dy_{kd+d}).
\]
\begin{center}\small\begin{tabular}{clll}\toprule
Khối&Chỉ số mốc&Đóng góp&Thông tin chuyển tiếp\\\midrule
0&$0,1,\ldots,d$&$h(w_0y_0+\cdots+w_dy_d)$&Cộng vào tổng\\
$k$&$kd,\ldots,kd+d$&$h(w_0y_{kd}+\cdots+w_dy_{kd+d})$&Nút chung nhận hai trọng số\\\bottomrule
\end{tabular}\end{center}
\textbf{Output}: trọng số chuẩn của bậc $d$, trọng số ghép trên toàn lưới, giá trị $I_n$. Các bậc 1--10 và ví dụ Boole nằm ở phần giải thích sau.

\subsection{Tự chia lưới: bảng tổng quát của quá trình kiểm sai số}
Có ngưỡng $\varepsilon$ và nguồn hàm thì tool tính lại trên lưới mịn, không chỉ làm tròn $I_n$. Bậc dùng Runge là $p=2$ cho hình thang/trung điểm, $p=4$ cho Simpson; Newton--Cotes tùy bậc.
\begin{center}\small\begin{tabular}{cllll}\toprule
Lượt&Lưới thô/mịn&Hai giá trị&Ước lượng lưới mịn&Quyết định\\\midrule
$k$&$n_k,2n_k$&$I_{n_k},I_{2n_k}$&$|I_{2n_k}-I_{n_k}|/(2^p-1)$&Đạt ngưỡng?\\\bottomrule
\end{tabular}\end{center}
App còn kiểm một lưới thăm dò không lồng để tránh kết luận từ hiệu hai lưới tình cờ nhỏ. Không đạt thì chuyển lưới mịn thành lưới thô của lượt sau. \textbf{Output}: số khoảng cuối, kết quả, lịch sử các lưới và lý do dừng. Runge là ước lượng theo mô hình hội tụ; cận lý thuyết chỉ dùng khi có cận đạo hàm phù hợp.

\subsection{Tích phân kép: ghép trọng số theo hai chiều}
\textbf{Input}: $f(x,y)$, hình chữ nhật $[a,b]\times[c,d]$, số khoảng $n$ mỗi chiều, quy tắc Simpson/hình thang và dung sai. Không dùng hai cột $x,y$ như một bảng giá trị hàm hai biến.
\textbf{Quá trình}: lập các nút $x_i,y_j$ $\to$ tính ma trận $f_{ij}=f(x_i,y_j)$ $\to$ nhân trọng số tích $W_iV_j$ $\to$ cộng. $W_i,V_j$ ở công thức dưới là trọng số \emph{đã chứa hệ số chung}: Simpson dùng $(1,4,2,\ldots,4,1)/3$ ở mỗi chiều, nên hệ số $1/9$ đã nằm trong tích; hình thang dùng $(1/2,1,\ldots,1,1/2)$.
\begin{align*}
I_n&=h_xh_y\sum_{i=0}^n\sum_{j=0}^nW_iV_jf_{ij}\\
 &=h_xh_y\{W_0[V_0f_{00}+\cdots+V_nf_{0n}]\\
 &\qquad+W_1[V_0f_{10}+\cdots+V_nf_{1n}]+\cdots+W_n[V_0f_{n0}+\cdots+V_nf_{nn}]\}.
\end{align*}
\begin{center}\small\begin{tabular}{clll}\toprule
Hàng $x_i$&Các giá trị cần có&Tổng theo y&Đóng góp hàng\\\midrule
$i$&$f_{i0},f_{i1},\ldots,f_{in}$&$T_i=V_0f_{i0}+\cdots+V_nf_{in}$&$W_iT_i$\\\bottomrule
\end{tabular}\end{center}
\textbf{Output}: giá trị tích phân, lưới và lịch sử tinh chỉnh theo dung sai; miền được hỗ trợ là chữ nhật. Các mục sau cho chi tiết cách tạo trọng số và ví dụ số.


\section{Nhìn đúng dữ liệu trước khi chọn công thức}
Đạo hàm số và tích phân số đều thay một phép toán trên hàm thật bằng phép toán trên dữ liệu hoặc một hàm xấp xỉ. Chúng không tự tạo thêm thông tin về độ trơn hay sai số đo. Bảng sau là đường đi thực tế trong app.
\begin{center}
\begin{tabular}{p{.26\textwidth}p{.28\textwidth}p{.35\textwidth}}
\toprule
Đề cho gì? & Chọn trong app & Việc đang làm\\\midrule
Bảng $(x_i,y_i)$, hỏi $f^{(s)}(z)$ & Đạo hàm từ bảng & Lấy đạo hàm đa thức qua $p+1$ mốc đã chọn.\\
Đã dựng một đa thức Newton/Stirling/Bessel & Nội suy đa thức, cấp đạo hàm $s$ & Giữ chính đa thức và tập mốc của câu nội suy, lấy đạo hàm theo đúng biến.\\
Biểu thức $f$, hỏi tích phân & Tích phân, nguồn Hàm & Tính $f$ tại các nút của quy tắc; có thể tinh chỉnh lưới.\\
Bảng hữu hạn, hỏi tích phân & Tích phân, nguồn Bảng & Dùng đúng lưới của bảng; không tự lấy giá trị hàm ở các nút mới.\\
Biểu thức $f(x,y)$ trên hình chữ nhật & Tích phân kép & Tích tensor của hai quy tắc một chiều, tinh chỉnh cả hai chiều.\\\bottomrule
\end{tabular}
\end{center}

Trong chương này $p$ là bậc đa thức dùng lấy đạo hàm; khi nói sai số tích phân, ký hiệu $r$ là bậc sai số toàn cục để tránh nhầm hai đại lượng. $n$ luôn là \emph{số khoảng chia}, không phải số nút. Với quy tắc chứa hai đầu, có $n+1$ nút; quy tắc trung điểm có $n$ nút.

\section{Đạo hàm từ p+1 mốc bất kỳ}
\label{sec:derivative-weights}
\subsection{Đầu vào, đầu ra và điều kiện}
Nhập hai dãy $x,y$, điểm tính $z$, bậc đa thức $p$ và cấp đạo hàm $s$. Trong màn hình này app yêu cầu
\[
1\le s\le p<m,
\]
với $m$ là số cặp nhập. Tất cả số phải hữu hạn, $x_i$ phân biệt. Mốc không cách đều vẫn hợp lệ. Đầu ra là
\[
f^{(s)}(z)\approx P_p^{(s)}(z)=\sum_{i=0}^{p}w_i y_i.
\]
Nếu đề cho sẵn một đa thức thì dùng Horner; không cần xấp xỉ bằng bảng. Nếu đề đã yêu cầu đạo hàm của Stirling chín mốc thì phải giữ tập chín mốc ấy, thay vì chuyển sang một cửa sổ ba mốc khác.

\subsection{Bước 1: chọn đúng các mốc sẽ dùng}
App sắp tăng $x$ cùng cả cặp và giữ ID dòng nguồn. Tìm hàng gần $z$ nhất, gọi chỉ số là $c$; khi hòa chọn hàng bên trái. Đặt
\[
\ell=\max\{0,\min(c-\lfloor p/2\rfloor,m-p-1)\}.
\]
Tập dùng là các hàng $\ell,\ldots,\ell+p$. Đây là \emph{cửa sổ liên tiếp quanh hàng gần nhất}; trên lưới không đều nó không nhất thiết là $p+1$ mốc gần $z$ nhất. Điểm ở biên làm cửa sổ lệch về một phía. Điểm nằm ngoài đoạn chọn vẫn tính được đạo hàm của $P_p$, nhưng đó là ngoại suy và phải được ghi rõ.

\subsection{Bước 2: từ cơ sở Lagrange đến trọng số}
Với các mốc đã chọn, đặt
\[
L_i(x)=\prod_{j\ne i}\frac{x-x_j}{x_i-x_j},\qquad
P_p(x)=\sum_i y_iL_i(x),\qquad w_i=L_i^{(s)}(z).
\]
Để tính $w_i$ mà không khai triển các lũy thừa lớn theo $x$, thay $x=z+u$:
\[
N_i(u)=\prod_{j\ne i}(z-x_j+u)=\sum_{k=0}^{p}a_k u^k,
\quad d_i=\prod_{j\ne i}(x_i-x_j),
\quad w_i=\frac{s!a_s}{d_i}.
\]
Chỉ cần giữ $a_0,\ldots,a_s$. Khởi tạo $(a_0,a_1,\ldots,a_s)=(1,0,\ldots,0)$. Mỗi khi nhân thêm $v+u$, $v=z-x_j$, cập nhật \emph{từ cấp cao xuống cấp thấp}:
\[
a_k\leftarrow va_k+a_{k-1}\quad(k=s,s-1,\ldots,1),\qquad a_0\leftarrow va_0.
\]
Đồng thời nhân $d_i$ với $x_i-x_j$. Cuối vòng qua các $j\ne i$, lấy $s!a_s/d_i$. Không có phép chia cho $z-x_j$, nên $z$ trùng một mốc không gây lỗi chia cho 0.

\subsection{Bước 3: nhân tung độ, cộng và kiểm tra}
Tính $q_i=w_i y_i$, rồi tổng dồn $T_i=T_{i-1}+q_i$, $T_{-1}=0$. Kết quả cuối là $T_p$. App đối chiếu các moment
\[
\sum_iw_i(x_i-z)^k=
\begin{cases}s!,&k=s,\\0,&0\le k\le p,\ k\ne s.\end{cases}
\]
Các đẳng thức này nói công thức lấy đúng đạo hàm của mọi đa thức bậc không quá $p$. Chúng kiểm tra cách dựng trọng số, không chứng minh các giá trị đo chính xác hay hàm thật đủ trơn.

\subsection{Ví dụ đầy đủ: ba điểm ở đầu bảng}
Ví dụ tự dựng: $x=(0,1,2)$, $y=(1,2,5)$, tính $f'(0)$ bằng đa thức bậc hai. Có $p=2,s=1,z=0$. Cả ba mốc đều được dùng.
\begin{center}
\begin{tabular}{rlll}
\toprule $i$&$L_i(x)$&$L_i'(x)$&$w_i=L_i'(0)$\\\midrule
0&$(x-1)(x-2)/2$&$x-3/2$&$-3/2$\\
1&$-x(x-2)$&$-2x+2$&$2$\\
2&$x(x-1)/2$&$x-1/2$&$-1/2$\\\bottomrule
\end{tabular}
\end{center}
Hàng $i=0$ cũng có thể làm theo đúng vòng nhân hệ số của app:
\begin{center}
\begin{tabular}{rllll}
\toprule Nhân tử & $v$ & $(a_0,a_1)$ sau nhân & Mẫu tích lũy\\\midrule
Khởi tạo&---&$(1,0)$&$1$\\
$u-1$&$-1$&$(-1,1)$&$(0-1)=-1$\\
$u-2$&$-2$&$(2,-3)$&$(-1)(0-2)=2$\\\bottomrule
\end{tabular}
\end{center}
Do đó $w_0=1!(-3)/2=-3/2$. Tương tự $w_1=2,w_2=-1/2$.
\begin{center}
\begin{tabular}{rrrrr}
\toprule $i$&$w_i$&$y_i$&$w_iy_i$&Tổng dồn\\\midrule
0&$-3/2$&1&$-3/2$&$-3/2$\\
1&2&2&4&$5/2$\\
2&$-1/2$&5&$-5/2$&0\\\bottomrule
\end{tabular}
\end{center}
Kết luận $f'(0)\approx0$. Kiểm tra: $\sum w_i=0$, $\sum w_ix_i=1$, $\sum w_ix_i^2=0$.

\subsection{Ví dụ mốc không đều và đạo hàm cấp hai}
Với $x=(0,1,3)$, $y=(1,3,13)$, $z=1$, $p=2,s=1$, thuật toán tổng quát cho
\[
(w_0,w_1,w_2)=(-2/3,1/2,1/6),\qquad
P'(1)=-\frac23+\frac32+\frac{13}{6}=3.
\]
Không có một $h$ chung nên không thay bằng công thức ba điểm giữa cách đều.

Với $x=(0,1,2)$, $y=(1,2,5)$, $z=1$, $p=2,s=2$, nhận
\[
(w_0,w_1,w_2)=(1,-2,1),\qquad P''(1)=1-4+5=2.
\]

\subsection{Khi đã đổi biến trong bài nội suy}
Nếu $x=a+ht$ và $Q(t)=P(a+ht)$ thì
\[
t_z=\frac{z-a}{h},\qquad P^{(s)}(z)=\frac{Q^{(s)}(t_z)}{h^s}.
\]
Thứ tự: giữ $Q$ vừa dựng; đạo hàm theo $t$; thế $t_z$; chia $h^s$. Ví dụ $Q(t)=16+16t+4t^2$, $x=4+2t$, tại $x=4$: $t_z=0$, $Q'(0)=16$, $P'(4)=16/2=8$; $Q''(0)=8$, $P''(4)=8/2^2=2$.

\subsection{Công thức quen thuộc và điều kiện sai số}
Các công thức dưới áp dụng cho giá trị hàm chính xác trên lưới đều $h>0$. $M_k$ phải là cận đúng của $|f^{(k)}|$ trên đoạn liên quan.
\begin{center}\small
\begin{tabular}{p{.2\textwidth}p{.5\textwidth}p{.2\textwidth}}
\toprule Công thức&Giá trị xấp xỉ&Cận đủ điều kiện\\\midrule
Hai điểm tiến&$[f(z+h)-f(z)]/h$&$M_2h/2$\\
Hai điểm lùi&$[f(z)-f(z-h)]/h$&$M_2h/2$\\
Ba điểm giữa, cấp 1&$[f(z+h)-f(z-h)]/(2h)$&$M_3h^2/6$\\
Ba điểm tiến, cấp 1&$[-3f(z)+4f(z+h)-f(z+2h)]/(2h)$&$M_3h^2/3$\\
Ba điểm lùi, cấp 1&$[3f(z)-4f(z-h)+f(z-2h)]/(2h)$&$M_3h^2/3$\\
Ba điểm giữa, cấp 2&$[f(z-h)-2f(z)+f(z+h)]/h^2$&$M_4h^2/12$\\\bottomrule
\end{tabular}\end{center}
Với trọng số bất kỳ và moment đúng đến bậc $p$, Taylor tại $z$ cho cận có thể thô:
\[
\left|f^{(s)}(z)-\sum_iw_if(x_i)\right|
\le\frac{M_{p+1}}{(p+1)!}\sum_i|w_i||x_i-z|^{p+1}.
\]
Nếu $|\delta y_i|\le\eta_i$, sai số do dữ liệu không vượt $\sum_i|w_i|\eta_i$. Trọng số thường tăng khi $h$ giảm, đặc biệt ở cấp đạo hàm cao. Vì thế giảm $h$ không tự bảo đảm tăng chính xác khi số liệu bị nhiễu. Không đạo hàm công thức phần dư chứa $\xi_x$ như thể $\xi_x$ là hằng số.

\section{Tích phân một chiều: khung tính chung}
\label{sec:quadrature-framework}
Mục tiêu $I=\int_a^bf(x)\,dx$. Các trường đầu vào là phương pháp; biểu thức hoặc bảng; hai cận $a,b$; $n$; bậc Newton--Cotes $d$ khi có; dung sai tự chia lưới và cận đạo hàm $M$ tùy chọn.

\begin{enumerate}
\item Chuẩn hóa cận về đoạn tăng. Nếu đảo hai cận, giữ hệ số dấu $-1$ để đổi dấu kết quả cuối. Hai cận bằng nhau cho $I=0$.
\item Chọn đúng loại nút và trọng số theo các mục tiếp theo. Đặt $h=(b-a)/n$ trên đoạn tăng.
\item Với biểu thức, tính $y_i=f(x_i)$. Với bảng, chọn các cặp trong đoạn và kiểm tra chúng khớp lưới; không chỉ kiểm tra số lượng.
\item Tính từng tích $w_iy_i$, cộng chúng rồi nhân $h$. Với bảng nhỏ, app cộng bằng Decimal của các số đã đọc; không làm tròn từng hạng để cộng tiếp.
\item Nếu có yêu cầu tinh chỉnh, làm đúng quy trình ở mục~\ref{sec:quadrature-adaptive}. Ghi rõ kết quả chính đang ở lưới nào.
\end{enumerate}

Giới hạn hiện tại là $10^6$ khoảng cho một chiều. App kiểm tra các nút vẫn phân biệt trong số float. Với bảng hữu hạn, app không cho tự tăng $n$ vì không có quyền tự sinh tung độ mới.

\section{Hình thang ghép}
\label{sec:trapezoid}
Cho $x_i=a+ih$, $i=0,\ldots,n$:
\[
T_n=h\left(\frac{y_0}{2}+\sum_{i=1}^{n-1}y_i+\frac{y_n}{2}\right).
\]
Nó tích phân từng đường thẳng nối hai mốc liên tiếp. Các bước: lấy nửa tung độ hai đầu; giữ nguyên các tung độ trong; cộng; nhân $h$. Không cần $n$ chẵn.

Ví dụ tự dựng $f(x)=x^2$, $[a,b]=[0,1]$, $n=2$, $h=1/2$:
\begin{center}\begin{tabular}{rrrr}
\toprule $x_i$&$y_i$&$w_i$&$w_iy_i$\\\midrule
0&0&$1/2$&0\\
$1/2$&$1/4$&1&$1/4$\\
1&1&$1/2$&$1/2$\\\bottomrule
\end{tabular}\end{center}
Tổng là $3/4$, nên $T_2=(1/2)(3/4)=3/8=0.375$. Giá trị đúng của ví dụ là $1/3$, sai lệch $1/24$.
Nếu $f\in C^2$, $|f''|\le M_2$, thì
\[
|I-T_n|\le\frac{(b-a)h^2}{12}M_2=\frac{(b-a)^3M_2}{12n^2}.
\]
Tăng $n$ chỉ cho hội tụ lý thuyết khi hàm/lưới phù hợp; các hàm có kỳ dị hoặc dao động chưa được lấy mẫu có thể làm ước lượng thực hành sai.

\section{Trung điểm ghép}
\label{sec:midpoint}
Các tung độ phải lấy tại tâm từng ô:
\[
\xi_i=a+(i+\tfrac12)h,\quad i=0,\ldots,n-1,
\qquad M_n=h\sum_{i=0}^{n-1}f(\xi_i).
\]
Đầu vào bảng có $n$ tâm ô, không có hai đầu mút. App kiểm tra vị trí tâm so với $a,b,n$.

Ví dụ $x^2$ trên $[0,1]$, $n=2$: hai tâm là $1/4,3/4$; các giá trị là $1/16,9/16$. Tổng $10/16$, nhân $h=1/2$ được $M_2=5/16=0.3125$.
Nếu $|f''|\le M_2^{\ast}$ thì
\[
|I-M_n|\le\frac{(b-a)h^2}{24}M_2^{\ast}.
\]
Ký hiệu $M_2^{\ast}$ là cận đạo hàm, khác ký hiệu $M_n$ của quy tắc trung điểm.

\section{Simpson ghép}
\label{sec:simpson}
Cần $n$ chẵn và $n+1$ nút cách đều, vì mỗi parabol dùng hai khoảng:
\[
S_n=\frac h3\left[y_0+4\sum_{i=1,3,\ldots,n-1}y_i
 +2\sum_{i=2,4,\ldots,n-2}y_i+y_n\right].
\]
\begin{enumerate}
\item Viết hàng trọng số nguyên $1,4,2,4,2,\ldots,4,1$.
\item Nhân từng tung độ; cộng các tích.
\item Nhân $h/3$. Nếu dùng trọng số app đã lưu là $1/3,4/3,2/3,\ldots$, chỉ nhân $h$, không chia 3 lần nữa.
\end{enumerate}
Ví dụ $x^2$ trên $[0,1]$, $n=2$: $S_2=(1/6)[0+4(1/4)+1]=1/3$. Trường hợp này đúng chính xác vì Simpson tích phân đúng đa thức đến bậc ba. Nếu $f\in C^4$, $|f^{(4)}|\le M_4$:
\[
|I-S_n|\le\frac{(b-a)h^4}{180}M_4.
\]
Không áp công thức Simpson này cho $n$ lẻ hoặc mốc không đều. Chọn một bậc nội suy khác không tự biến bảng không đều thành lưới Simpson.

\section{Newton--Cotes đóng bậc 1 đến 10}
\label{sec:newton-cotes}
\subsection{Công thức cho một khối và cách sinh trọng số}
Một khối bậc $d$ gồm $d$ khoảng, $d+1$ mốc đều. Đặt $x=x_0+ht$, $0\le t\le d$. Nội suy Lagrange theo các nút $0,1,\ldots,d$ rồi tích phân:
\[
I_{\rm khối}\approx h\sum_{j=0}^{d}w_j f(x_0+jh),\qquad
w_j=\int_0^d\prod_{k\ne j}\frac{t-k}{j-k}\,dt.
\]
Một cách kiểm tra độc lập là hệ moment
\[
\sum_{j=0}^{d}w_jj^r=\frac{d^{r+1}}{r+1},\qquad r=0,\ldots,d.
\]
App lấy trọng số chuẩn bằng \texttt{scipy.integrate.newton\_cotes(d,equal=1)}. Bảng dưới được kiểm lại bằng hệ moment hữu tỷ, không lấy số hiển thị làm chuẩn.
{\small\begin{longtable}{rcll}
\toprule
$d$ & $p$ & Hệ số chung $\kappa_d$ & Vector nguyên $v$, với $w_j=\kappa_d v_j$\\
\midrule
\endhead
1&2&$\frac{1}{2}$&\begin{minipage}[t]{.55\textwidth}\raggedright\small $(1,\allowbreak 1)$\end{minipage}\\[4pt]
2&4&$\frac{1}{3}$&\begin{minipage}[t]{.55\textwidth}\raggedright\small $(1,\allowbreak 4,\allowbreak 1)$\end{minipage}\\[4pt]
3&4&$\frac{3}{8}$&\begin{minipage}[t]{.55\textwidth}\raggedright\small $(1,\allowbreak 3,\allowbreak 3,\allowbreak 1)$\end{minipage}\\[4pt]
4&6&$\frac{2}{45}$&\begin{minipage}[t]{.55\textwidth}\raggedright\small $(7,\allowbreak 32,\allowbreak 12,\allowbreak 32,\allowbreak 7)$\end{minipage}\\[4pt]
5&6&$\frac{5}{288}$&\begin{minipage}[t]{.55\textwidth}\raggedright\small $(19,\allowbreak 75,\allowbreak 50,\allowbreak 50,\allowbreak 75,\allowbreak 19)$\end{minipage}\\[4pt]
6&8&$\frac{1}{140}$&\begin{minipage}[t]{.55\textwidth}\raggedright\small $(41,\allowbreak 216,\allowbreak 27,\allowbreak 272,\allowbreak 27,\allowbreak 216,\allowbreak 41)$\end{minipage}\\[4pt]
7&8&$\frac{7}{17280}$&\begin{minipage}[t]{.55\textwidth}\raggedright\small $(751,\allowbreak 3577,\allowbreak 1323,\allowbreak 2989,\allowbreak 2989,\allowbreak 1323,\allowbreak 3577,\allowbreak 751)$\end{minipage}\\[4pt]
8&10&$\frac{4}{14175}$&\begin{minipage}[t]{.55\textwidth}\raggedright\small $(989,\allowbreak 5888,\allowbreak -928,\allowbreak 10496,\allowbreak -4540,\allowbreak 10496,\allowbreak -928,\allowbreak 5888,\allowbreak 989)$\end{minipage}\\[4pt]
9&10&$\frac{9}{89600}$&\begin{minipage}[t]{.55\textwidth}\raggedright\small $(2857,\allowbreak 15741,\allowbreak 1080,\allowbreak 19344,\allowbreak 5778,\allowbreak 5778,\allowbreak 19344,\allowbreak 1080,\allowbreak 15741,\allowbreak 2857)$\end{minipage}\\[4pt]
10&12&$\frac{5}{299376}$&\begin{minipage}[t]{.55\textwidth}\raggedright\small $(16067,\allowbreak 106300,\allowbreak -48525,\allowbreak 272400,\allowbreak -260550,\allowbreak 427368,\allowbreak -260550,\allowbreak 272400,\allowbreak -48525,\allowbreak 106300,\allowbreak 16067)$\end{minipage}\\[4pt]
\bottomrule
\end{longtable}
}

Ở đây $p$ là bậc sai số toàn cục dùng trong Runge: $p=d+1$ nếu $d$ lẻ, $p=d+2$ nếu $d$ chẵn. Độ chính xác đa thức của quy tắc là đến bậc $p-1$. Sự đối xứng tạo thêm một bậc đúng khi $d$ chẵn.

\subsection{Ghép nhiều khối}
Đầu vào app cần $1\le d\le10$ và $n$ chia hết cho $d$. Chia lưới thành các khối bắt đầu ở $i=0,d,2d,\ldots,n-d$:
\begin{enumerate}
\item Khởi tạo trọng số toàn lưới $W_0=\cdots=W_n=0$.
\item Với mỗi khối bắt đầu i, cộng $w_j$ vào $W_{i+j}$, $j=0,\ldots,d$.
\item Các mốc chung giữa hai khối nhận tổng hai đóng góp.
\item Tính $I_n=h\sum_{i=0}^n W_i y_i$.
\end{enumerate}
Ví dụ $d=2,n=4$ cho trọng số $(1,4,2,4,1)/3$, chính là Simpson ghép.

\subsection{Ví dụ Boole, $d=4$}
Với $f=x^2$ trên $[0,1]$, $n=d=4$, $h=1/4$:
\[
I\approx\frac{2h}{45}[7y_0+32y_1+12y_2+32y_3+7y_4].
\]
\begin{center}\begin{tabular}{rrrr}
\toprule $i$&$y_i$&Trọng số nguyên&Đóng góp trong ngoặc\\\midrule
0&0&7&0\\
1&$1/16$&32&2\\
2&$1/4$&12&3\\
3&$9/16$&32&18\\
4&1&7&7\\\bottomrule
\end{tabular}\end{center}
Tổng $30$, hệ số ngoài $2h/45=1/90$, kết quả $1/3$. Với $d=1,2,3$ lần lượt có hình thang, Simpson $1/3$, Simpson $3/8$; các bậc còn lại dùng cùng quy trình nhưng trọng số khác.

Một số bậc cao có trọng số âm. Chúng có thể khuếch đại sai số dữ liệu; tăng bậc không phải bảo đảm tốt hơn. App không tự thay bậc hay cắt khối lẻ ở cuối.

\section{Cận lý thuyết, Runge và tự chia lưới}
\label{sec:quadrature-adaptive}
\subsection{Ba chế độ khác nhau}
\begin{center}\begin{tabular}{p{.24\textwidth}p{.63\textwidth}}
\toprule Cấu hình&App thực hiện và cách đọc kết quả\\\midrule
Không nhập dung sai&Giữ n đã chọn; tính thêm lưới 2n nếu được để đối chiếu. Kết quả chính vẫn là $I_n$. Runge in ra cho sai số lưới $2n$.\\
Nhập dung sai, không có M&Liên tiếp tăng lưới và dùng Runge, kèm một lưới kiểm tra không lồng. Kết quả chính là lưới cuối thực sự đạt tiêu chuẩn thực hành.\\
Nhập dung sai và M&Chọn n nhỏ nhất theo cận $C/n^p\le\varepsilon$, tôn trọng số khoảng mỗi khối. Giá trị chính dùng n đó.\\\bottomrule
\end{tabular}\end{center}

Nếu $I_h=I+Ch^p+o(h^p)$, hai lưới cách nhau hệ số 2 cho
\[
E_{h/2}\approx\frac{|I_{h/2}-I_h|}{2^p-1}.
\]
Đây là ước lượng tiệm cận khi đã ở miền sai số có quy luật đó. Một hiệu bằng 0 cũng có thể do bỏ sót dao động; nó không tự là chứng minh $I_h=I$.

\subsection{Đúng thứ tự vòng tự chia lưới trong code}
\begin{enumerate}
\item Tính $I_n$ ở n đầu vào.
\item Tính $I_{2n}$ và $E=|I_{2n}-I_n|/(2^p-1)$. Đổi n thành $2n$, giữ giá trị mới.
\item Nếu $E>\varepsilon$, lặp lại bước 2.
\item Nếu $E\le\varepsilon$, tạo thêm lưới $n+q$, với $q=1$ cho hình thang/trung điểm, $q=2$ cho Simpson, $q=d$ cho Newton--Cotes. Gần giới hạn tài nguyên có thể dùng $n-q$.
\item Tính $I_{n+q}$. Nếu $|I_{n+q}-I_n|\le\varepsilon$, nhận $I_n$; nếu chưa đạt, chuyển điểm xuất phát sang lưới kiểm tra và tiếp tục.
\end{enumerate}
App cho tối đa 18 vòng ở một chiều và không vượt $10^6$ khoảng. Hết giới hạn trả lỗi chưa đạt, không nhận im giá trị cuối. Kiểm tra thêm lưới không lồng giúp giảm bỏ sót dao động nhưng vẫn không là cận bảo đảm.

\subsection{Ví dụ đầy đủ: hình thang tự chia cho $x^2$}
Chọn $[0,1]$, n ban đầu 2, $\varepsilon=10^{-3}$, không cho M. Công thức chính xác của hình thang trên hàm này là $T_n=1/3+1/(6n^2)$.
\begin{center}\begin{tabular}{rrrl}
\toprule n mới&$T_n$&Ước lượng Runge cho n mới&Quyết định\\\midrule
2&0.375&---&Bước xuất phát\\
4&0.34375&0.010416666667&Chia tiếp\\
8&0.3359375&0.002604166667&Chia tiếp\\
16&0.333984375&0.000651041667&Thử lưới 17\\\bottomrule
\end{tabular}\end{center}
Lưới 17 cho $T_{17}=0.333910034602$; chênh với lưới 16 khoảng $7.43404\cdot10^{-5}<10^{-3}$. App nhận $T_{16}=0.333984375$. Giá trị được nhận là lưới 16, không phải lưới kiểm tra 17.

Nếu bổ sung $M_2=2$, cận hình thang là $1/(6n^2)$. Điều kiện $1/(6n^2)\le10^{-3}$ cho n nhỏ nhất bằng 13. App chọn n=13 và nhận $T_{13}\approx0.334319526627$, với cận $1/1014\approx0.000986193294$. Đây là chế độ có cận nếu $M_2$ đúng.

\subsection{Bảng dữ liệu và những lưới không được suy diễn}
Với bảng ở hai đầu đoạn, app có thể lấy mỗi hai nút để tạo lưới thô nếu số khoảng phù hợp quy tắc. Khi đó Runge so giá trị bảng hiện tại với bảng thô và áp dụng cho lưới hiện tại. Trung điểm không tự có tâm của lưới thô trong tập tâm lưới mịn nên code không làm phép đối chiếu đó.

Không được nội suy thêm các y mới rồi coi chúng như phép đo độc lập để tinh chỉnh tích phân của bảng. Nếu tự quyết định dùng một mô hình nội suy trước khi tích phân, đó là một bước mô hình hóa khác cần nêu trong bài.

\section{Tích phân kép trên miền chữ nhật}
\label{sec:double-integral}
App nhận biểu thức $f(x,y)$ và hình chữ nhật $[a,b]\times[c,d]$, số khoảng n cho \emph{mỗi chiều}, chọn hình thang hoặc Simpson, dung sai Runge. Không có tùy chọn miền tam giác, miền cong hoặc bảng hai chiều trong màn hình này.
\[
h_x=\frac{b-a}{n},\quad h_y=\frac{d-c}{n},\qquad
I_n=h_xh_y\sum_{i=0}^{n}\sum_{j=0}^{n}w_iw_j f(x_i,y_j).
\]
Với Simpson $w=(1,4,2,\ldots,4,1)/3$ và n chẵn. Với hình thang hai đầu có trọng số $1/2$, các nút trong có trọng số 1.
\begin{enumerate}
\item Chuẩn hóa từng cặp cận; đảo một chiều đổi dấu, đảo cả hai giữ dấu.
\item Lập hai dãy nút và cùng một vector trọng số cho mỗi chiều.
\item Tính các $f_{ij}$ và từng tích $w_iw_jf_{ij}$.
\item Cộng hai chiều rồi nhân $h_xh_y$.
\item Tăng gấp đôi n ở cả hai chiều; so Runge với bậc 2 hoặc 4; khi đủ nhỏ kiểm thêm lưới n+1 hoặc n+2 không lồng.
\end{enumerate}
Giới hạn hiện tại: không quá $2000\times2000$ khoảng và tối đa 12 vòng tinh chỉnh. Miền có chiều dài 0 trả tích phân 0. Một dung sai chung không phải cận tổng hợp đã được chứng minh của hai chiều.

Ví dụ tự dựng $f=x^2+y^2$ trên $[0,1]^2$, Simpson n=2, $h_x=h_y=1/2$:
\begin{center}\begin{tabular}{r|rrr}
\toprule $x\backslash y$&0&$1/2$&1\\\midrule
0&0&$1/4$&1\\
$1/2$&$1/4$&$1/2$&$5/4$\\
1&1&$5/4$&2\\\bottomrule
\end{tabular}\end{center}
Nhân theo vector nguyên $(1,4,1)$ ở từng chiều: tổng có trọng số nguyên là 24; hệ số ngoài $h_xh_y/9=1/36$ nên $I_2=2/3$. Tinh chỉnh n=4 và kiểm n=6 đều cho $2/3$, nên app nhận lưới 4 trong chế độ tự kiểm. Đây là ví dụ đúng với đa thức đang dùng; không suy kết quả luôn chính xác cho mọi hàm.

\section{Checklist chép lời giải và các ca dừng}
\begin{enumerate}
\item Ghi dữ liệu là biểu thức hay bảng; miền và thứ tự cận.
\item Ghi số khoảng, số nút, h; với trung điểm phải ghi vị trí tâm.
\item Viết trọng số theo đúng quy ước: đã chứa $1/3$, $3/8$ hay còn để ngoài.
\item Bảng có các giá trị, trọng số, đóng góp và tổng. Với đạo hàm ghi cả tập mốc và cách tạo trọng số.
\item Nếu có tinh chỉnh, ghi giá trị trên mỗi lưới, ước lượng cho lưới nào, kiểm tra không lồng và lý do dừng.
\item Chỉ ghi cận bảo đảm khi đã có giả thiết/cận đạo hàm đúng. Trạng thái đạt Runge là kết luận thực hành.
\end{enumerate}
Các lỗi cần xử lý: mốc trùng/không hữu hạn; bảng không khớp lưới; Simpson n lẻ; Newton--Cotes n không chia hết d; thiếu tung độ; hàm không thực tại nút; bước lưới không phân biệt được trong float; kết quả vượt miền số; chưa đạt dung sai trong giới hạn tài nguyên.

\paragraph{Nguồn đối chiếu trong repo.} Các hàm \texttt{differentiate}, \texttt{quadrature\_weights}, \texttt{integrate}, \texttt{double\_integral} ở \texttt{pps/calculus.py}; trường nhập tại \texttt{pps/gui.py}.

Số liệu kiểm tra được tái tạo bằng script trong thư mục \texttt{verification/}: \texttt{calculus\_verification.py}. Kết quả lưu trong \texttt{calculus\_examples.json}.

\chapter{Bài toán Cauchy: Euler, Runge--Kutta và Adams}
\label{ch:ode}
\section{Input và output dùng chung cho các tool ODE}
\subsection*{Input: các ô đúng tên trên giao diện}
Bảng này được lấy từ biểu mẫu của tool. Giá trị mặc định là ví dụ nhập, cần thay bằng dữ liệu đề. Hai dãy $x,y$ là bảng dữ liệu riêng bên dưới nếu chức năng dùng bảng.
\subsubsection*{Ô chung của các lựa chọn trong nhóm}
\begin{center}\small\begin{tabular}{@{}p{.32\linewidth}p{.17\linewidth}p{.43\linewidth}@{}}\toprule
Ô nhập & Khóa JSON & Cách nhập / mặc định\\\midrule
Vế phải x′;\allowbreak{} y′;\allowbreak{} z′ (1–3 hàm) & \texttt{expr} & \texttt{x}
\par Biến độc lập t;\allowbreak{} một phương trình x′=x nhập x. Hệ: y;\allowbreak{} -x\\[3pt]
Giá trị đầu x,\allowbreak{} y,\allowbreak{} z & \texttt{initial} & \texttt{1}\\[3pt]
Thời điểm đầu t0 & \texttt{a} & \texttt{0}\\[3pt]
Thời điểm cuối & \texttt{b} & \texttt{1}\\[3pt]
Bước h & \texttt{h} & \texttt{0.1}\\[3pt]
Dung sai bước ẩn & \texttt{tol} & \texttt{1e-10}\\[3pt]
\bottomrule\end{tabular}\end{center}
\subsubsection*{Ô riêng: RK2 tham số}
\begin{center}\small\begin{tabular}{@{}p{.32\linewidth}p{.17\linewidth}p{.43\linewidth}@{}}\toprule
Ô nhập & Khóa JSON & Cách nhập / mặc định\\\midrule
Tham số biết trước & \texttt{parameter} & \texttt{alpha2}
\par alpha2,\allowbreak{} beta11,\allowbreak{} r1 hoặc r2.\\[3pt]
Giá trị tham số & \texttt{value} & \texttt{1}\\[3pt]
\bottomrule\end{tabular}\end{center}
\subsubsection*{Ô riêng: RK3 tham số}
\begin{center}\small\begin{tabular}{@{}p{.32\linewidth}p{.17\linewidth}p{.43\linewidth}@{}}\toprule
Ô nhập & Khóa JSON & Cách nhập / mặc định\\\midrule
alpha2 & \texttt{alpha2} & \texttt{0.5}\\[3pt]
alpha3 & \texttt{alpha3} & \texttt{1}\\[3pt]
\bottomrule\end{tabular}\end{center}
\subsubsection*{Ô riêng: RK tùy chỉnh}
\begin{center}\small\begin{tabular}{@{}p{.32\linewidth}p{.17\linewidth}p{.43\linewidth}@{}}\toprule
Ô nhập & Khóa JSON & Cách nhập / mặc định\\\midrule
Ma trận A (JSON) & \texttt{rk\_a} & \texttt{[[0,\allowbreak{}0],\allowbreak{}[1,\allowbreak{}0]]}\\[3pt]
Trọng số b & \texttt{rk\_b} & \texttt{0.5 0.5}\\[3pt]
Vị trí c & \texttt{rk\_c} & \texttt{0 1}
\par Cho RK2/\allowbreak{}RK3 hoặc RK hiện nhiều tầng;\allowbreak{} A tam giác dưới.\\[3pt]
\bottomrule\end{tabular}\end{center}

\paragraph{Đổi đề sang tên biến của tool.} Biến độc lập là $t$; các thành phần trạng thái lần lượt $x,y,z$. Ví dụ đề $y'=y$ thì nhập biểu thức \texttt{x}, điều kiện đầu một số; đề hệ dao động nhập \texttt{y; -x}, đầu \texttt{1 0}. Nếu dùng tên không có trong hệ, tool từ chối.
\paragraph{Input lưới.} $a$ là thời điểm đầu, $b$ là cuối, $h$ cùng chiều $a\to b$ và chia hết đoạn trong dung sai máy. $u_0$ là vector đầu. Dung sai bước ẩn không phải sai số toàn cục ODE; số bước được xác định từ $a,b,h$.
\paragraph{Output chung.} Bảng các thời điểm $t_n$ và trạng thái $u_n$; bảng hệ số RK nếu có; phần dư cuối mỗi bước ẩn và đồ thị. Tool hiện không lưu mọi tầng RK hoặc mọi vòng nội bộ của bộ giải ẩn; các khuôn dưới đây hướng dẫn cách tự lập bài tính tay theo đúng công thức.

\subsection{Euler hiện, Euler ẩn, Hình thang: khuôn một bước}
Đặt $u'=F(t,u)$, $F_n=F(t_n,u_n)$. Mỗi dòng nhận $t_n,u_n$, tạo $t_{n+1}=t_n+h$ rồi xử lý:
\begin{center}\small\begin{tabular}{@{}p{.21\linewidth}p{.43\linewidth}p{.27\linewidth}@{}}\toprule
Tool&Quá trình trong bước&Giá trị chuyển tiếp\\\midrule
Euler hiện&Tính $F_n$, nhân $h$, cộng vào $u_n$&$u_{n+1}=u_n+hF_n$\\
Euler ẩn&Dự đoán $u_n+hF_n$; giải $R(v)=v-u_n-hF(t_{n+1},v)=0$&Nghiệm mới và phần dư\\
Hình thang&Dự đoán như trên; giải $R(v)=v-u_n-\frac h2[F_n+F(t_{n+1},v)]=0$&Nghiệm mới và phần dư\\\bottomrule
\end{tabular}\end{center}
\begin{center}\small\begin{tabular}{cllll}\toprule
Bước&Nhận từ trước&Độ dốc cũ&Phép tính/ẩn phải giải&Output hàng\\\midrule
0&$t_0,u_0$&$F(t_0,u_0)$&Công thức của tool chọn&$t_1,u_1,\|R(u_1)\|$ nếu ẩn\\
$n$&$t_n,u_n$&$F(t_n,u_n)$&Cùng quy trình trên&$t_{n+1},u_{n+1}$\\\bottomrule
\end{tabular}\end{center}
Không dùng $u_{n+1}$ của một thành phần để cập nhật tùy tiện thành phần còn lại trong cùng bước hệ vector.

\subsection{RK2, RK3, RK4: bảng tầng tổng quát thay cho bảng số riêng}
Mỗi $K_j$ dưới đây là \emph{độ dốc}, chưa nhân $h$. Mọi tầng của một bước đều xuất phát từ cùng $t_n,u_n$.
\begin{center}\small\begin{tabular}{clll}\toprule
RK2&Thời điểm&Trạng thái thế vào F&Độ dốc\\\midrule
1&$t_n$&$u_n$&$K_1=F(t_n,u_n)$\\
2&$t_n+h$&$u_n+hK_1$&$K_2=F(t_n+h,u_n+hK_1)$\\\bottomrule
\end{tabular}\end{center}
Output bước: $u_{n+1}=u_n+\frac h2(K_1+K_2)$.
\begin{center}\small\begin{tabular}{clll}\toprule
RK3&Thời điểm&Trạng thái thế vào F&Độ dốc\\\midrule
1&$t_n$&$u_n$&$K_1$\\
2&$t_n+h/2$&$u_n+hK_1/2$&$K_2$\\
3&$t_n+h$&$u_n-hK_1+2hK_2$&$K_3$\\\bottomrule
\end{tabular}\end{center}
Output bước: $u_{n+1}=u_n+\frac h6(K_1+4K_2+K_3)$.
\begin{center}\small\begin{tabular}{clll}\toprule
RK4&Thời điểm&Trạng thái thế vào F&Độ dốc\\\midrule
1&$t_n$&$u_n$&$K_1$\\
2&$t_n+h/2$&$u_n+hK_1/2$&$K_2$\\
3&$t_n+h/2$&$u_n+hK_2/2$&$K_3$\\
4&$t_n+h$&$u_n+hK_3$&$K_4$\\\bottomrule
\end{tabular}\end{center}
Output bước: $u_{n+1}=u_n+\frac h6(K_1+2K_2+2K_3+K_4)$. Sau khi có $u_{n+1}$ mới bắt đầu bảng tầng của bước sau.

\subsection{RK2 tham số, RK3 tham số và RK tùy chỉnh}
\textbf{Input riêng}: RK2 chọn tham số biết trước \texttt{alpha2}, \texttt{beta11}, \texttt{r1} hoặc \texttt{r2} và giá trị; RK3 nhập \texttt{alpha2,alpha3}; tùy chỉnh nhập $A,b,c$. Các mục sau đưa công thức chuyển tham số thành bảng hệ số và các nhánh ngoại lệ.
\textbf{Quá trình}: tạo và kiểm bảng $\to$ với mỗi tầng $j$, điền hàng tổng quát
\[
T_j=t_n+c_jh,\quad
U_j=u_n+h(a_{j1}K_1+\cdots+a_{j,j-1}K_{j-1}),\quad K_j=F(T_j,U_j).
\]
Cuối bước $u_{n+1}=u_n+h(b_1K_1+b_2K_2+\cdots+b_sK_s)$. Output gồm bảng hệ số thực sự dùng và nghiệm trên lưới. Kiểm cấu trúc RK hiện không tự chứng minh bậc cao của một bảng tùy ý.

\subsection{Adams 2, 3, 4, 5: bảng lịch sử và thứ tự cập nhật}
\textbf{Input riêng}: chọn số $s$ qua tên Adams $s$, không nhập thêm $s$ giá trị đầu; tool dùng $s-1$ bước RK4 khởi tạo. Nếu đoạn ngắn thì có thể chỉ chạy các bước khởi tạo.
\begin{center}\small\begin{tabular}{cllll}\toprule
Vị trí lịch sử&Thời điểm&Trạng thái&Độ dốc&Đóng góp AB / AM cũ\\\midrule
0 mới nhất&$t_n$&$u_n$&$F_n$&$\beta_0F_n$ / $a_1F_n$\\
$j$&$t_{n-j}$&$u_{n-j}$&$F_{n-j}$&$\beta_jF_{n-j}$ / $a_{j+1}F_{n-j}$\\
$s-1$ cũ nhất&$t_{n-s+1}$&$u_{n-s+1}$&$F_{n-s+1}$&Đóng góp cuối\\\bottomrule
\end{tabular}\end{center}
Dự báo của từng tool, viết rõ các phép nhân:
\begin{align*}
\mathrm{AB2}:\quad u^P&=u_n+\tfrac h2(3F_n-F_{n-1}),\\
\mathrm{AB3}:\quad u^P&=u_n+\tfrac h{12}(23F_n-16F_{n-1}+5F_{n-2}),\\
\mathrm{AB4}:\quad u^P&=u_n+\tfrac h{24}(55F_n-59F_{n-1}+37F_{n-2}-9F_{n-3}),\\
\mathrm{AB5}:\quad u^P&=u_n+\tfrac h{720}(1901F_n-2774F_{n-1}+2616F_{n-2}\\
&\hspace{38mm}-1274F_{n-3}+251F_{n-4}).
\end{align*}
Sau dự báo, tính phần đuôi $T=a_1F_n+a_2F_{n-1}+\cdots+a_sF_{n-s+1}$; giải phương trình $v=u_n+h[a_0F(t_{n+1},v)+T]$ từ điểm đầu $u^P$. Tool giải bước hiệu chỉnh ẩn, không chỉ thế dự báo vào một lần. Output bước là $u_{n+1}$ đã hiệu chỉnh và phần dư. Đưa giá trị này vào đầu lịch sử, loại giá trị cũ nhất rồi tiếp tục. Các vector AM đúng theo từng $s$ và ví dụ mọi bước khởi tạo được cho bên dưới.


\section{Chuyển đề về đúng tên biến của tool}
App giải hệ
\[
\mathbf u'(t)=\mathbf F(t,\mathbf u(t)),\qquad
\mathbf u(t_0)=\mathbf u_0,
\]
với một đến ba thành phần. Trên giao diện, biến độc lập luôn là \texttt{t}; các thành phần lần lượt là \texttt{x}, \texttt{y}, \texttt{z}. Vì vậy, nếu sách viết $y'=t+y$, $y(0)=1$, trong tool nhập vế phải \texttt{t+x} và giá trị đầu \texttt{1}. Không nhập \texttt{t+y} cho một phương trình một ẩn vì khi đó \texttt{y} không phải thành phần đã khai báo.

\begin{center}\begin{tabular}{p{.3\textwidth}p{.55\textwidth}}
\toprule Trường nhập&Ý nghĩa\\\midrule
Vế phải&1--3 biểu thức, ngăn bằng dấu chấm phẩy; thứ tự trùng thứ tự giá trị đầu.\\
Giá trị đầu&Vector $\mathbf u_0$ gồm đúng số thành phần.\\
Thời điểm đầu/cuối&$a=t_0,b$; app dùng lưới đều cố định.\\
Bước h&Phải cùng chiều a đến b; $(b-a)/h$ phải là số nguyên trong dung sai lưới.\\
Dung sai bước ẩn&Điều khiển giải hệ đại số của Euler ẩn, hình thang và Adams; không tự chọn h.\\
Tham số/bảng RK&Có thêm ở các lựa chọn tham số và RK tùy chỉnh.\\\bottomrule
\end{tabular}\end{center}

App cho tối đa 20.000 khoảng thời gian. Nếu a=b hữu hạn và h hữu hạn khác 0, kết quả chỉ gồm điều kiện đầu. Nếu b<a có thể nhập h âm. Không tự thêm một bước cuối ngắn hơn để bù độ dài đoạn; hãy chọn h sao cho lưới chạm đúng b. Mọi giá trị hàm ở các tầng phải hữu hạn trong miền thực.

\section{Ba khái niệm phải tách biệt}
\begin{enumerate}
\item \textbf{Tính từng bước:} từ $\mathbf u_n$ tìm $\mathbf u_{n+1}$. Euler/RK hiện dùng một số phép tính hữu hạn, không có vòng hội tụ nội bộ.
\item \textbf{Giải bước ẩn:} một bước có thể chứa $\mathbf u_{n+1}$ ở cả hai vế. Bộ giải đại số phải tìm nghiệm của phương trình đó. Dung sai ở đây kiểm tra phần dư phương trình bước.
\item \textbf{Hội tụ khi h tiến về 0:} nghiệm trên lưới tiến về nghiệm liên tục, cần các giả thiết về hàm, tính nhất quán, ổn định và bậc của phương pháp. Giải bước ẩn rất chính xác không tự làm sai số toàn cục nhỏ.
\end{enumerate}
Với $\mathbf F$ Lipschitz theo trạng thái, bài toán Cauchy có nghiệm duy nhất cục bộ theo các giả thiết thông thường. Các kết quả bậc sai số dưới đây giả sử nghiệm đủ trơn và các bước nằm trong miền ổn định; chúng không phải cận số được tự chứng minh từ dữ liệu đã nhập.

\section{Euler hiện}
\label{sec:ode-euler-explicit}
\subsection{Công thức và thứ tự thao tác}
\[
\mathbf k_1=\mathbf F(t_n,\mathbf u_n),\qquad
\mathbf u_{n+1}=\mathbf u_n+h\mathbf k_1.
\]
\begin{enumerate}
\item Lấy $t_n,\mathbf u_n$ từ hàng nghiệm đã có.
\item Thế vào vế phải để tính vector độ dốc $\mathbf k_1$.
\item Nhân cả vector với h rồi cộng vào trạng thái cũ.
\item Ghi hàng mới tại $t_{n+1}=t_n+h$; dùng hàng này cho bước sau.
\end{enumerate}
Đây là tiếp tuyến tại đầu trái của bước. Sai số toàn cục điển hình là $O(h)$.

Ví dụ tự dựng $u'=u,u(0)=1,h=0.1$, tương ứng nhập \texttt{x}. Hai bước đầy đủ:
\begin{center}\begin{tabular}{rrrrl}
\toprule n&$t_n$&$u_n$&$k_1$&Phép cập nhật\\\midrule
0&0&1&1&$u_1=1+0.1(1)=1.1$\\
1&0.1&1.1&1.1&$u_2=1.1+0.1(1.1)=1.21$\\\bottomrule
\end{tabular}\end{center}
Không có epsilon để dừng hai phép tính này. Với phương trình thử $u'=\lambda u$, hệ số tăng trưởng là $R(z)=1+z$, $z=h\lambda$; ổn định tuyệt đối đòi $|1+z|\le1$. Vì vậy một h lớn có thể làm nghiệm số tăng dù nghiệm thật đang giảm.

\section{Euler ẩn}
\label{sec:ode-euler-implicit}
\subsection{Từ phương trình vi phân đến hệ đại số ở mỗi bước}
\[
\mathbf u_{n+1}=\mathbf u_n+h\mathbf F(t_{n+1},\mathbf u_{n+1}).
\]
Đặt ẩn $\mathbf v=\mathbf u_{n+1}$ và phần dư
\[
\mathbf R(\mathbf v)=\mathbf v-\mathbf u_n-h\mathbf F(t_{n+1},\mathbf v),
\quad D\mathbf R=I-hJ_F(t_{n+1},\mathbf v).
\]
\begin{enumerate}
\item Tính $\mathbf F(t_n,\mathbf u_n)$ để có dự đoán Euler hiện $\mathbf v_0=\mathbf u_n+h\mathbf F(t_n,\mathbf u_n)$.
\item Lập đúng $\mathbf R$ và Jacobian ở trên. App lấy Jacobian của vế phải bằng đạo hàm ký hiệu.
\item Gọi \texttt{scipy.optimize.root} với dự đoán đó, Jacobian và dung sai đã nhập.
\item Kiểm lại phần dư tại nghiệm cuối và tính không suy biến của Jacobian bước; chỉ nhận khi các kiểm tra đạt.
\end{enumerate}
App không lưu đầy đủ các lần gọi nội bộ của \texttt{root}. Bảng phần dư là phần dư cuối mỗi bước thời gian, không phải một nhật ký Newton. Không nên tự gắn các vòng Newton tính tay vào đó rồi nói tool đã chạy đúng các vòng ấy.

Ví dụ $u'=u$, $u_0=1$, $h=0.1$: dự đoán là 1.1; phương trình cần giải là $v=1+0.1v$, hay $0.9v=1$. Do đó
\[
u_1=\frac{10}{9}=1.111111111111\ldots,\qquad
R(u_1)=0,\quad R'(v)=0.9.
\]
Đây là giải đại số chính xác của ví dụ tuyến tính; phép gọi thư viện được đối chiếu với kết quả này.

Nếu dùng lặp đơn $v^{(k+1)}=u_n+hF(t_{n+1},v^{(k)})$ thì điều kiện đủ $|h|L<1$ giúp co trên đoạn bất biến thích hợp. \emph{Đó không phải thuật toán lặp đơn mà code hiện dùng.} Bộ giải root có thể hội tụ trong những trường hợp lặp đơn không co, và cũng có thể thất bại ở bài phi tuyến khó.

Euler ẩn có bậc toàn cục 1. Với phương trình thử, $R(z)=1/(1-z)$, ổn định trên nửa mặt phẳng $\operatorname{Re}z\le0$; điều này không thay thế việc tìm đúng nghiệm bước phi tuyến.

\section{Hình thang cho ODE}
\label{sec:ode-trapezoid}
Đây là phương pháp hình thang \emph{ẩn} cho phương trình vi phân:
\[
\mathbf u_{n+1}=\mathbf u_n+\frac h2\left[
\mathbf F(t_n,\mathbf u_n)+\mathbf F(t_{n+1},\mathbf u_{n+1})\right].
\]
Đặt $\theta=1/2$; Euler ẩn tương ứng $\theta=1$. Hai lựa chọn dùng cùng dạng
\[
\mathbf R(\mathbf v)=\mathbf v-\mathbf u_n-h[(1-\theta)\mathbf F_n+
\theta\mathbf F(t_{n+1},\mathbf v)],\quad
D\mathbf R=I-h\theta J_F.
\]
Các bước nhập, dự đoán, giải và kiểm tra như Euler ẩn; chỉ thay $\theta$. Với $u'=u$, $u_0=1,h=0.1$:
\[
v=1+0.05(1+v),\quad0.95v=1.05,
\quad u_1=\frac{21}{19}=1.105263157895\ldots.
\]
Bậc toàn cục là 2 dưới giả thiết đủ trơn và ổn định. Hệ số thử $R(z)=(1+z/2)/(1-z/2)$; ổn định trên nửa mặt phẳng trái nhưng không triệt tiêu mạnh các thành phần rất cứng như Euler ẩn.

\subsection{Chính xác app kiểm tra gì khi nhận một bước ẩn?}
Đặt $\tau$ là dung sai đã nhập. Tại nghiệm cuối $\widehat{\mathbf v}$, app yêu cầu các số hữu hạn và
\[
\|\mathbf R(\widehat{\mathbf v})\|_\infty
\le\max(\tau,32\epsilon_{\rm mach})\left(1+\|\widehat{\mathbf v}\|_\infty\right).
\]
Tiếp theo, nếu các giá trị kỳ dị của $D\mathbf R$ là $\sigma_{\max},\sigma_{\min}$, app từ chối khi
\[
\sigma_{\min}\le64\epsilon_{\rm mach}\max(1,\sigma_{\max}).
\]
Đây là kiểm tra phần dư và tính khả nghịch số của bước, không là chứng nhận nghiệm duy nhất toàn cục của bài phi tuyến. Không đạt thì báo giảm h hoặc kiểm tra miền hàm; không âm thầm chuyển sang Euler hiện.

\section{Khung Runge--Kutta hiện tổng quát}
\label{sec:ode-rk-general}
Một bảng Butcher gồm ma trận A, vector b và c, với s tầng:
\[
\mathbf K_i=\mathbf F\left(t_n+c_i h,
\mathbf u_n+h\sum_{j<i}a_{ij}\mathbf K_j\right),\quad
\mathbf u_{n+1}=\mathbf u_n+h\sum_{i=1}^{s}b_i\mathbf K_i.
\]
Trong quy ước của app, $K_i$ là \emph{độ dốc}, chưa nhân h. Nếu dùng sách định nghĩa $k_i=hF(\cdots)$ thì phải đổi quy ước nhất quán, không nhân h hai lần.
\begin{enumerate}
\item Tầng 1: tổng các tầng trước rỗng nên chỉ dùng $t_n+c_1h$ và $\mathbf u_n$.
\item Tầng i: tính trạng thái trung gian từ các $K_j$ đã có, rồi thế vào F để lấy $K_i$ mới.
\item Khi đủ s tầng, nhân từng $K_i$ với $b_i$, cộng, nhân h rồi cộng trạng thái cũ.
\item Ghi trạng thái mới và dùng nó cho bước kế tiếp; không dùng trạng thái tầng cuối thay cho tổng trọng số.
\end{enumerate}
Không có phép lặp hội tụ giữa các tầng RK hiện. Số tầng là s, được xác định bởi bảng hệ số.

\section{RK2 mặc định: Heun hai tầng}
\label{sec:ode-rk2}
\[
K_1=F(t_n,u_n),\quad K_2=F(t_n+h,u_n+hK_1),
\quad u_{n+1}=u_n+\frac h2(K_1+K_2).
\]
Bảng hệ số là
\[
\begin{array}{c|cc}0&0&0\\1&1&0\\\hline&1/2&1/2\end{array}.
\]
Với $u'=u,u_0=1,h=0.1$:
\begin{center}\begin{tabular}{rlll}
\toprule Tầng&Thời điểm&Trạng thái thế vào F&Độ dốc\\\midrule
1&0&1&$K_1=1$\\
2&0.1&$1+0.1(1)=1.1$&$K_2=1.1$\\\bottomrule
\end{tabular}\end{center}
Tổng độ dốc $(1+1.1)/2=1.05$, nên $u_1=1.105$. Bậc toàn cục 2; với phương trình thử $R(z)=1+z+z^2/2$.

\section{RK3 mặc định: Kutta ba tầng}
\label{sec:ode-rk3}
\begin{align*}
K_1&=F(t_n,u_n),\\
K_2&=F(t_n+h/2,u_n+hK_1/2),\\
K_3&=F(t_n+h,u_n+h(-K_1+2K_2)),\\
u_{n+1}&=u_n+\frac h6(K_1+4K_2+K_3).
\end{align*}
Với cùng ví dụ, toàn bộ các tầng là:
\begin{center}\begin{tabular}{rlll}
\toprule Tầng&Thời điểm&Trạng thái trung gian&Độ dốc\\\midrule
1&0&1&1\\
2&0.05&$1+0.05(1)=1.05$&1.05\\
3&0.1&$1+0.1[-1+2(1.05)]=1.11$&1.11\\\bottomrule
\end{tabular}\end{center}
Vậy $u_1=1+0.1(1+4\cdot1.05+1.11)/6=1.105166666667$. Bậc toàn cục 3; $R(z)=1+z+z^2/2+z^3/6$ cho phương trình thử.

\section{RK4 cổ điển}
\label{sec:ode-rk4}
\begin{align*}
K_1&=F(t_n,u_n),&K_2&=F(t_n+h/2,u_n+hK_1/2),\\
K_3&=F(t_n+h/2,u_n+hK_2/2),&K_4&=F(t_n+h,u_n+hK_3),\\
u_{n+1}&=u_n+\frac h6(K_1+2K_2+2K_3+K_4).
\end{align*}
\begin{center}\begin{tabular}{rlll}
\toprule Tầng&Thời điểm&Trạng thái trung gian, $u'=u$&$K_i$\\\midrule
1&0&1&1\\
2&0.05&$1+0.05(1)=1.05$&1.05\\
3&0.05&$1+0.05(1.05)=1.0525$&1.0525\\
4&0.1&$1+0.1(1.0525)=1.10525$&1.10525\\\bottomrule
\end{tabular}\end{center}
Tổng có trọng số là $(1+2.1+2.105+1.10525)/6=1.051708333333\ldots$; nhân h và cộng 1 được
\[
u_1=1.105170833333\ldots=\frac{265241}{240000}.
\]
Sai số toàn cục điển hình $O(h^4)$. Đó không phải lời hứa bốn chữ số đúng với mọi h và mọi phương trình.

\section{RK2 tham số: bốn cách nhập cùng một họ}
\label{sec:ode-rk2-parameter}
Đặt $\alpha=c_2\ne0$. Họ hai tầng bậc hai:
\[
A=\begin{pmatrix}0&0\\\alpha&0\end{pmatrix},\quad
c=(0,\alpha)^T,\quad b_2=\frac1{2\alpha},\quad b_1=1-\frac1{2\alpha}.
\]
Nhập một trong bốn tham số mà giao diện nhận:
\begin{center}\begin{tabular}{lll}
\toprule Tên trường tham số&Giá trị cho&App suy ra\\\midrule
\texttt{alpha2}&$\alpha$&$c_2=\alpha$\\
\texttt{beta11}&$a_{21}$&$c_2=a_{21}$\\
\texttt{r2}&$b_2\ne0$&$c_2=1/(2b_2)$\\
\texttt{r1}&$b_1\ne1$&$c_2=1/[2(1-b_1)]$\\\bottomrule
\end{tabular}\end{center}
Sau khi lập bảng, app chạy cùng thuật toán RK hiện tổng quát. $\alpha=1$ là Heun; $\alpha=1/2$ là trung điểm; $\alpha=2/3$ là một lựa chọn có trọng số $(1/4,3/4)$.

Ví dụ tự dựng $\alpha=3/4$, $u'=u$, $h=0.1$: $K_1=1$, $K_2=1+0.1(3/4)=1.075$; $b=(1/3,2/3)$. Tổng $1/3+(2/3)1.075=1.05$, nên $u_1=1.105$. Kết quả một bước trùng RK2 mặc định trên phương trình tuyến tính này, nhưng không có nghĩa mọi họ cho cùng kết quả trên mọi F phi tuyến.

\section{RK3 tham số và hai nhánh suy biến được hỗ trợ}
\label{sec:ode-rk3-parameter}
Với $c_2\ne0$, $c_3\ne0$, $c_2\ne c_3$, đặt
\begin{align*}
b_2&=\frac{3c_3-2}{6c_2(c_3-c_2)},&
b_3&=\frac{2-3c_2}{6c_3(c_3-c_2)},&b_1&=1-b_2-b_3,\\
a_{21}&=c_2,&a_{32}&=\frac1{6b_3c_2},&a_{31}&=c_3-a_{32}.
\end{align*}
Cần $b_3\ne0$ trong nhánh này. Nhập c2,c3, lập bảng trước, rồi tính ba tầng theo mục~\ref{sec:ode-rk-general}.

Ví dụ $c_2=3/10,c_3=4/5$ cho
\[
b=(1/12,4/9,17/36),\quad a_{21}=3/10,
\quad a_{32}=20/17,\quad a_{31}=-32/85.
\]
Với $u'=u,h=0.1$: $K_1=1$, $K_2=1.03$,
\[
K_3=1+0.1\left[-\frac{32}{85}+\frac{20}{17}(1.03)\right]
=1.083529411765\ldots.
\]
Thế các b vào tổng cho $u_1=1.105166666667\ldots$, được đối chiếu với app.

Không dùng công thức có mẫu bằng 0 để kết luận mọi trường hợp trùng c đều vô hiệu. Code có hai bảng riêng khi $c_2=2/3$:
\[
\begin{array}{c|ccc}
0&0&0&0\\2/3&2/3&0&0\\0&-1&1&0\\\hline&0&3/4&1/4
\end{array}
\qquad(c_3=0),
\qquad
\begin{array}{c|ccc}
0&0&0&0\\2/3&2/3&0&0\\2/3&0&2/3&0\\\hline&1/4&3/8&3/8
\end{array}
\qquad(c_3=c_2).
\]
Trong ví dụ $u'=u,h=0.1$, bảng thứ nhất cho $(K_1,K_2,K_3)=(1,16/15,151/150)$; bảng thứ hai cho $(1,16/15,241/225)$. Hai tổng đều cho $u_1=1.105166666667\ldots$. Nếu $c_3=0$ hoặc $c_3=c_2$ nhưng $c_2$ khác $2/3$ trong kiểm tra số của code, app báo bộ tham số không thuộc nhánh được hỗ trợ.

\section{RK tùy chỉnh và hệ nhiều phương trình}
\label{sec:ode-rk-custom}
Ma trận A nhập dưới dạng JSON, b,c là các dãy. Code nhận $1\le s\le10$, A kích thước $s\times s$, b,c có s phần tử, mọi số hữu hạn. Điều kiện cấu trúc:
\[
a_{ij}=0\ (j\ge i),\qquad c_i=\sum_ja_{ij},\qquad\sum_i b_i=1.
\]
Đây là RK \emph{hiện}; không nhập bảng RK ẩn rồi mong tool giải các tầng đồng thời. Những kiểm tra trên xác nhận kích thước, tính hiện và điều kiện nhất quán cấp một, \emph{không tự xác nhận bảng tùy chỉnh đạt bậc 2,3 hay cao hơn}.

Ví dụ nhập A=$[[0,0],[1,0]]$, b=$[0.5,0.5]$, c=$[0,1]$ tái tạo RK2 Heun. Đề cho một bảng khác thì giữ đúng bảng và tự nêu bậc nếu đã chứng minh các điều kiện bậc.

Với hệ nhiều phương trình, mỗi $K_i$ là vector. Ví dụ tự dựng
\[
x'=y,\quad y'=-x,\quad (x(0),y(0))=(1,0),\quad h=0.1,
\]
nhập vế phải \texttt{y; -x}, điều kiện đầu \texttt{1 0}. RK4 có:
\begin{center}\begin{tabular}{rlll}
\toprule Tầng&t&Trạng thái trung gian $(x,y)$&Vector độ dốc\\\midrule
1&0&$(1,0)$&$(0,-1)$\\
2&0.05&$(1,-0.05)$&$(-0.05,-1)$\\
3&0.05&$(0.9975,-0.05)$&$(-0.05,-0.9975)$\\
4&0.1&$(0.995,-0.09975)$&$(-0.09975,-0.995)$\\\bottomrule
\end{tabular}\end{center}
Kết quả $(x_1,y_1)=(0.995004166667,-0.099833333333)$. Hai thành phần phải cập nhật từ cùng các vector tầng, không dùng $x$ mới để tính lại $y$ trong cùng bước.

\section{Adams 2--5: khởi tạo, dự báo rồi giải hiệu chỉnh}
\label{sec:ode-adams}
\subsection{Số bước và các hệ số đúng với code}
Gọi s là lựa chọn \texttt{Adams s}. Ở bước n, có s độ dốc cũ $F_n,F_{n-1},\ldots,F_{n-s+1}$.
\[
\mathbf u_{n+1}^P=\mathbf u_n+h\sum_{j=0}^{s-1}\beta_j\mathbf F_{n-j}
\]
là dự báo Adams--Bashforth. Hiệu chỉnh Adams--Moulton giải
\[
\mathbf v=\mathbf u_n+h\left[a_0\mathbf F(t_{n+1},\mathbf v)
+\sum_{j=0}^{s-1}a_{j+1}\mathbf F_{n-j}\right].
\]
\begin{center}\small\begin{tabular}{rll}
\toprule s&Vector $\beta$ của dự báo&Vector a của hiệu chỉnh\\\midrule
2&$(3,-1)/2$&$(5,8,-1)/12$\\
3&$(23,-16,5)/12$&$(9,19,-5,1)/24$\\
4&$(55,-59,37,-9)/24$&$(251,646,-264,106,-19)/720$\\
5&$(1901,-2774,2616,-1274,251)/720$&$(475,1427,-798,482,-173,27)/1440$\\\bottomrule
\end{tabular}\end{center}
Tên \texttt{Adams s} ở đây mô tả s giá trị cũ ở AB; AM dùng thêm giá trị mới nên vector a có s+1 phần tử. Không tự thay AM bằng một công thức s phần tử có cùng tên trong tài liệu khác.

\subsection{Từng giai đoạn thực sự chạy}
\begin{enumerate}
\item Bắt đầu chỉ có $\mathbf u_0$. Tính s−1 bước bằng RK4 để đủ lịch sử. Nếu toàn bộ đoạn quá ngắn, có thể chưa có bước Adams nào; mọi bước đang có đều là khởi tạo RK4.
\item Từ n=s−1, xếp lịch sử mới nhất trước. Tính từng $F_{n-j}$, nhân $\beta_j$, cộng và lập $\mathbf u_{n+1}^P$.
\item Tính phần đuôi đã biết $\mathbf T=\sum_{j=0}^{s-1}a_{j+1}\mathbf F_{n-j}$.
\item Lập $\mathbf R(\mathbf v)=\mathbf v-\mathbf u_n-h[a_0\mathbf F(t_{n+1},\mathbf v)+\mathbf T]$ và Jacobian $I-ha_0J_F$.
\item Gọi bộ giải root, dùng dự báo làm giá trị đầu. Kiểm phần dư và Jacobian như các bước ẩn ở trên.
\item Nhận giá trị hiệu chỉnh, thêm nó vào lịch sử, dịch cửa sổ để tính bước tiếp theo.
\end{enumerate}
App không chỉ thế dự báo vào AM đúng một lần. Nó giải phương trình hiệu chỉnh đến khi kiểm tra phần dư đạt. Cũng không có quy tắc ``s mốc cho phép s vòng''; các mốc thời gian cũ được dùng lặp lại để tính nhiều bước sau.

Các công thức AB s có bậc s và AM ở bảng có bậc s+1 theo lý thuyết đủ trơn. Sai số toàn cục của triển khai còn phụ thuộc chất lượng khởi tạo RK4 và việc giải hiệu chỉnh. Đặc biệt không tự ghi bậc toàn chương trình là 6 chỉ vì AM của lựa chọn s=5 có bậc 6.

\subsection{Ví dụ số cho đủ cả bốn lựa chọn}
Các bảng sau dùng cùng ví dụ tự dựng $u'=u$, $u_0=1,h=0.1$. Đối với hàm này $F_i=u_i$; phương trình hiệu chỉnh là tuyến tính nên có thể giải chính xác bằng một phép chia. Bảng ghi mọi bước khởi tạo rồi một bước Adams, được đối chiếu với đường tính thật của app.
\examplebreak
\subsubsection*{Ví dụ Adams 2: đủ 1 bước khởi tạo và một bước Adams}
Ví dụ tự dựng: $u^{\prime}=u$, $u(0)=1$, $h=0.1$. Gọi $R_4=265241/240000$ là hệ số tăng trưởng một bước RK4. Ở bước Adams dưới đây, $n=s-1$.
\begin{center}\begin{tabular}{rrll}\toprule $i$&$t_i$&$u_i$&Cách có giá trị\\\midrule
0&0&1&Điều kiện đầu\\
1&0.1&1.10517083333&RK4: $u_1 = R_4u_{0}$\\
\bottomrule\end{tabular}\end{center}
Ở bước Adams, xếp lịch sử từ mới nhất đến cũ nhất:
\begin{center}\begin{tabular}{rlll}\toprule $j$&$f_{n-j}$&$\beta_jf_{n-j}$&$a_{j+1}f_{n-j}$\\\midrule
0&1.10517083333&1.65775625&0.736780555556\\
1&1&-0.5&-0.0833333333333\\
\bottomrule\end{tabular}\end{center}
Dự báo: $u_{n+1}^P=u_n+0.1\sum\beta_jf_{n-j}=1.22094645833$.
Hiệu chỉnh giải phương trình tuyến tính $(1-0.1\,5/12)u_{n+1}=u_n+0.1\sum a_{j+1}f_{n-j}$.
Kết quả: $u_{n+1}=1.22140753623$ tại $t_{n+1}=0.2$. Đối chiếu app: $1.22140753623$.
Đây là phép giải đại số của phương trình hiệu chỉnh trong ví dụ tuyến tính, không phải nhật ký các lần gọi nội bộ của bộ giải phi tuyến.
\examplebreak
\subsubsection*{Ví dụ Adams 3: đủ 2 bước khởi tạo và một bước Adams}
Ví dụ tự dựng: $u^{\prime}=u$, $u(0)=1$, $h=0.1$. Gọi $R_4=265241/240000$ là hệ số tăng trưởng một bước RK4. Ở bước Adams dưới đây, $n=s-1$.
\begin{center}\begin{tabular}{rrll}\toprule $i$&$t_i$&$u_i$&Cách có giá trị\\\midrule
0&0&1&Điều kiện đầu\\
1&0.1&1.10517083333&RK4: $u_1 = R_4u_{0}$\\
2&0.2&1.22140257085&RK4: $u_2 = R_4u_{1}$\\
\bottomrule\end{tabular}\end{center}
Ở bước Adams, xếp lịch sử từ mới nhất đến cũ nhất:
\begin{center}\begin{tabular}{rlll}\toprule $j$&$f_{n-j}$&$\beta_jf_{n-j}$&$a_{j+1}f_{n-j}$\\\midrule
0&1.22140257085&2.34102159413&0.966943701923\\
1&1.10517083333&-1.47356111111&-0.230243923611\\
2&1&0.416666666667&0.0416666666667\\
\bottomrule\end{tabular}\end{center}
Dự báo: $u_{n+1}^P=u_n+0.1\sum\beta_jf_{n-j}=1.34981528582$.
Hiệu chỉnh giải phương trình tuyến tính $(1-0.1\,3/8)u_{n+1}=u_n+0.1\sum a_{j+1}f_{n-j}$.
Kết quả: $u_{n+1}=1.34985892504$ tại $t_{n+1}=0.3$. Đối chiếu app: $1.34985892504$.
Đây là phép giải đại số của phương trình hiệu chỉnh trong ví dụ tuyến tính, không phải nhật ký các lần gọi nội bộ của bộ giải phi tuyến.
\examplebreak
\subsubsection*{Ví dụ Adams 4: đủ 3 bước khởi tạo và một bước Adams}
Ví dụ tự dựng: $u^{\prime}=u$, $u(0)=1$, $h=0.1$. Gọi $R_4=265241/240000$ là hệ số tăng trưởng một bước RK4. Ở bước Adams dưới đây, $n=s-1$.
\begin{center}\begin{tabular}{rrll}\toprule $i$&$t_i$&$u_i$&Cách có giá trị\\\midrule
0&0&1&Điều kiện đầu\\
1&0.1&1.10517083333&RK4: $u_1 = R_4u_{0}$\\
2&0.2&1.22140257085&RK4: $u_2 = R_4u_{1}$\\
3&0.3&1.34985849706&RK4: $u_3 = R_4u_{2}$\\
\bottomrule\end{tabular}\end{center}
Ở bước Adams, xếp lịch sử từ mới nhất đến cũ nhất:
\begin{center}\begin{tabular}{rlll}\toprule $j$&$f_{n-j}$&$\beta_jf_{n-j}$&$a_{j+1}f_{n-j}$\\\midrule
0&1.34985849706&3.09342572243&1.21112304042\\
1&1.22140257085&-3.00261465334&-0.447847609312\\
2&1.10517083333&1.70380503472&0.162705706019\\
3&1&-0.375&-0.0263888888889\\
\bottomrule\end{tabular}\end{center}
Dự báo: $u_{n+1}^P=u_n+0.1\sum\beta_jf_{n-j}=1.49182010744$.
Hiệu chỉnh giải phương trình tuyến tính $(1-0.1\,251/720)u_{n+1}=u_n+0.1\sum a_{j+1}f_{n-j}$.
Kết quả: $u_{n+1}=1.49182437726$ tại $t_{n+1}=0.4$. Đối chiếu app: $1.49182437726$.
Đây là phép giải đại số của phương trình hiệu chỉnh trong ví dụ tuyến tính, không phải nhật ký các lần gọi nội bộ của bộ giải phi tuyến.
\examplebreak
\subsubsection*{Ví dụ Adams 5: đủ 4 bước khởi tạo và một bước Adams}
Ví dụ tự dựng: $u^{\prime}=u$, $u(0)=1$, $h=0.1$. Gọi $R_4=265241/240000$ là hệ số tăng trưởng một bước RK4. Ở bước Adams dưới đây, $n=s-1$.
\begin{center}\begin{tabular}{rrll}\toprule $i$&$t_i$&$u_i$&Cách có giá trị\\\midrule
0&0&1&Điều kiện đầu\\
1&0.1&1.10517083333&RK4: $u_1 = R_4u_{0}$\\
2&0.2&1.22140257085&RK4: $u_2 = R_4u_{1}$\\
3&0.3&1.34985849706&RK4: $u_3 = R_4u_{2}$\\
4&0.4&1.49182424008&RK4: $u_4 = R_4u_{3}$\\
\bottomrule\end{tabular}\end{center}
Ở bước Adams, xếp lịch sử từ mới nhất đến cũ nhất:
\begin{center}\begin{tabular}{rlll}\toprule $j$&$f_{n-j}$&$\beta_jf_{n-j}$&$a_{j+1}f_{n-j}$\\\midrule
0&1.49182424008&3.93883038944&1.47835638236\\
1&1.34985849706&-5.20070482063&-0.748046583789\\
2&1.22140257085&4.43776267409&0.408830582743\\
3&1.10517083333&-1.9555383912&-0.132773995949\\
4&1&0.348611111111&0.01875\\
\bottomrule\end{tabular}\end{center}
Dự báo: $u_{n+1}^P=u_n+0.1\sum\beta_jf_{n-j}=1.64872033636$.
Hiệu chỉnh giải phương trình tuyến tính $(1-0.1\,95/288)u_{n+1}=u_n+0.1\sum a_{j+1}f_{n-j}$.
Kết quả: $u_{n+1}=1.64872076496$ tại $t_{n+1}=0.5$. Đối chiếu app: $1.64872076496$.
Đây là phép giải đại số của phương trình hiệu chỉnh trong ví dụ tuyến tính, không phải nhật ký các lần gọi nội bộ của bộ giải phi tuyến.


\section{Đánh giá hội tụ và cách chép một bài ODE}
\begin{enumerate}
\item Ghi phương trình sau khi đổi tên biến và điều kiện đầu, không chỉ chép biểu thức trong ô nhập.
\item Ghi h, các $t_n$, số bước; kiểm tra lưới kết thúc đúng ở b.
\item Với Euler/RK hiện: công thức từng tầng, trạng thái trung gian, độ dốc, tổng trọng số và nghiệm mới. Với hệ, ghi vector hoặc tách cột đúng thứ tự.
\item Với phương pháp ẩn: phương trình dư theo ẩn mới, dự đoán, Jacobian, nghiệm bước và phần dư cuối. Nếu tool không cung cấp nhật ký bộ giải nội bộ thì không gán một dãy vòng tự nghĩ thành trace thực tế.
\item Với Adams: đánh dấu các bước RK4 khởi tạo; ghi lịch sử, dự báo AB, phương trình hiệu chỉnh AM và nghiệm đã hiệu chỉnh.
\item Nếu được phép kiểm thêm: chạy lại h/2 và so nghiệm tại các thời điểm chung. Với bậc p và đã ở miền tiệm cận, $\|U_{h/2}-U_h\|/(2^p-1)$ là gợi ý sai số lưới mịn, không tự là cận bảo đảm.
\end{enumerate}
Các ca từ chối gồm: sai số thành phần/vế phải; dùng biến trạng thái không tồn tại; h sai chiều hoặc không chia hết đoạn; hàm/Jacobian ngoài miền; bảng RK không hiện; tham số RK ở nhánh không hỗ trợ; bước ẩn không đạt phần dư; Jacobian bước suy biến; số tràn; quá giới hạn lưới. Tăng số chữ số in ra không sửa các nguyên nhân đó.

\paragraph{Nguồn đối chiếu trong repo.} \texttt{pps/differential.py}: các bảng \texttt{RK}, \texttt{AB}, \texttt{AM}, các hàm \texttt{custom\_rk2}, \texttt{custom\_rk3}, \texttt{tableau}, \texttt{rkstep}, \texttt{ode}. Số liệu được kiểm bằng \texttt{verification/calculus\_verification.py}. Các bảng tầng trong chương là phép tính theo công thức được đối chiếu với kết quả app; chúng không được nhận là log nội bộ mà app chưa ghi.

\chapter{Bài toán biên và trị riêng Sturm--Liouville}
\label{ch:bvp}
\section{Input, bảng tổng quát và output của hai tool bài toán biên}
\subsection*{Input: các ô đúng tên trên giao diện}
Bảng này được lấy từ biểu mẫu của tool. Giá trị mặc định là ví dụ nhập, cần thay bằng dữ liệu đề. Hai dãy $x,y$ là bảng dữ liệu riêng bên dưới nếu chức năng dùng bảng.
\subsubsection*{Ô chung của các lựa chọn trong nhóm}
\begin{center}\small\begin{tabular}{@{}p{.32\linewidth}p{.17\linewidth}p{.43\linewidth}@{}}\toprule
Ô nhập & Khóa JSON & Cách nhập / mặc định\\\midrule
Hàm p(x) & \texttt{p} & \texttt{1}\\[3pt]
Hàm q(x) & \texttt{q} & \texttt{0}\\[3pt]
Đầu đoạn a & \texttt{a} & \texttt{0}\\[3pt]
Cuối đoạn b & \texttt{b} & \texttt{1}\\[3pt]
Bước h & \texttt{h} & \texttt{0.05}\\[3pt]
\bottomrule\end{tabular}\end{center}
\subsubsection*{Ô riêng: Biên tổng quát}
\begin{center}\small\begin{tabular}{@{}p{.32\linewidth}p{.17\linewidth}p{.43\linewidth}@{}}\toprule
Ô nhập & Khóa JSON & Cách nhập / mặc định\\\midrule
Hàm f(x) & \texttt{f} & \texttt{2}\\[3pt]
Biên trái: alpha beta gamma & \texttt{left} & \texttt{0 1 0}
\par alpha*u′(a) + beta*u(a) = gamma\\[3pt]
Biên phải: alpha beta gamma & \texttt{right} & \texttt{0 1 0}
\par alpha*u′(b) + beta*u(b) = gamma\\[3pt]
\bottomrule\end{tabular}\end{center}
\subsubsection*{Ô riêng: Trị riêng}
\begin{center}\small\begin{tabular}{@{}p{.32\linewidth}p{.17\linewidth}p{.43\linewidth}@{}}\toprule
Ô nhập & Khóa JSON & Cách nhập / mặc định\\\midrule
Hàm r(x) & \texttt{f} & \texttt{1}\\[3pt]
Số trị riêng nhỏ nhất & \texttt{count} & \texttt{5}\\[3pt]
\bottomrule\end{tabular}\end{center}

\paragraph{Chọn đúng nhánh trước khi nhập.} \emph{Biên tổng quát} nhận $p,q,f$ và hai bộ $\alpha,\beta,\gamma$. \emph{Trị riêng} nhận $p,q,r$, số trị riêng; biên cố định $u(a)=u(b)=0$. Ô thứ ba cùng khóa JSON \texttt{f} nhưng nhãn trên GUI đổi thành $r(x)$ ở trị riêng.

\subsection{Biên tổng quát: dữ liệu hàm thành từng hàng hệ sai phân}
\textbf{Input}: phương trình $(pu')'-qu=-f$, đoạn $a<b$, $h>0$ chia hết đoạn, hai điều kiện $\alpha u'+\beta u=\gamma$. Đạo hàm ở cả hai đầu theo chiều $x$ tăng. Không dùng dấu pháp tuyến đầu trái khi chưa đổi về $u'$.
\textbf{Quá trình}: tạo $N+1$ nút $\to$ tính hệ số tại nút/nửa bước $\to$ lập bảng hệ $\to$ thay hai hàng biên $\to$ giải $\to$ thế lại kiểm phần dư.
\begin{center}\small\begin{tabular}{cllll}\toprule
Hàng&$A_i$ nhân $u_{i-1}$&$B_i$ nhân $u_i$&$C_i$ nhân $u_{i+1}$&Vế phải $F_i$\\\midrule
0&0&Hệ số biên trái&Hệ số biên trái&Biên trái\\
$1\le i<N$&$p(x_i-h/2)$&$-(A_i+C_i+h^2q(x_i))$&$p(x_i+h/2)$&$-h^2f(x_i)$\\
$N$&Hệ số biên phải&Hệ số biên phải&0&Biên phải\\\bottomrule
\end{tabular}\end{center}
Hàng trong nghĩa là $A_iu_{i-1}+B_iu_i+C_iu_{i+1}=F_i$, không có phép tổng ngầm nào khác. Công thức Robin nửa ô và Dirichlet được viết đầy đủ ở phần sau.
\begin{center}\small\begin{tabular}{cllll}\toprule
Nút&Hoành độ&Nghiệm&Vế trái đã thế lại&Phần dư\\\midrule
$i$&$x_i=a+ih$&$u_i$&$A_iu_{i-1}+B_iu_i+C_iu_{i+1}$&Vế trái $-F_i$\\\bottomrule
\end{tabular}\end{center}
\textbf{Output}: bảng hệ số/nghiệm theo nút, phần dư lớn nhất và đồ thị. Đây là nghiệm hệ sai phân; cần giảm $h$ để khảo sát sai số so với nghiệm liên tục. Không có ô dung sai lặp trong tool BVP.

\subsection{Trị riêng: dữ liệu hàm thành ma trận và các mode}
\textbf{Input}: $-(pu')'+qu=\lambda ru$, biên 0 tại hai đầu, $p$ dương trên các mặt lưới và $r$ dương tại nút trong, số trị riêng $1\le k\le N-1$.
\textbf{Quá trình}: bỏ hai giá trị biên 0 $\to$ lập $Ku=\lambda Ru$ $\to$ đổi sang ma trận đối xứng $R^{-1/2}KR^{-1/2}$ $\to$ lấy các trị riêng nhỏ nhất $\to$ trả hàm riêng ra các nút.
\begin{center}\small\begin{tabular}{clll}\toprule
Nút trong&Đường chéo K&Đường chéo phụ K&Trọng số R\\\midrule
$i$&$[p(x_i-h/2)+p(x_i+h/2)]/h^2+q(x_i)$&$-p(x_i+h/2)/h^2$&$r(x_i)$\\\bottomrule
\end{tabular}\end{center}
Đường chéo phụ của ma trận đối xứng sau đổi biến phải chia thêm $\sqrt{r_ir_{i+1}}$. Sau khi giải, thêm 0 ở hai đầu và chuẩn hóa mỗi mode có $\max_i|u_i|=1$.
\begin{center}\small\begin{tabular}{cllll}\toprule
Nút&$x_i$&Mode 1&Mode 2&$\cdots$, mode $k$\\\midrule
0&$a$&0&0&0\\
$i$&$a+ih$&$u_i^{(1)}$&$u_i^{(2)}$&$\cdots,u_i^{(k)}$\\
$N$&$b$&0&0&0\\\bottomrule
\end{tabular}\end{center}
\textbf{Output}: bảng $\lambda_1,\ldots,\lambda_k$ và bảng hàm riêng trên lưới. Dấu mode tùy ý; không coi hai vector khác nhau bởi dấu là hai nghiệm khác nhau. Sai số ma trận nhỏ không tự là cận sai số bài liên tục.


\section{Nhận dạng phương trình và ô nhập}
Nhóm \textbf{Bài toán biên} có hai nhánh khác nhau. Nhánh \textbf{Biên tổng quát} tìm một hàm $u$ từ vế phải và điều kiện tại hai đầu. Nhánh \textbf{Trị riêng} tìm các số $\lambda$ sao cho tồn tại nghiệm khác 0 với biên thuần nhất. Cùng nhập $p,q$, nhưng ô thứ ba là $f$ ở bài toán biên và là trọng số $r$ ở trị riêng.

\begin{center}\small
\begin{tabular}{@{}p{.20\linewidth}p{.35\linewidth}p{.35\linewidth}@{}}\toprule
Ô nhập & Biên tổng quát & Trị riêng\\\midrule
$p,q$ & Hệ số của $(pu')'-qu$ & Hệ số của $-(pu')'+qu$\\
$f$ hoặc $r$ & Vế phải $-f$ & Trọng số $r$ trong $\lambda r u$\\
$a,b,h$ & $a<b$, $h>0$, lưới chia đúng đoạn & Như cột trái, ít nhất một nút trong\\
Biên trái/phải & Ba số $\alpha,\beta,\gamma$ & Cố định $u(a)=u(b)=0$\\
Số trị riêng & Không dùng & Từ 1 đến số nút trong\\\bottomrule
\end{tabular}\end{center}

\textbf{Kiểm tra lưới}: $N=(b-a)/h$ phải là số nguyên trong dung sai máy; tool không tạo một bước cuối ngắn hơn. Lưới $x_i=a+ih$, $i=0,\ldots,N$. Hai nhánh giới hạn 10\,000 khoảng. Hàm và các giá trị trên lưới phải hữu hạn; $p$ còn được lấy tại nửa bước.

\section{Sai phân hữu hạn cho bài toán biên}
\subsection{Phương trình đang được giải}
\[
 (p(x)u'(x))'-q(x)u(x)=-f(x),\qquad a<x<b,
\]
\[
 \alpha_Lu'(a)+\beta_Lu(a)=\gamma_L,\qquad
 \alpha_Ru'(b)+\beta_Ru(b)=\gamma_R.
\]
Đạo hàm $u'$ tại \emph{cả hai đầu} đều theo chiều $x$ tăng. Nếu đề dùng đạo hàm pháp tuyến ngoài thì tại đầu trái phải đổi dấu trước khi nhập. Không nhập dấu phương trình theo trí nhớ: bài $-u''=2$ tương ứng $p=1,q=0,f=2$.

\subsection{Nút trong: tính thông lượng rồi lập hàng}
Tại $x_i$, $i=1,\ldots,N-1$, đặt $p_{i\pm1/2}=p(x_i\pm h/2)$. Xấp xỉ
\[
 (pu')'(x_i)\simeq
 \frac{p_{i+1/2}(u_{i+1}-u_i)-p_{i-1/2}(u_i-u_{i-1})}{h^2}.
\]
Nhân $h^2$ và chuyển vế:
\begin{align*}
 A_i&=p_{i-1/2},&C_i&=p_{i+1/2},\\
 B_i&=-(A_i+C_i+h^2q_i),&F_i&=-h^2f_i,\\
 A_iu_{i-1}+B_iu_i+C_iu_{i+1}&=F_i.
\end{align*}
Mỗi hàng dùng hai giá trị $p$ tại nửa bước và một giá trị $q,f$ tại nút. Không thế $p(x_i)$ thay cho cả hai giá trị nửa bước nếu $p$ biến thiên.

\subsection{Biên Dirichlet và biên Robin nửa ô}
Nếu $\alpha=0$ thì cần $\beta\ne0$, biên cho trực tiếp $u=\gamma/\beta$. Hàng ma trận tương ứng có đường chéo 1, vế phải $\gamma/\beta$, hệ số ngoài đường chéo bằng 0.

Nếu $\alpha\ne0$, đặt $d=+1$ tại trái, $d=-1$ tại phải; $x_b$ là đầu biên và $p_h=p(x_b+dh/2)$. Công thức nửa ô trong mã cho
\begin{align*}
 B_b&=-p_h-\frac{h^2q(x_b)}2
          +dh\frac\beta\alpha p(x_b),\\
 F_b&=-\frac{h^2f(x_b)}2
          +dh\frac\gamma\alpha p(x_b).
\end{align*}
Tại trái, $C_0=p_h$, $A_0=0$; tại phải, $A_N=p_h$, $C_N=0$. Đây là công thức thông lượng trên nửa ô; không phải hàng đạo hàm một phía $(u_1-u_0)/h$ đơn thuần. Hạng $h^2f/2$ và dấu $d$ cần được giữ.

\subsection{Trình tự giải và tiêu chuẩn hoàn tất}
\begin{enumerate}
\item Kiểm tra dữ liệu, tạo $N+1$ nút, lấy $p$ tại các mặt và $q,f$ tại nút.
\item Lập từng hàng nút trong theo $A_i,B_i,C_i,F_i$.
\item Thay hàng đầu và cuối bằng điều kiện biên đúng loại.
\item Giải hệ ba đường chéo bằng \texttt{scipy.linalg.solve\_banded}. Báo cáo không chứa nhật ký từng phép khử của bộ giải thư viện; không được gọi một bảng tự tính tay là nhật ký Thomas thực tế của nhánh này.
\item Kiểm tra nghiệm hữu hạn và tính phần dư $r_i=A_iu_{i-1}+B_iu_i+C_iu_{i+1}-F_i$; xuất $\max_i|r_i|$ cùng bảng hệ số/nghiệm.
\end{enumerate}
Đây là phép giải hữu hạn, không nhập dung sai vòng lặp. Phần dư đánh giá hệ rời rạc. Muốn đánh giá nghiệm liên tục, so với nghiệm biết trước hoặc tính lại khi giảm $h$ với giả thiết trơn phù hợp.

\section{Ví dụ Dirichlet: thấy đủ mọi hàng}
\textbf{Ví dụ tự dựng}: $u''=2$, $u(0)=0$, $u(1)=1$, chọn $h=0.5$. Nhập $p=1,q=0,f=-2$, biên trái \texttt{0 1 0}, biên phải \texttt{0 1 1}. Nghiệm liên tục để đối chiếu là $u=x^2$.

$N=2$, nút $0,0.5,1$. Hàng trong: $A=C=1$, $B=-2$, $F=-0.25(-2)=0.5$.
\[
 \begin{bmatrix}1&0&0\\1&-2&1\\0&0&1\end{bmatrix}
 \begin{bmatrix}u_0\\u_1\\u_2\end{bmatrix}
 =\begin{bmatrix}0\\0.5\\1\end{bmatrix}.
\]
Từ biên, $u_0=0,u_2=1$; hàng giữa cho $-2u_1+1=0.5$, nên $u_1=0.25$.
\begin{center}\begin{tabular}{rrrrrrr}\toprule
$x_i$&$A_i$&$B_i$&$C_i$&$F_i$&$u_i$&$r_i$\\\midrule
0&0&1&0&0&0&0\\
0.5&1&-2&1&0.5&0.25&0\\
1&0&1&0&1&1&0\\\bottomrule
\end{tabular}\end{center}
Ví dụ bậc hai này cho giá trị nút đúng trong số học chính xác; không suy ra bài toán bất kỳ đều có sai số bằng 0.

\section{Ví dụ Robin: kiểm tra dấu ở hai đầu}
Vẫn $u=x^2$, $p=1,q=0,f=-2$, nhưng biên
\[
 u'(0)+u(0)=0,\qquad u'(1)+u(1)=3.
\]
Nhập trái \texttt{1 1 0}, phải \texttt{1 1 3}; $h=0.5$. Tại trái $d=+1$:
\[
 B_0=-1+0.5=-0.5,\quad C_0=1,\quad F_0=0.25.
\]
Tại phải $d=-1$:
\[
 A_2=1,\quad B_2=-1-0.5=-1.5,\quad F_2=0.25-0.5(3)=-1.25.
\]
Toàn hệ:
\[
 \begin{bmatrix}-.5&1&0\\1&-2&1\\0&1&-1.5\end{bmatrix}
 \begin{bmatrix}u_0\\u_1\\u_2\end{bmatrix}
 =\begin{bmatrix}.25\\.5\\-1.25\end{bmatrix}.
\]
Từ hàng đầu $u_1=.25+.5u_0$; hàng giữa suy ra $u_2=1$; hàng cuối suy ra $u_1=.25$, rồi $u_0=0$. Thế lại ba hàng lần lượt cho $0.25,0.5,-1.25$, nên phần dư đều bằng 0 trong số học chính xác.

\subsection{Các trường hợp phải dừng hoặc xem lại}
Hai biên không được có đồng thời $\alpha=\beta=0$. Hai biên Neumann thuần ($\beta_L=\beta_R=0$) với $q=0$ bị tool từ chối vì nghiệm không duy nhất nếu tồn tại: cộng một hằng số vẫn thỏa phương trình và biên. Bài toán đó cần điều kiện tương thích và một chuẩn hóa bổ sung, chưa có ô nhập ở nhánh này.

$p>0$, $q\ge0$ và biên phù hợp là các giả thiết thường giúp bài toán elliptic có tính đặt đúng; mã BVP không bắt buộc mọi hàm $p,q$ theo đúng hai bất đẳng thức này. Hệ vẫn có thể suy biến với các hệ số/biên khác. Không hiểu việc dữ liệu được chấp nhận là một chứng minh tồn tại/duy nhất của bài toán liên tục.

\section{Trị riêng: từ phương trình vi phân đến ma trận đối xứng}
\subsection{Bài toán duy nhất được hỗ trợ ở GUI}
\[
 -(p(x)u')'+q(x)u=\lambda r(x)u,\qquad u(a)=u(b)=0.
\]
Các giá trị $p_{i+1/2}>0$ và $r_i>0$ trên lưới được kiểm tra; $q_i$ phải hữu hạn. Tool chưa có lựa chọn biên Robin/tuần hoàn ở nhánh trị riêng. Số trị riêng trả về là số nhỏ nhất theo thứ tự đại số, không phải số nhỏ nhất theo trị tuyệt đối.

Bỏ hai nút biên bằng 0, vector chưa biết là $u=(u_1,\ldots,u_{N-1})^T$. Lập
\[
 Ku=\lambda Ru,\qquad R=\operatorname{diag}(r_1,\ldots,r_{N-1}),
\]
\[
 K_{ii}=\frac{p_{i-1/2}+p_{i+1/2}}{h^2}+q_i,
 \qquad K_{i,i+1}=K_{i+1,i}=-\frac{p_{i+1/2}}{h^2}.
\]
Đặt $v=R^{1/2}u$, ta có bài toán đối xứng chuẩn
\[
 Bv=\lambda v,\quad B=R^{-1/2}KR^{-1/2},\quad
 B_{ii}=\frac{K_{ii}}{r_i},\quad
 B_{i,i+1}=\frac{K_{i,i+1}}{\sqrt{r_ir_{i+1}}}.
\]
Không được chỉ chia đường chéo cho $r_i$ rồi giữ nguyên đường chéo phụ khi $r$ không hằng.

\subsection{Các bước bộ tính thực hiện}
\begin{enumerate}
\item Tạo lưới; kiểm tra số trị riêng không vượt $N-1$.
\item Tính đường chéo và đường chéo phụ của $B$ bằng số thập phân độ chính xác 80 chữ số, từ các giá trị hàm đã được đánh giá bằng số máy.
\item Chọn $s=\max|B_{ij}|$ (nếu tất cả bằng 0 thì lấy 1), đưa $B/s$ về số máy. Chia thang trước khi gọi thư viện giúp tránh mất độ lớn quá nhỏ/lớn; không biến toàn bộ phép đánh giá hàm thành độ chính xác tùy ý.
\item Gọi \texttt{eigh\_tridiagonal}, chỉ lấy các chỉ số trị riêng được yêu cầu. Trả lại $\lambda=s\widetilde\lambda$.
\item Đổi $u=R^{-1/2}v$, thêm 0 tại hai đầu, chuẩn hóa từng hàm riêng sao cho $\max_i|u_i|=1$.
\end{enumerate}
Dấu của hàm riêng tùy ý: $u$ và $-u$ biểu diễn cùng một mode. App không lưu vòng lặp nội bộ của LAPACK. Các bảng trị riêng và hàm riêng là kết quả giải ma trận; sổ tay chỉ diễn giải bằng tay các ví dụ nhỏ.

\section{Ví dụ hai trị riêng và hai hàm riêng}
\textbf{Ví dụ tự dựng}: $p=1,q=0,r=1$, $[a,b]=[0,1]$, $h=1/3$, yêu cầu 2 trị riêng. Có hai nút trong, $h^{-2}=9$:
\[
 K=B=\begin{bmatrix}18&-9\\-9&18\end{bmatrix},\quad R=I.
\]
Giải tay
\[
 \det(B-\lambda I)=(18-\lambda)^2-81=0
 \Longrightarrow\lambda_1=9,\quad\lambda_2=27.
\]
Với $\lambda_1=9$, phương trình cho $u_1=u_2$; với $\lambda_2=27$, $u_1=-u_2$. Sau chuẩn hóa (chọn dấu để dễ đọc):
\begin{center}\begin{tabular}{rrr}\toprule
$x_i$&Mode 1, $\lambda=9$&Mode 2, $\lambda=27$\\\midrule
0&0&0\\
$1/3$&1&1\\
$2/3$&1&-1\\
1&0&0\\\bottomrule
\end{tabular}\end{center}
Thế lại: $B(1,1)^T=9(1,1)^T$ và $B(1,-1)^T=27(1,-1)^T$. Nếu đổi $r=4$ thì $B=K/4$, trị riêng thành $2.25,6.75$; các hàm riêng sau chuẩn hóa cực đại vẫn như bảng.

\subsection{Sai số rời rạc và giảm bước}
Với ví dụ $p=r=1,q=0$, nghiệm liên tục có $\lambda_k=(k\pi)^2$. Trị riêng \emph{rời rạc} đúng là
\[
 \lambda_{k,h}=\frac4{h^2}\sin^2\!\left(\frac{k\pi}{2N}\right),
 \qquad k=1,\ldots,N-1.
\]
Do đó $9$ và $27$ ở lưới thô không phải lỗi bộ giải khi khác $\pi^2$ và $4\pi^2$. Khai triển $\sin$ cho sai số bậc $h^2$ khi giữ $k$ cố định trong ví dụ trơn này.
\begin{center}\begin{tabular}{rrr}\toprule
$h$&$\lambda_{1,h}$&Sai lệch tuyệt đối so với $\pi^2$\\\midrule
$1/4$&9.37258300203&0.497021399414\\
$1/8$&9.74341983856&0.126184562534\\
$1/16$&9.83793643355&0.0316679675433\\\bottomrule
\end{tabular}\end{center}
Tỷ số sai lệch gần 4 khi giảm nửa bước phù hợp bậc hai. Đây là kiểm tra theo một nghiệm biết trước, không phải chứng nhận sai số cho mọi bài. Muốn so hàm riêng giữa hai lưới cần thống nhất dấu, chuẩn hóa và các nút so sánh.

\section{Bảng tự kiểm trước khi chép bài}
\begin{enumerate}
\item Viết lại đúng dấu $-f$ hoặc dấu $-(pu')'$ theo nhánh đã chọn.
\item Kiểm tra $N,h$, số nút trong, $p$ tại nửa bước và $r$ tại nút.
\item Với Robin, kiểm tra thứ tự $\alpha,\beta,\gamma$ và dấu đạo hàm đầu trái.
\item Với BVP, thế nghiệm vào từng hàng biên và hàng trong; ghi phần dư hệ.
\item Với trị riêng, kiểm tra $Ku-\lambda Ru$, biên 0, chuẩn hóa và dấu tương đương. Không nhận phần dư ma trận là sai số của phương trình liên tục.
\end{enumerate}

\appendix
\chapter{Các hàm Python bổ trợ và giải hệ tuyến tính}
\label{app:api}

\section{Đọc đúng tên hàm và thứ tự hệ số}
Các hàm dưới đây nằm trong \texttt{pps/chapter1\_tools.py}; có thể gọi bằng
\begin{verbatim}
from pps import chapter1_tools as ch1
\end{verbatim}
Chúng không phải tất cả đều có nút riêng trong GUI. Đặc biệt \textbf{hệ số đơn thức của mô-đun này theo bậc tăng}; ô Horner trên GUI lại theo bậc giảm. Ví dụ $2x^3-3x+1$ là \texttt{[1,-3,0,2]} khi gọi \texttt{horner\_division}, và \texttt{2 0 -3 1} khi nhập GUI.

\begin{center}\small\begin{tabular}{@{}p{.38\textwidth}p{.54\textwidth}@{}}\toprule
API & Vai trò và nơi đọc thuật toán\\\midrule
\texttt{chebyshev\_value(n,z)} & Tính $T_n(z)$ bằng truy hồi; xem ngay dưới bảng\\
\texttt{horner\_division(c,a)} & Chia đa thức đơn thức cho $x-a$; Chương \ref{ch:interp}\\
\texttt{newton\_horner(nodes,c,z)} & Horner dạng Newton theo đúng thứ tự mốc; Chương \ref{ch:interp}\\
\texttt{newton\_append(nodes,c,x,y)} & Thêm một mốc mà không tính lại hệ số cũ; Chương \ref{ch:interp}\\
\texttt{interpolation\_error\_bound} & Cận tại một điểm khi đề cho $M_m$; Chương \ref{ch:interp}\\
\texttt{interpolation\_global\_bound} & Cận toàn dải lưới đều; cần tên mẫu \texttt{'cách đều'}\\
\texttt{chebyshev\_global\_bound} & Cận trên $[a,b]$ cho đúng $m$ nghiệm Chebyshev; không lấy min/max mốc thay $a,b$\\
\texttt{fill\_missing} & Điền ô thiếu theo giả thiết bậc đa thức; quy trình dưới đây\\
\bottomrule\end{tabular}\end{center}

\begin{center}\small\begin{tabular}{@{}p{.38\textwidth}p{.54\textwidth}@{}}\toprule
API (tiếp) & Vai trò và nơi đọc thuật toán\\\midrule
\texttt{scan\_level}, \texttt{monotone\_branches} & Quét mức và nhánh đơn điệu của \emph{bảng}; Chương \ref{ch:inverse}\\
\texttt{inverse\_by\_swap} & Wrapper hàm ngược; dùng đúng tập mốc đã chọn\\
\texttt{bisect\_model} & Chia đôi một mô hình đã có, trả lịch sử; Chương \ref{ch:inverse}\\
\texttt{inverse\_controlled} & Wrapper lặp có chứng nhận, mặc định tối đa 1000 vòng; GUI mặc định 100\\
\texttt{spline\_degree} & API cũ bậc 1--4, biết $p-1$ đạo hàm trái; Chương \ref{ch:spline}\\
\texttt{spline\_boundary} & API cũ bậc 1--4, đúng $p-1$ điều kiện biên; GUI mới xét cả hệ thiếu/thừa\\
\texttt{spline\_at} & Giá trị/đạo hàm, kiểm tra hai phía tại nút; không tự ngoại suy\\
\texttt{thomas} & Giải hệ ba đường chéo có đổi thang từng hàng; ví dụ dưới đây\\
\texttt{gaussian\_elimination} & Hệ vuông, đổi thang hàng và chọn trụ theo cột; ví dụ dưới đây\\
\texttt{transformed\_fit} & Bốn mô hình biến đổi cố định; Chương \ref{ch:ls}\\
\texttt{original\_scale\_fit} & API cũ gọi tối ưu SciPy; khác nhánh GUI LM có đầy đủ vết bước thử\\\bottomrule
\end{tabular}\end{center}

Một \texttt{Report} thường chứa \texttt{value}, \texttt{lines}, \texttt{tables}. Tuy nhiên kiểu \texttt{value} phụ thuộc hàm: vector nghiệm, số, đối tượng spline hoặc hàm dự đoán. Không giả sử mọi kết quả đều có cùng cấu trúc.

\section{Truy hồi Chebyshev và cận sai số}
\[
 T_0(z)=1,\quad T_1(z)=z,\quad
 T_{k+1}(z)=2zT_k(z)-T_{k-1}(z).
\]
Tính theo $k=1,\ldots,n-1$, giữ hai giá trị cũ trước khi gán cặp mới. Ví dụ $z=0.5$: $T_0=1,T_1=.5,T_2=-.5,T_3=-1$. Hàm chấp nhận $z$ ngoài $[-1,1]$ nhưng có thể tràn vì độ lớn tăng nhanh. $n$ phải nguyên không âm.

Với $m$ mốc và cận đạo hàm cấp $m$ là $M_m$:
\[
 |f(z)-P(z)|\le\frac{M_m}{m!}\prod_{i=0}^{m-1}|z-x_i|.
\]
$M_m$ phải đúng trên đoạn chứa mọi mốc và điểm tính, kể cả khi ngoại suy. Với $m\ge2$ mốc đều bước $h$, cận toàn dải mà API dùng là
\[
 |f-P|\le\frac{M_m h^m}{4m}.
\]
Với $m$ nghiệm Chebyshev trên $[a,b]$, cận là
\[
 \frac{2M_m}{m!}\left(\frac{b-a}{4}\right)^m.
\]
Ví dụ $x=(0,1,2)$, $z=.5$, $M_3=6$: cận điểm bằng $.375$; cận lưới đều bằng $.5$. Với 3 nghiệm Chebyshev trên $[0,2]$ cùng $M_3$, cận bằng $.25$. Đây là cận khi có giả thiết đạo hàm, không phải sai số thực tế suy từ bảng đo.

\section{Điền ô thiếu khi đề cho trước bậc}
\texttt{fill\_missing(x,y,degree,tolerance=1e-9)} nhận các ô thiếu dạng \texttt{None} hoặc \texttt{nan}. Quy trình:
\begin{enumerate}
\item Kiểm tra hoành độ hữu hạn, phân biệt, đủ $d+1$ tung độ biết; $0\le d<m$.
\item Chọn $d+1$ ô biết đầu tiên theo thứ tự mảng để dựng đa thức Newton bậc không quá $d$.
\item Tính giá trị mô hình ở \emph{mọi} hoành độ, kể cả các ô biết còn lại.
\item Với các ô biết, kiểm $\max|P(x_i)-y_i|\le\tau\max|y_i|$. Nếu thang tung độ bằng 0 thì vế phải là 0; không tự thêm một ngưỡng tuyệt đối 1.
\item Chỉ điền các ô thiếu; giữ nguyên các ô đã biết nếu phép kiểm đạt.
\end{enumerate}
Ví dụ $x=(0,1,2,3)$, $y=(1,\mathrm{None},5,10)$, $d=2$. Dùng ba cặp biết $(0,1),(2,5),(3,10)$: tỷ sai phân cấp 1 là $2,5$, cấp 2 là $(5-2)/(3-0)=1$. Vì vậy
\[
 P(x)=1+2x+x(x-2)=1+x^2,\qquad y(1)=2.
\]
Kết quả $(1,2,5,10)$. Nếu có thêm một ô biết không nằm trên đa thức trong ngưỡng kiểm thì hàm báo mâu thuẫn, không sửa ô đó để ép giả thiết.

\section{Thomas: thứ tự khử tiến và thế ngược}
Nhập bốn vector cùng độ dài $n$: đường chéo dưới $a$, chính $b$, trên $c$ và vế phải $d$, với $a_0=c_{n-1}=0$. Hệ là
\[
 a_ix_{i-1}+b_ix_i+c_ix_{i+1}=d_i.
\]
API chia mỗi hàng cho $s_i=\max(|a_i|,|b_i|,|c_i|)$ trước khi khử; hàng toàn 0 bị từ chối. Các công thức dưới đây áp dụng cho hệ \emph{sau đổi thang}:
\begin{align*}
 m_i&=a_i/\widetilde b_{i-1},\\
 \widetilde b_i&=b_i-m_i c_{i-1},\qquad
 \widetilde d_i=d_i-m_i\widetilde d_{i-1},\quad i=1,\ldots,n-1,\\
 x_{n-1}&=\widetilde d_{n-1}/\widetilde b_{n-1},\\
 x_i&=(\widetilde d_i-c_ix_{i+1})/\widetilde b_i,\quad i=n-2,\ldots,0.
\end{align*}
Thomas không chọn đổi hàng. Pivot gần 0 có thể xuất hiện ở một hệ vẫn khả nghịch; khi đó cần bộ giải có chọn trụ. Dung sai pivot mặc định $10^{-12}$ xét trên hệ đã đổi thang. Không gọi lỗi pivot là một bằng chứng hệ gốc vô nghiệm.

\subsection{Ví dụ trọn vẹn ba ẩn}
\[
 \begin{bmatrix}2&-1&0\\-1&2&-1\\0&-1&2\end{bmatrix}x
 =\begin{bmatrix}1\\0\\1\end{bmatrix}.
\]
Chia cả ba hàng cho 2: $b=(1,1,1)$, $a=(0,-.5,-.5)$, $c=(-.5,-.5,0)$, $d=(.5,0,.5)$.
\begin{center}\begin{tabular}{rrrrr}\toprule
$i$&$m_i$&$\widetilde b_i$&$\widetilde d_i$&Nguồn vế phải\\\midrule
0&---&1&$1/2$&Khởi tạo\\
1&$-1/2$&$3/4$&$1/4$&$0-(-1/2)(1/2)$\\
2&$-2/3$&$2/3$&$2/3$&$1/2-(-2/3)(1/4)$\\\bottomrule
\end{tabular}\end{center}
Thế ngược lần lượt $x_2=1$, $x_1=(1/4+1/2)/(3/4)=1$, $x_0=(1/2+1/2)/1=1$. Kiểm hệ gốc cho phần dư $(0,0,0)$.
\begin{verbatim}
r = ch1.thomas([0,-1,-1], [2,2,2], [-1,-1,0], [1,0,1])
print(r.value)  # [1. 1. 1.]
\end{verbatim}

\section{API giải hệ: input và output}
\texttt{gaussian\_elimination(A,b,pivot\_tolerance=1e-12)} cần $A$ vuông hữu hạn, $b$ cùng cỡ. Không dùng API này để giải trực tiếp bài chữ nhật $\Phi c\approx z$; dùng nhánh hồi quy và các cách giải ở Chương \ref{ch:ls}.
Output là vector nghiệm, bảng chọn pivot thực tế và phần dư theo hệ gốc. Phần này không nhắc lại thuật toán khử.
Hạng thiếu trong dung sai bị từ chối; hàm không phân loại riêng vô nghiệm/vô số nghiệm như bộ báo cáo spline tổng quát. Nhật ký API này lưu bảng chọn pivot, không lưu mọi ma trận sau từng phép khử.

\section{Phân biệt API cũ và luồng lời giải hiện tại}
\texttt{spline\_degree} trả \texttt{PPoly}, truyền đạo hàm đầu trái cho bậc 1--4. \texttt{spline\_boundary} cũ nhận đúng $p-1$ điều kiện; GUI mới qua \texttt{solve\_spline} kiểm cả hạng và điều kiện dư. Không dùng giới hạn của API cũ để kết luận GUI chỉ làm được biên trái.

\texttt{original\_scale\_fit} cũ gọi \texttt{scipy.optimize.least\_squares}, với tối đa 10\,000 lần đánh giá và các ngưỡng nội bộ $10^{-12}$. Nó không cung cấp nhật ký LM của GUI. Nhánh GUI \textbf{SSE thang y} dùng \texttt{solve\_least\_squares}, ghi bước thử, tham số, hệ Jacobian, hệ số giảm chấn và bước nhận/từ chối. Hai cách không nhất thiết có cùng số vòng, cùng quỹ đạo hay cùng tiêu chuẩn dừng.

Các wrapper nội suy ngược cũ gọi bộ giải mới nhưng chuyển trạng thái thất bại thành ngoại lệ qua \texttt{require\_success}. Đọc báo cáo của \texttt{solve\_inverse} nếu cần giữ lịch sử thất bại để kiểm tra. Không lấy giá trị cuối của một phép lặp thất bại làm đáp số được công nhận.

\chapter{Áp dụng vào đề trong repo và khả năng chạy lại}
\label{app:exam}

\section{Đề 20241: kiểm dữ liệu trước khi chọn công thức}
Nguồn ảnh trong \texttt{kienthuc/}: \texttt{20241\_MI3042\_GK.jpg}. Workbook: \texttt{20241GKtest2.xlsx}. Bài giải và nhật ký: \texttt{examples/de\_20241}.

Bản kiểm đã chọn \textbf{Sheet3} theo ghi chú ảnh đề 2; có 106 cặp, bước $h=0.117$, tham số $\bar x=4,\bar y=6$. Các chỉ số mốc bên dưới bắt đầu từ 0. Các tệp JSON lưu đủ dữ liệu và lựa chọn để mở lại bằng \textbf{Mở bài}.

\subsection{Horner với đa thức cho sẵn}
\[
 P(x)=3.5x^5-2.78x^4+5.63x^3-15.28x+1.54.
\]
Nhập \texttt{3.5 -2.78 5.63 0 -15.28 1.54}, điểm 2.57, cấp 3. Hệ số $x^2$ bằng 0 phải giữ đúng cột. Các lượt Horner lặp và hệ số nhân giai thừa đã giải thích ở Chương \ref{ch:interp}; kết quả
\[
 P''(2.57)=1054.696646,\qquad P'''(2.57)=1249.3386.
\]
Tệp chạy lại: \texttt{A\_Horner.json}. Không dùng bảng dữ liệu đo để giải câu đã cho đa thức tường minh này.

\subsection{Stirling 9 mốc và đạo hàm cùng đa thức}
Nhập toàn bộ bảng, chọn Stirling, điểm 4, cấp 1, số mốc 9. Cửa sổ là các chỉ số 22--30, gốc giữa $x_0=4.042$; do đó
\[
 t=\frac{4-4.042}{.117}=-\frac{14}{39}.
\]
Lập sai phân trên đúng 9 mốc, cộng công thức Stirling và đạo hàm theo $t$ rồi chia $h$. Kết quả cùng đa thức:
\[
 P(4)=5.442986424022284,\qquad P'(4)=-2.673281525949639.
\]
Tệp \texttt{C\_Stirling\_9.json} giữ đúng đầu vào. Đạo hàm này là đạo hàm của mô hình 9 mốc, không phải dùng tiếp một công thức đạo hàm khác trên một tập mốc khác.

\subsection{Tìm mọi nghiệm ở mức 6}
Quét toàn bộ bảng trước khi lặp. Bốn khoảng đổi dấu là các cặp chỉ số 14--15, 23--24, 62--63, 82--83. Quét bảng chỉ chứng minh các khoảng có nghiệm nếu hàm thật liên tục; không chứng minh không có nghiệm tiếp xúc bị bỏ sót giữa hai dòng cùng dấu.

Mỗi nhánh dưới đây khai báo một tập 9 mốc, giữ đúng gốc đầu/cuối của tập đó, $\varepsilon_x=\varepsilon_y=10^{-6}$. Chọn chế độ có cận co và kiểm cả ba giả thiết ở Chương \ref{ch:inverse}.
\begin{center}\small
\begin{tabular}{@{}rllrr@{}}\toprule
Nhánh&Mốc gốc&Hướng&$\widehat x$&Số vòng\\\midrule
1&11--19&Tiến&2.72873297033552&13\\
2&22--30&Tiến&3.76076318833793&9\\
3&58--66&Tiến&8.35257882568784&6\\
4&78--86&Lùi&10.6415671769223&4\\\bottomrule
\end{tabular}\end{center}
\begin{center}\small
\begin{tabular}{rrrr}\toprule
Nhánh&Cận $q$&Cận vị trí vòng cuối $B$&Phần dư vòng cuối\\\midrule
1&$453609/1087240$&$2.97435\cdot10^{-7}$&$5.34363\cdot10^{-7}$\\
2&$24456/61985$&$7.12262\cdot10^{-7}$&$4.84535\cdot10^{-7}$\\
3&$51935/417732$&$4.85921\cdot10^{-8}$&$1.68418\cdot10^{-7}$\\
4&$20987/247968$&$1.87384\cdot10^{-7}$&$5.05965\cdot10^{-7}$\\\bottomrule
\end{tabular}\end{center}
Mọi nhánh đều đạt hai ngưỡng. Bảng trên là tổng kết; các tệp \texttt{D1\_Lap\_tien} đến \texttt{D4\_Lap\_lui} giữ \textbf{đủ mọi vòng và từng phép thế}, dưới dạng JSON đầu vào, TXT và HTML. Ví dụ nhỏ đủ ba/năm vòng trong Chương \ref{ch:inverse} giúp đọc các bảng dài này mà không cần in lặp hàng chục trang.

Nhánh 4 có thêm lượt thử tiến không đủ chứng nhận trong tệp \texttt{D4\_Thu\_tien\_chua\_du\_chung\_nhan}. Lượt đó không được tính là đáp số thành công; đổi sang lùi là một lựa chọn công khai trong lượt khác. Điều này minh họa tại sao ``có đủ mốc'' không có nghĩa mọi hướng lặp đều hội tụ.

\section{Ảnh 20251 và workbook test2: không trộn hai đề}
Ảnh \texttt{20251\_MI3042\_GK.jpg} chứa hai biến thể cạnh nhau: bên trái ghi 20251, mô hình $a\exp(bx+c/x)$, Stirling 7 và nhắc test1; bên phải in 20241, mô hình $a\exp(bx^2+c/x)$, Bessel 8. Repo cung cấp \texttt{20251GKtest2.xlsx}, có 104 cặp, $h=.115$, $\bar x=2.55$, $\bar y=1.23$. Vì vậy mọi kết quả dưới đây nói rõ mô hình chạy trên \textbf{test2}; không nhận là đáp số test1 chưa có dữ liệu.

\subsection{Bình phương tối thiểu: miền log quyết định cách giải}
Dữ liệu có cả tung độ âm và đổi dấu. Phép lấy $\ln y$ trên nguyên bảng không xác định; không được âm thầm bỏ dòng âm. Nếu đề bắt buộc tuyến tính hóa log, phải nêu giả thiết dữ liệu không thỏa. Nếu mục tiêu cho phép cực tiểu SSE trên thang $y$ gốc, chọn \textbf{Hàm mũ -- SSE thang y} và công khai sự lựa chọn đó.

Với $a\exp(bx^2+c/x)$, cơ sở nhập \texttt{x\^{}2;1/x}, không thêm hằng 1 vì đã có biên độ $a$. Bộ giải LM cho nghiệm dừng tìm được:
\begin{align*}
 a&=1.5172857551154062,\\
 b&=-0.2524591562528406,\\
 c&=1.9912315060358277,\\
 \mathrm{SSE}_y&=0.0012003355260271063.
\end{align*}
Với mô hình khác $a\exp(bx+c/x)$ trên \emph{cùng test2}, cơ sở \texttt{x;1/x}, kết quả là
\[
 (a,b,c)\simeq(11.6727348564,-1.29073115703,0.981281915998),
\]
\[
 \mathrm{SSE}_y\simeq0.0619256671698.
\]
Các kết quả không phải chứng nhận cực tiểu toàn cục. Nhật ký bước LM, cả bước thử bị từ chối, nằm ở \texttt{01a\_mu\_x2\_SSE\_104} và \texttt{01b\_mu\_x\_SSE\_104}; lỗi miền log ở \texttt{01c\_log\_khong\_ap\_dung}.

\subsection{Nội suy ngược: chọn mốc tường minh}
Mức $1.23$ có khoảng đổi dấu $[2.035,2.15]$. Hai lựa chọn \emph{khác tập mốc}, đều bậc 3 và có chứng nhận, $\varepsilon_x=\varepsilon_y=10^{-9}$:
\begin{center}\small\begin{tabular}{llll}\toprule
Lựa chọn&Chỉ số mốc&Kết quả&Số vòng\\\midrule
Tiến&9,10,11,12&2.130827273369&5\\
Lùi&7,8,9,10&2.130891279015&5\\\bottomrule
\end{tabular}\end{center}
Sự khác nhau ở chữ số sau không phải lỗi đổi biến: hai mô hình khác nhau. Nếu đề quy định tập mốc thì phải giữ đúng tập đó. Hai ca giữ \emph{nguyên 104 mốc} vẫn dựng đúng đa thức tối đa bậc 103 và thử phép lặp, nhưng tràn số tại vòng 3/2; không có đáp số thành công. Không dùng kết quả cục bộ thay lén cho ca toàn bảng thất bại.

\subsection{Stirling, Bessel và một bài spline tự đặt}
Tại $2.55$, cấp đạo hàm 1:
\begin{center}\small\begin{tabular}{lrr}\toprule
Mô hình&$P(2.55)$&$P'(2.55)$\\\midrule
Stirling 7&0.643602714824106&$-0.987821541452920$\\
Bessel 8&0.644033162232002&$-0.985994845413192$\\\bottomrule
\end{tabular}\end{center}
Tệp chạy tương ứng \texttt{03a\_Stirling\_7}, \texttt{03b\_Bessel\_8}. Hai tập mốc khác nhau; không lấy độ chênh này làm cận sai số đã chứng minh.

Ảnh này \textbf{không có câu spline}. Bài spline sau là tự đặt trên đúng 104 mốc để kiểm công cụ: cubic tự nhiên, không ngoại suy, tính tại $2.55$ và tích phân $[2.035,2.955]$. Kết quả:
\[
 S(2.55)=0.644019495981,\quad S'(2.55)=-0.981423912732,
\]
\[
 \int_{2.035}^{2.955}S(x)\,dx=0.696713440263.
\]
Tệp \texttt{04\_spline\_104} giữ hệ moment và quá trình Thomas. Không gán điều kiện tự nhiên này cho đề nếu đề cho các đạo hàm khác.

\section{Chạy lại và kiểm nguồn}
Mở JSON trong GUI rồi bấm \textbf{Tính kết quả}, hoặc từ thư mục gốc chạy:
\begin{verbatim}
.venv\Scripts\python.exe exam_cli.py examples\de_20241\A_Horner.json
\end{verbatim}
Các chương đối chiếu chủ yếu với các mô-đun:
\begin{itemize}
\item \texttt{interpolation.py}, \texttt{polynomial\_model.py}, \texttt{polynomial\_worked.py}: đa thức, đạo hàm, Horner.
\item \texttt{inverse\_solution.py}, \texttt{inverse\_model.py}, \texttt{inverse\_document.py}, \texttt{inverse\_export.py}: mô hình, chứng nhận, vòng lặp và báo cáo.
\item \texttt{spline\_worked.py}, \texttt{least\_squares\_worked.py}, \texttt{linear\_algebra.py}: spline, hồi quy, khử và phân rã.
\item \texttt{calculus.py}, \texttt{differential.py}, \texttt{chapter1\_tools.py}, \texttt{gui.py}: giải tích số, ODE, BVP, trị riêng, helper và đầu vào thực tế.
\end{itemize}
Tài liệu trong \texttt{kienthuc/real} cho quy ước trình bày và đối chiếu môn học, gồm PDF Chương 1 hiệu chỉnh, tài liệu nội suy ngược, spline và bình phương tối thiểu. Khi PDF lý thuyết mô tả một họ phương pháp rộng hơn, sổ tay này vẫn ghi giới hạn của mã hiện tại.

Các script trong \texttt{kienthuc/final/verification} tạo lại số ví dụ, bảng hệ số Newton--Cotes, bảng Adams và tệp kết quả JSON. \texttt{coverage.csv} ánh xạ mọi lựa chọn GUI vào chương. \texttt{Kiem\_tra\_va\_nguon.md} ghi lệnh chạy, kết quả kiểm tra và hạn chế của môi trường. Khi mã thay đổi, cần chạy lại các kiểm tra và biên dịch PDF, không chỉ sửa một con số trong bản in.

\end{document}
